% =====================================================================
%  DXC OASIS -- Intelligent IT Operations Assistant
%  Analyse complete du projet (architecture, backend, base de donnees,
%  workflows, MCP, actions sensibles, technologies, fichier par fichier)
%
%  Compilation :  pdflatex analyse-oasis.tex   (x2 pour la table des matieres)
%  Aucune dependance exotique : article + tikz + listings + booktabs.
% =====================================================================

\documentclass[11pt,a4paper]{report}

% ---------------------------------------------------------------- encodage
\usepackage[utf8]{inputenc}
\usepackage[T1]{fontenc}
\usepackage[french]{babel}
\usepackage{lmodern}
\usepackage{microtype}

% ---------------------------------------------------------------- maths
% amsmath : \text{} ; amssymb : \varnothing
\usepackage{amsmath}
\usepackage{amssymb}

% ---------------------------------------------------------------- mise en page
\usepackage[a4paper,margin=2.3cm,headheight=14pt]{geometry}
\usepackage{fancyhdr}
\usepackage{titlesec}
\usepackage{setspace}
\onehalfspacing

% ---------------------------------------------------------------- tableaux
\usepackage{booktabs}
\usepackage{longtable}
\usepackage{array}
\usepackage{tabularx}
\usepackage{multirow}

% ---------------------------------------------------------------- couleurs / boites
\usepackage{xcolor}
\usepackage{mdframed}

% ---------------------------------------------------------------- code
\usepackage{listings}

% ---------------------------------------------------------------- schemas
\usepackage{tikz}
\usetikzlibrary{arrows.meta,positioning,shapes.geometric,fit,backgrounds,calc}

% ---------------------------------------------------------------- palette
\definecolor{oasisblue}{HTML}{1D4ED8}
\definecolor{oasisdark}{HTML}{0F172A}
\definecolor{oasisgrey}{HTML}{64748B}
\definecolor{oasisline}{HTML}{E2E8F0}
\definecolor{oasisbg}{HTML}{F8FAFC}
\definecolor{oasisred}{HTML}{DC2626}
\definecolor{oasisgreen}{HTML}{059669}
\definecolor{oasisamber}{HTML}{D97706}
\definecolor{oasisviolet}{HTML}{7C3AED}

% ---------------------------------------------------------------- liens
% (chargé APRÈS la palette : hyperref a besoin des couleurs déjà définies)
\usepackage{hyperref}
\hypersetup{
  colorlinks=true,
  linkcolor=oasisblue,
  urlcolor=oasisblue,
  citecolor=oasisblue,
  pdftitle={DXC OASIS -- Analyse complete du projet},
  pdfauthor={Kawthar Khayet}
}

% ---------------------------------------------------------------- styles titres
\titleformat{\chapter}[display]
  {\normalfont\huge\bfseries\color{oasisdark}}
  {\color{oasisblue}\MakeUppercase{\chaptertitlename} \thechapter}{10pt}{\Huge}
\titleformat{\section}
  {\normalfont\Large\bfseries\color{oasisblue}}{\thesection}{1em}{}
\titleformat{\subsection}
  {\normalfont\large\bfseries\color{oasisdark}}{\thesubsection}{1em}{}
\titleformat{\subsubsection}
  {\normalfont\normalsize\bfseries\color{oasisgrey}}{\thesubsubsection}{1em}{}

\pagestyle{fancy}
\fancyhf{}
\fancyhead[L]{\small\color{oasisgrey}DXC OASIS --- Analyse complète du projet}
\fancyhead[R]{\small\color{oasisgrey}\thepage}
\renewcommand{\headrulewidth}{0.4pt}

% ---------------------------------------------------------------- listings
\lstdefinestyle{oasis}{
  basicstyle=\ttfamily\scriptsize,
  keywordstyle=\color{oasisblue}\bfseries,
  commentstyle=\color{oasisgreen}\itshape,
  stringstyle=\color{oasisamber},
  numberstyle=\tiny\color{oasisgrey},
  numbers=left,
  numbersep=7pt,
  breaklines=true,
  breakatwhitespace=false,
  showstringspaces=false,
  frame=single,
  rulecolor=\color{oasisline},
  backgroundcolor=\color{oasisbg},
  captionpos=b,
  tabsize=2,
  extendedchars=true,
  literate={é}{{\'e}}1 {è}{{\`e}}1 {à}{{\`a}}1 {ê}{{\^e}}1 {ç}{{\c c}}1
           {É}{{\'E}}1 {î}{{\^i}}1 {ô}{{\^o}}1 {û}{{\^u}}1 {ù}{{\`u}}1
           {â}{{\^a}}1 {ë}{{\"e}}1 {ï}{{\"i}}1 {œ}{{oe}}2 {«}{{\guillemotleft}}1
           {»}{{\guillemotright}}1 {–}{{--}}1 {—}{{---}}1 {…}{{\ldots}}1
}
\lstset{style=oasis}

% ---------------------------------------------------------------- boites
\newmdenv[
  linecolor=oasisblue,linewidth=2pt,topline=false,bottomline=false,rightline=false,
  backgroundcolor=oasisbg,skipabove=8pt,skipbelow=8pt,
  innerleftmargin=10pt,innerrightmargin=10pt,innertopmargin=8pt,innerbottommargin=8pt
]{infobox}

\newmdenv[
  linecolor=oasisred,linewidth=2pt,topline=false,bottomline=false,rightline=false,
  backgroundcolor=oasisbg,skipabove=8pt,skipbelow=8pt,
  innerleftmargin=10pt,innerrightmargin=10pt,innertopmargin=8pt,innerbottommargin=8pt
]{warnbox}

\newmdenv[
  linecolor=oasisgreen,linewidth=2pt,topline=false,bottomline=false,rightline=false,
  backgroundcolor=oasisbg,skipabove=8pt,skipbelow=8pt,
  innerleftmargin=10pt,innerrightmargin=10pt,innertopmargin=8pt,innerbottommargin=8pt
]{keybox}

% ---------------------------------------------------------------- raccourcis
\newcommand{\file}[1]{\texttt{\small #1}}
\newcommand{\fn}[1]{\texttt{\small\textcolor{oasisviolet}{#1}}}
\newcommand{\cls}[1]{\texttt{\small\textcolor{oasisblue}{#1}}}
\newcommand{\tbl}[1]{\texttt{\small\textcolor{oasisamber}{#1}}}
\newcommand{\route}[1]{\texttt{\small\textcolor{oasisgreen}{#1}}}
\newcommand{\arrowto}{\ensuremath{\rightarrow}}

% styles TikZ communs
\tikzset{
  box/.style={rectangle,rounded corners=3pt,draw=oasisline,line width=1pt,
              fill=white,minimum height=9mm,minimum width=26mm,align=center,
              font=\scriptsize\sffamily},
  bluebox/.style={box,draw=oasisblue,fill=oasisbg},
  redbox/.style={box,draw=oasisred,fill=white},
  greenbox/.style={box,draw=oasisgreen,fill=white},
  ambox/.style={box,draw=oasisamber,fill=white},
  vibox/.style={box,draw=oasisviolet,fill=white},
  flow/.style={-{Stealth[length=2.4mm]},draw=oasisgrey,line width=0.8pt,font=\tiny\sffamily},
  flowb/.style={flow,draw=oasisblue},
  flowr/.style={flow,draw=oasisred},
  lbl/.style={font=\tiny\sffamily\itshape,text=oasisgrey,align=center}
}

% =====================================================================
\begin{document}

% ---------------------------------------------------------------- page de titre
\begin{titlepage}
\centering
\vspace*{2cm}

{\color{oasisblue}\rule{\textwidth}{2pt}}\\[6mm]

{\Huge\bfseries\color{oasisdark} DXC OASIS}\\[3mm]
{\LARGE\color{oasisgrey} Intelligent IT Operations Assistant}\\[6mm]

{\color{oasisblue}\rule{\textwidth}{2pt}}\\[16mm]

{\Large\bfseries Analyse complète du projet}\\[4mm]
{\large\color{oasisgrey}
Architecture · Arborescence · Backend · Base de données\\
Workflows métier · MCP · Actions sensibles · Technologies\\
Analyse fichier par fichier}\\[20mm]

\begin{minipage}{0.8\textwidth}
\centering\small\color{oasisgrey}
Document d'analyse technique construit \emph{exclusivement} à partir du code
source réel du dépôt. Chaque affirmation renvoie à un fichier, une fonction ou
une ligne existante. Ce qui n'existe pas dans le code est signalé comme tel,
jamais complété par hypothèse.
\end{minipage}

\vfill

\begin{tabular}{rl}
\textbf{Dépôt} & \file{DXC-OASIS-Intelligent-IT-Operations-Assistant}\\
\textbf{Branche analysée} & \file{kaouthar}\\
\textbf{Auteur} & Kawthar Khayet\\
\textbf{Contexte} & PFA --- DXC Technology\\
\end{tabular}

\vspace{1cm}
\end{titlepage}

% ---------------------------------------------------------------- avertissement
\chapter*{Note de méthode}
\addcontentsline{toc}{chapter}{Note de méthode}

Ce document a été rédigé après lecture intégrale du code du dépôt. Trois règles
ont été tenues du début à la fin.

\begin{keybox}
\textbf{1. Rien n'est inventé.} Toute classe, fonction, table, route ou variable
citée existe dans le code. Les noms sont recopiés tels quels, y compris quand
ils sont maladroits ou incohérents.

\textbf{2. Ce qui est vide est dit vide.} Plusieurs fichiers du dépôt sont
présents mais ne contiennent aucune ligne (Dockerfiles, \file{nginx.conf},
\file{ci.yml}, \file{orchestrator.py}, \file{permissions.py}\ldots). Ils sont
listés au chapitre~\ref{ch:fichiers} plutôt que décrits comme s'ils
fonctionnaient.

\textbf{3. Ce qui est mort est dit mort.} Certains modules existent, compilent,
mais ne sont importés par personne (\file{base\_persona.py},
\file{llm\_router.py}, \file{rbac\_middleware.py}). Ils sont signalés comme
vestiges de la V1 et non présentés comme faisant partie du chemin d'exécution.
\end{keybox}

\noindent
Cette rigueur a un intérêt direct pour une soutenance : elle permet de répondre
sans hésiter à la question « et ça, ça marche vraiment ? ».

\tableofcontents
\clearpage

% =====================================================================
\chapter{Vue d'ensemble de l'application}
\label{ch:vue}

% ---------------------------------------------------------------------
\section{Quel est l'objectif de l'application ?}

OASIS est un \textbf{copilote d'exploitation informatique spécialisé Storage \&
Backup}. Il ne se contente pas de répondre à des questions : il lit l'état
\emph{réel} d'une machine de stockage, détecte les anomalies, ouvre des tickets
d'incident, conduit un diagnostic assisté par LLM et \textbf{propose} une
remédiation que l'ingénieur exécute d'un clic.

L'objectif se lit précisément dans quatre décisions de conception visibles dans
le code :

\begin{enumerate}
  \item \textbf{Aucune donnée simulée.} Le commentaire d'en-tête de
  \file{app/mcp/storage/ssh.py} l'énonce : « Aucun tool ne renvoie de données
  inventées. Quand la VM est injoignable ou qu'une commande échoue, on remonte
  une erreur explicite et l'interface l'affiche telle quelle plutôt que de
  combler le vide. » Les outils exécutent réellement \texttt{df}, \texttt{lvs},
  \texttt{systemctl}, \texttt{journalctl} et lisent \texttt{/proc/diskstats}
  sur une VM Ubuntu par SSH.

  \item \textbf{Le modèle propose, l'humain décide.} Les outils qui modifient la
  VM sont regroupés dans \fn{REMEDIATION\_TOOLS}
  (\file{app/mcp/storage/server.py}). Quand le LLM les appelle, l'agent
  n'exécute rien : il enregistre une ligne \tbl{action\_requests} au statut
  \texttt{pending}. Le seul chemin d'exécution est
  \route{POST /storage/actions/\{id\}/confirm}.

  \item \textbf{Tout est audité.} Chaque appel d'outil est écrit dans
  \tbl{mcp\_calls} (arguments réels, résultat brut, statut, corrélation), chaque
  requête HTTP dans \tbl{audit\_logs}, chaque étape du cycle de vie d'un
  incident dans \tbl{incident\_ticket\_notes}.

  \item \textbf{Le cycle se referme sur une vérification.} Un ticket n'est pas
  clos parce qu'une commande s'est terminée sans erreur, mais parce que
  l'indicateur qui l'avait ouvert est redevenu sain
  (\file{app/services/incident\_verification.py}).
\end{enumerate}

\begin{infobox}
\textbf{Formulation en une phrase pour la soutenance.} \emph{OASIS transforme
une conversation en langage naturel en actions vérifiées sur une infrastructure
de stockage réelle, sans jamais laisser le modèle de langage agir seul.}
\end{infobox}

% ---------------------------------------------------------------------
\section{Quelle est son architecture globale ?}

L'architecture réelle comporte \textbf{sept étages}, du navigateur jusqu'aux
commandes Linux :

\begin{longtable}{@{}p{2.6cm}p{4.3cm}p{7.6cm}@{}}
\toprule
\textbf{Étage} & \textbf{Technologie} & \textbf{Rôle exact dans le code}\\
\midrule
\endhead
Interface & React 18 + Vite + Tailwind & SPA mono-espace \file{/storage}, 8 pages, JWT en \texttt{localStorage} ou \texttt{sessionStorage}\\
API HTTP & FastAPI & 11 routeurs montés sous \route{/api/v1} dans \file{app/main.py}\\
Sécurité & JWT HS256 + RBAC & \file{core/security.py}, \file{api/deps.py}, \file{middleware/audit\_middleware.py}\\
Services & Modules Python & \file{app/services/}, \file{app/dashboard/kpi\_service.py}\\
Agent & Persona Storage + Groq & \file{app/personas/storage/agent.py} : RAG + tool calling + garde-fou\\
MCP & SDK \texttt{mcp} 2.0.0, transport stdio & Client \file{mcp/storage/client.py} \arrowto{} serveur \file{mcp/storage/tool\_server.py}\\
Infrastructure & Paramiko / SSH & \file{mcp/storage/ssh.py} \arrowto{} VM Ubuntu (LVM, systemd, restic, MinIO)\\
\bottomrule
\end{longtable}

\noindent
Trois magasins de données accompagnent ces étages :

\begin{itemize}
  \item \textbf{PostgreSQL 16} --- toutes les données métier (utilisateurs,
  incidents, actions, conversations, demandes, métriques, audit).
  \item \textbf{ChromaDB} --- la base vectorielle du RAG, interrogée en HTTP sur
  le port 8001 (\file{app/rag/vectorstore.py}).
  \item \textbf{Système de fichiers} --- les documents déposés (Data Hub) et les
  pièces jointes des demandes, sous \file{settings.upload\_dir}.
\end{itemize}

% ---------------------------------------------------------------------
\section{Schéma d'architecture adapté au code réel}

Le schéma ci-dessous ne reprend pas un diagramme générique : il suit exactement
le chemin d'exécution d'une question posée dans le chat.

\begin{figure}[h!]
\centering
\begin{tikzpicture}[node distance=7mm]

\node[box,fill=oasisbg] (user) {Ingénieur Storage\\\tiny(navigateur)};
\node[bluebox,below=of user] (front) {\textbf{Frontend React}\\\tiny ChatPage · IncidentsPage\\\tiny lib/api.ts (fetch + Bearer)};
\node[bluebox,below=of front] (api) {\textbf{API FastAPI}\\\tiny /api/v1/storage/chat\\\tiny deps.get\_current\_user + require\_role};
\node[vibox,below=of api] (agent) {\textbf{StoragePersona}\\\tiny personas/storage/agent.py\\\tiny handle\_message()};
\node[ambox,below=of agent] (client) {\textbf{MCP Client}\\\tiny StorageMCPClient\\\tiny transport stdio};
\node[ambox,below=of client] (server) {\textbf{MCP Server}\\\tiny tool\_server.py\\\tiny MCPServer("oasis-storage-mcp")};
\node[greenbox,below=of server] (tool) {\textbf{Tool Python}\\\tiny get\_capacity · list\_backups\\\tiny run\_backup · restore\_from\_backup};
\node[redbox,below=of tool] (ssh) {\textbf{SSH (paramiko)}\\\tiny mcp/storage/ssh.py};
\node[redbox,below=of ssh] (vm) {\textbf{VM Ubuntu Storage}\\\tiny df · lvs · systemctl · journalctl\\\tiny /proc/diskstats · restic · LVM};

% cote droit : dependances
\node[box,right=22mm of agent,draw=oasisviolet] (groq) {\textbf{Groq API}\\\tiny openai/gpt-oss-120b\\\tiny tool calling};
\node[box,right=22mm of api,draw=oasisblue] (pg) {\textbf{PostgreSQL}\\\tiny SQLAlchemy + Alembic};
\node[box,below=8mm of groq,draw=oasisviolet] (chroma) {\textbf{ChromaDB}\\\tiny RAG storage\_kb};

\draw[flowb] (user) -- node[right]{HTTPS} (front);
\draw[flowb] (front) -- node[right]{JSON + JWT} (api);
\draw[flowb] (api) -- node[right]{await} (agent);
\draw[flowb] (agent) -- node[right]{call\_tool()} (client);
\draw[flowb] (client) -- node[right]{stdio JSON-RPC} (server);
\draw[flowb] (server) -- node[right]{invoke} (tool);
\draw[flowr] (tool) -- node[right]{run()} (ssh);
\draw[flowr] (ssh) -- node[right]{commande} (vm);

\draw[flow,<->] (agent) -- node[above]{prompt / tool\_calls} (groq);
\draw[flow,<->] (agent.east) to[bend right=18] node[right,pos=0.6]{retrieve} (chroma.west);
\draw[flow,<->] (api) -- node[above]{ORM} (pg);

\end{tikzpicture}
\caption{Chemin complet d'une question du chat, tel qu'il est codé.}
\end{figure}

\clearpage

\begin{keybox}
\textbf{Formule condensée demandée, adaptée au vrai code :}
\medskip

\noindent
\texttt{Utilisateur} \arrowto{} \texttt{React (ChatPage)} \arrowto{}
\texttt{FastAPI /storage/chat} \arrowto{} \texttt{StoragePersona + Groq}
\arrowto{} \texttt{StorageMCPClient (stdio)} \arrowto{}
\texttt{MCPServer} \arrowto{} \texttt{Tool Python} \arrowto{}
\texttt{SSH paramiko} \arrowto{} \texttt{VM Ubuntu} \arrowto{}
\texttt{résultat JSON} \arrowto{} \texttt{Groq (reformulation)} \arrowto{}
\texttt{ChatTurnOut} \arrowto{} \texttt{React}
\end{keybox}

% ---------------------------------------------------------------------
\section{Une nuance essentielle : deux chemins vers les outils}

Le code comporte \textbf{deux voies distinctes} pour atteindre les mêmes
fonctions d'outil. Les confondre serait une erreur d'analyse.

\begin{longtable}{@{}p{4.2cm}p{5.1cm}p{5.2cm}@{}}
\toprule
& \textbf{Voie MCP (asynchrone)} & \textbf{Voie directe (in-process)}\\
\midrule
\endhead
Appelant & \file{personas/storage/agent.py}, \file{api/v1/actions.py}, \file{services/incident\_verification.py} & \file{dashboard/kpi\_service.py}, \file{services/incident\_watcher.py}, \file{services/metrics\_service.py}\\
Mécanisme & Sous-processus MCP, JSON-RPC sur stdio & \texttt{import} Python classique\\
Audit \tbl{mcp\_calls} & Oui (si \texttt{db} et \texttt{user\_id} fournis) & Non\\
Pourquoi ? & Le LLM doit voir un catalogue d'outils décrits & Pas de LLM impliqué : lecture déterministe et rapide\\
Exemple & \fn{mcp.call\_tool("get\_backup\_job\_status", ...)} & \fn{list\_backups()} importé depuis \file{tools/backup.py}\\
\bottomrule
\end{longtable}

\begin{infobox}
\textbf{Conséquence à connaître.} Le dashboard et la surveillance de fond
n'apparaissent \emph{pas} dans \tbl{mcp\_calls}. Seuls les appels déclenchés par
une conversation ou par une confirmation d'action y figurent --- ce qui est
cohérent avec l'objectif de cette table : prouver qu'une réponse de l'agent
s'appuie sur un relevé réel.
\end{infobox}

% ---------------------------------------------------------------------
\section{Quelles sont les différentes parties du projet ?}

\begin{longtable}{@{}p{2.6cm}p{3.2cm}p{8.7cm}@{}}
\toprule
\textbf{Partie} & \textbf{Emplacement} & \textbf{Contenu réel}\\
\midrule
\endhead
Backend & \file{backend/} & Application FastAPI, 14 migrations Alembic, couche MCP, RAG, personas, services\\
Frontend & \file{frontend/} & SPA React 18 / TypeScript / Vite / Tailwind, 8 pages, 4276 lignes\\
Laboratoire & \file{vm-lab/} & 4 scripts shell qui construisent l'infrastructure de démonstration dans une VM Ubuntu\\
Orchestration & \file{docker-compose.yml} & \textbf{Uniquement} PostgreSQL 16 et ChromaDB. Ni backend ni frontend n'y sont déclarés\\
Reverse proxy & \file{nginx/nginx.conf} & \textcolor{oasisred}{Fichier présent mais vide (0 octet)}\\
CI & \file{.github/workflows/ci.yml} & \textcolor{oasisred}{Fichier présent mais vide (0 octet)}\\
Images Docker & \file{backend/Dockerfile}, \file{frontend/Dockerfile} & \textcolor{oasisred}{Présents mais vides (0 octet)}\\
\bottomrule
\end{longtable}

\begin{warnbox}
\textbf{À dire tel quel en soutenance.} Le projet \emph{n'est pas} conteneurisé
de bout en bout aujourd'hui. Docker sert à faire tourner les deux dépendances
(PostgreSQL et ChromaDB) ; le backend et le frontend se lancent en local, comme
le montre \file{.claude/launch.json} : \texttt{uvicorn app.main:app -{}-reload
-{}-port 8000} et \texttt{npm run dev} sur le port 5173.
\end{warnbox}

% ---------------------------------------------------------------------
\section{Quel est le rôle du frontend ?}

Le frontend est une \textbf{application monopage} qui ne dessert qu'un seul
espace de travail. \file{src/App.tsx} le dit explicitement en commentaire :
« Cette version ne dessert qu'un seul espace de travail. Les autres rôles sont
renvoyés vers Storage plutôt que de se voir proposer une liste de personas. »

\subsection{Routage}

Tout est sous \route{/storage}, protégé par le composant \cls{RequireAuth} :

\begin{longtable}{@{}p{4.0cm}p{3.6cm}p{6.9cm}@{}}
\toprule
\textbf{Route} & \textbf{Composant} & \textbf{Ce que la page fait}\\
\midrule
\endhead
\route{/login} & \cls{LoginPage} & Formulaire e-mail / mot de passe, option « se souvenir »\\
\route{/storage/dashboard} & \cls{DashboardPage} & KPI capacité, sauvegardes, incidents, performance, insights, activité\\
\route{/storage/chat} & \cls{ChatPage} & Fils de conversation persistés, propositions d'action inline\\
\route{/storage/incidents} & \cls{IncidentsPage} & Liste filtrable, fiche détaillée, cycle de vie en 5 étapes, bouton « Investigate »\\
\route{/storage/requests} & \cls{RequestsPage} & Demandes ingénieur \(\leftrightarrow\) administrateur, pièces jointes, transitions de statut\\
\route{/storage/data-hub} & \cls{DataHubPage} & Dépôt de documents indexés dans le RAG\\
\route{/storage/reports} & \cls{ReportsPage} & Téléchargement des rapports PDF\\
\route{/storage/settings} & \cls{SettingsPage} & Écran de réglages\\
\bottomrule
\end{longtable}

\subsection{Responsabilités concrètes}

\begin{itemize}
  \item \textbf{Porter le jeton.} \file{lib/api.ts} lit le JWT dans
  \texttt{localStorage} ou \texttt{sessionStorage} et ajoute l'en-tête
  \texttt{Authorization: Bearer}.
  \item \textbf{Traiter la déconnexion globale.} Toute réponse 401 déclenche
  \fn{onUnauthorized()}, câblé sur \fn{signOut} par \cls{AuthProvider}.
  \item \textbf{Rendre les boutons de confirmation.} \cls{ActionProposal} est le
  \emph{seul} composant qui appelle \route{/storage/actions/\{id\}/confirm} ou
  \route{/reject}. Il est partagé entre le chat et la fiche d'incident.
  \item \textbf{Ne dupliquer aucune règle métier.} Les phrases affichées sur le
  dashboard (\fn{headline}, \fn{insights}) sont calculées côté backend dans
  \file{kpi\_service.py} ; le frontend les affiche.
  \item \textbf{Télécharger des fichiers authentifiés.} \fn{api.download} récupère
  la réponse en \emph{blob} car un simple \texttt{<a href>} n'enverrait pas
  l'en-tête \texttt{Authorization}.
\end{itemize}

% ---------------------------------------------------------------------
\section{Quel est le rôle du backend ?}

Le backend est le \textbf{centre de décision}. Il fait cinq choses que le
frontend ne fait jamais.

\begin{enumerate}
  \item \textbf{Authentifier et autoriser} --- \fn{get\_current\_user},
  \fn{require\_role}.
  \item \textbf{Orchestrer l'agent} --- construire le prompt (RAG + date + ticket),
  boucler sur les appels d'outils, intercepter les outils de remédiation.
  \item \textbf{Décider ce qui est un incident} --- les seuils sont dans
  \file{personas/storage/incident\_rules.py} (85\,\% / 95\,\%) et
  \file{tools/health.py} (90\,\% / 95\,\% / 24\,h), \textbf{jamais} dans le
  modèle.
  \item \textbf{Exécuter les actions confirmées} --- et elles seules.
  \item \textbf{Tenir trois tâches de fond} --- collecte de métriques,
  surveillance des sauvegardes, premier relevé de santé.
\end{enumerate}

\subsection{Le point d'entrée : \file{app/main.py}}

\begin{lstlisting}[language=Python,caption={Cycle de vie de l'application (extrait réel)}]
@asynccontextmanager
async def lifespan(_: FastAPI):
    tasks = [
        asyncio.create_task(metrics_collector.run_forever()),
        asyncio.create_task(incident_watcher.run_forever()),
        asyncio.create_task(_prime_storage_health()),
    ]
    try:
        yield
    finally:
        for task in tasks:
            task.cancel()
\end{lstlisting}

\noindent
Onze routeurs sont montés sous \route{/api/v1} : \texttt{cross\_domain},
\texttt{admin}, \texttt{incidents}, \texttt{auth}, \texttt{storage},
\texttt{storage/actions}, \texttt{storage/dashboard}, \texttt{storage/documents},
\texttt{storage/reports}, \texttt{requests}, \texttt{storage/search}.

\begin{warnbox}
\textbf{Point de sécurité à connaître.} Le middleware CORS est configuré avec
\texttt{allow\_origins=["*"]} et \texttt{allow\_credentials=True}. C'est
acceptable en développement ; en production, l'origine devrait être restreinte.
\end{warnbox}

% ---------------------------------------------------------------------
\section{Quelle base de données est utilisée et pourquoi ?}

\subsection{PostgreSQL 16}

Déclarée dans \file{docker-compose.yml} (\texttt{image: postgres:16}), pilotée
par SQLAlchemy en mode \textbf{synchrone} (\file{app/database/session.py} :
\fn{create\_engine} + \fn{sessionmaker}, pas d'\texttt{async\_sessionmaker}).

Quatre raisons \emph{vérifiables dans le code} justifient PostgreSQL plutôt
qu'un SQLite ou un MySQL :

\begin{enumerate}
  \item \textbf{JSONB.} Six colonnes l'utilisent :
  \tbl{incident\_tickets.mcp\_data}, \tbl{incident\_tickets.recommendations},
  \tbl{action\_requests.parameters}, \tbl{action\_requests.result},
  \tbl{mcp\_calls.arguments}, \tbl{mcp\_calls.result}. Le relevé brut d'un outil
  est stocké tel quel, et \emph{interrogé} : \file{incident\_service.py} exécute
  \texttt{IncidentTicket.mcp\_data[key].astext == str(identifier)} pour retrouver
  un ticket ouvert par identifiant technique. Aucun moteur sans JSON indexable ne
  permettrait cette requête.

  \item \textbf{UUID natif.} \texttt{sqlalchemy.dialects.postgresql.UUID} est
  utilisé comme clé primaire de sept tables.

  \item \textbf{\texttt{date\_bin}.} \file{metrics\_service.read\_series} agrège
  les métriques par intervalle de largeur fixe avec \fn{func.date\_bin}, une
  fonction \textbf{PostgreSQL 14+}. Le graphique de performance en dépend
  directement.

  \item \textbf{Types ENUM natifs.} \texttt{IncidentSeverity},
  \texttt{IncidentStatus}, \texttt{ActionRequestStatus}, \texttt{MCPCallStatus},
  \texttt{RequestStatus}, \texttt{RequestDirection}, \texttt{RequestPriority},
  \texttt{DocumentStatus}, \texttt{AccessRequestStatus} sont mappés en
  \texttt{Enum} SQLAlchemy.
\end{enumerate}

\subsection{ChromaDB}

Seconde base, purement vectorielle, atteinte en HTTP
(\fn{chromadb.HttpClient(host, port)}). Elle stocke les \emph{chunks} de la base
de connaissances Storage, vectorisés par \texttt{intfloat/multilingual-e5-small}
(384 dimensions).

\begin{keybox}
\textbf{Garde-fou remarquable.} \fn{get\_or\_create\_collection} grave le nom du
modèle d'embedding dans les métadonnées de la collection et \textbf{refuse} de
servir une collection indexée avec un autre modèle
(\cls{EmbeddingModelMismatch}). Raison donnée dans le code : l'ancien modèle
\texttt{bge-small-en-v1.5} produisait lui aussi des vecteurs de 384 dimensions,
donc Chroma aurait accepté la requête sans erreur et renvoyé « des réponses
fausses, mais parfaitement plausibles ».
\end{keybox}

% ---------------------------------------------------------------------
\section{Comment le frontend communique avec le backend ?}

Un unique module, \file{frontend/src/lib/api.ts}, concentre toute la
communication. Il n'y a ni Axios, ni React Query, ni client généré :
\texttt{fetch} natif enveloppé.

\begin{lstlisting}[caption={Contrat de communication (extraits de lib/api.ts)}]
const BASE_URL = import.meta.env.VITE_API_URL ?? '/api/v1';
const TOKEN_KEY = 'oasis.token';

// injection systematique du jeton
if (token) headers.set('Authorization', `Bearer ${token}`);

// FormData => pas de Content-Type impose (le navigateur pose la frontiere)
if (init.body && !(init.body instanceof FormData)) {
  headers.set('Content-Type', 'application/json');
}

// 401 => deconnexion globale, pas d'affichage d'erreur local
if (response.status === 401) { onUnauthorized(); throw new ApiError(...); }
\end{lstlisting}

\begin{longtable}{@{}p{2.4cm}p{12.5cm}@{}}
\toprule
\textbf{Verbe} & \textbf{Usage}\\
\midrule
\endhead
\fn{get} & Lectures : dashboard, incidents, conversations, demandes, documents\\
\fn{post} & Connexion, envoi d'un message, investigation, confirmation d'action\\
\fn{patch} & Renommage de conversation, statut d'incident, statut de demande\\
\fn{delete} & Suppression de conversation, de document, de demande, de pièce jointe\\
\fn{upload} & \texttt{FormData} avec un champ \texttt{file} : Data Hub et pièces jointes\\
\fn{download} & Récupération en \emph{blob} puis \texttt{URL.createObjectURL} : PDF et pièces jointes\\
\bottomrule
\end{longtable}

\noindent
En développement, \file{vite.config.ts} proxifie \route{/api} vers
\texttt{http://127.0.0.1:8000}, ce qui évite tout problème CORS côté navigateur.

% ---------------------------------------------------------------------
\section{Comment le backend communique avec les autres services ?}

\begin{longtable}{@{}p{2.9cm}p{2.9cm}p{8.6cm}@{}}
\toprule
\textbf{Service} & \textbf{Protocole} & \textbf{Où, et comment}\\
\midrule
\endhead
PostgreSQL & TCP / psycopg2 & \file{database/session.py} : \fn{create\_engine(settings.database\_url, pool\_pre\_ping=True)}\\
Groq (LLM) & HTTPS & \file{llm/groq\_client.py} : SDK \texttt{groq}, \fn{chat.completions.create} avec \texttt{tools} et \texttt{tool\_choice="auto"}\\
Serveur MCP & \textbf{stdio} & \file{mcp/storage/client.py} : \fn{stdio\_client(StdioServerParameters(command=sys.executable, args=["-m","app.mcp.storage.tool\_server"]))}\\
VM Storage & SSH (22) & \file{mcp/storage/ssh.py} : \texttt{paramiko.SSHClient}, authentification par clé, timeouts 5\,s / 15\,s\\
ChromaDB & HTTP (8001) & \file{rag/vectorstore.py} : \fn{chromadb.HttpClient}\\
Système de fichiers & local & \file{settings.upload\_dir} pour le Data Hub et \file{upload\_dir/requests} pour les pièces jointes\\
\bottomrule
\end{longtable}

\begin{infobox}
\textbf{Anthropic Claude est configuré mais pas branché.} \file{core/config.py}
déclare \texttt{anthropic\_api\_key}, et \file{orchestrator/llm\_router.py}
contient un \texttt{raise NotImplementedError("Claude pas encore branche")}. Le
seul LLM réellement appelé est \textbf{Groq}. Ne pas l'annoncer autrement.
\end{infobox}

% ---------------------------------------------------------------------
\section{Où intervient l'authentification ?}

L'authentification intervient à \textbf{quatre endroits distincts}.

\begin{enumerate}
  \item \textbf{À la connexion} --- \route{POST /api/v1/auth/login}.
  \fn{auth\_service.login} \arrowto{} \fn{authenticate\_user} (recherche par
  e-mail, \fn{verify\_password} bcrypt, contrôle \texttt{is\_active}) \arrowto{}
  \fn{create\_access\_token(subject=email, role=role.name)}. Le jeton porte
  \texttt{sub}, \texttt{role} et \texttt{exp}, signé HS256, durée
  \texttt{jwt\_expire\_minutes}.

  \item \textbf{À chaque requête protégée} --- la dépendance
  \fn{get\_current\_user} (\file{api/deps.py}) extrait le \texttt{Bearer} via
  \texttt{HTTPBearer}, décode, relit l'utilisateur en base par son e-mail et
  \textbf{revérifie} \texttt{is\_active}. Un compte désactivé après émission du
  jeton est donc rejeté immédiatement.

  \item \textbf{Au contrôle de rôle} --- \fn{require\_role("storage\_engineer",
  "administrator")} renvoie 403 sinon. Utilisé par les routeurs
  \texttt{storage}, \texttt{dashboard}, \texttt{knowledge}, \texttt{reports},
  \texttt{search}, \texttt{requests}, \texttt{admin} et par
  \route{/incidents/\{id\}/investigate}.

  \item \textbf{Dans le middleware d'audit} --- \fn{\_extract\_user\_id} redécode
  le jeton pour attribuer chaque ligne de \tbl{audit\_logs} à un utilisateur,
  y compris pour les requêtes refusées en 401 ou 403.
\end{enumerate}

\begin{warnbox}
\textbf{Deux points à énoncer honnêtement.}

\emph{(a)} Les routes de \file{api/v1/actions.py} exigent seulement
\fn{get\_current\_user}, \textbf{pas} \fn{require\_role}. Un utilisateur
authentifié quel que soit son rôle peut donc confirmer une action dont il
connaît l'identifiant. Le filtrage par rôle est présent partout ailleurs.

\emph{(b)} \fn{create\_cross\_domain\_token} produit un jeton à portée réduite
(\texttt{scope: cross\_domain\_read\_only}), mais \textbf{aucun code du dépôt ne
vérifie ce scope}. Le commentaire du fichier le dit d'ailleurs : « que le futur
MCP Client devra vérifier ». C'est une brique préparée, pas une brique active.
\end{warnbox}

% ---------------------------------------------------------------------
\section{Où intervient le MCP Server ?}

Oui, le projet utilise réellement MCP, avec le SDK officiel
\texttt{mcp==2.0.0}. Le chapitre~\ref{ch:mcp} y est entièrement consacré. En
résumé :

\begin{itemize}
  \item Le \textbf{serveur} vit dans \file{app/mcp/storage/tool\_server.py} :
  quatre lignes qui instancient \cls{MCPServer("oasis-storage-mcp")} et
  enregistrent les dix fonctions du \fn{TOOL\_REGISTRY}.
  \item Le \textbf{registre} vit dans \file{app/mcp/storage/server.py} : un
  dictionnaire nom \arrowto{} fonction, plus deux ensembles de classification du
  risque.
  \item Le \textbf{client} vit dans \file{app/mcp/storage/client.py} : il lance
  le serveur comme \textbf{sous-processus} et dialogue avec lui en JSON-RPC sur
  l'entrée/sortie standard.
  \item Il intervient \textbf{à trois moments} : pendant une conversation du
  chat, pendant une investigation d'incident, et à la confirmation d'une action.
\end{itemize}

% ---------------------------------------------------------------------
\section{Comment les outils/actions sont exécutés ?}

Deux régimes, décidés par un seul \texttt{if} dans
\file{personas/storage/agent.py} :

\begin{lstlisting}[language=Python,caption={Le coeur du garde-fou (agent.py)}]
if name in REMEDIATION_TOOLS:
    if audit_db is None:
        result = {"status": "denied", "reason": "confirmation_required", ...}
    else:
        result, action_id = self._propose_action(audit_db, tool_name=name,
                                                 arguments=arguments,
                                                 created_by=created_by,
                                                 incident=incident)
else:
    result = await mcp.call_tool(name, db=audit_db, user_id=created_by,
                                 correlation_id=correlation_id, **arguments)
\end{lstlisting}

\begin{longtable}{@{}p{3.6cm}p{5.4cm}p{5.5cm}@{}}
\toprule
& \textbf{Outils de lecture} & \textbf{Outils de remédiation}\\
\midrule
\endhead
Liste & \fn{get\_capacity}, \fn{list\_backups}, \fn{get\_backup\_job\_status}, \fn{get\_backup\_logs}, \fn{list\_snapshots}, \fn{get\_storage\_health} & \fn{run\_backup}, \fn{create\_snapshot}, \fn{restore\_from\_backup}, \fn{repair\_backup\_repository}\\
Depuis le chat & Exécutés immédiatement via MCP & \textbf{Jamais exécutés} : une ligne \tbl{action\_requests} \texttt{pending} est créée\\
Audit & \tbl{mcp\_calls} & \tbl{action\_requests}, puis \tbl{mcp\_calls} à la confirmation\\
Exécution réelle & N/A & Uniquement \route{POST /storage/actions/\{id\}/confirm}\\
\bottomrule
\end{longtable}

% ---------------------------------------------------------------------
\section{Quels sont les principaux flux de données ?}

\begin{longtable}{@{}p{3.4cm}p{11.5cm}@{}}
\toprule
\textbf{Flux} & \textbf{Trajet}\\
\midrule
\endhead
Conversation & \cls{ChatPage} \arrowto{} \route{POST /storage/chat} \arrowto{} \fn{storage\_persona.handle\_message} (RAG Chroma + Groq + MCP) \arrowto{} \fn{chat\_service.save\_turn} \arrowto{} \tbl{chat\_messages}\\
Détection & \fn{incident\_watcher.run\_forever} (300\,s) \arrowto{} \fn{list\_backups} (SSH) \arrowto{} \fn{check\_backup\_status} \arrowto{} \fn{create\_ticket\_from\_diagnosis} \arrowto{} \tbl{incident\_tickets}\\
Investigation & \route{POST /incidents/\{id\}/investigate} \arrowto{} \fn{incident\_investigation.start} \arrowto{} agent avec \cls{IncidentContext} \arrowto{} fil de chat + \tbl{action\_requests} + synthèse \texttt{root\_cause}\\
Confirmation & \cls{ActionProposal} \arrowto{} \route{POST /storage/actions/\{id\}/confirm} \arrowto{} \fn{confirm\_and\_execute\_action} (MCP, \texttt{confirm=True}) \arrowto{} \fn{\_close\_the\_loop} \arrowto{} \fn{incident\_verification.verify}\\
Dashboard & \route{GET /storage/dashboard/overview} \arrowto{} \fn{kpi\_service.get\_overview} \arrowto{} cache 60\,s de \fn{get\_storage\_health} + agrégats SQL\\
Métriques & \fn{metrics\_collector} (300\,s) \arrowto{} \fn{sample\_disk\_performance} (2 lectures de \texttt{/proc/diskstats}) \arrowto{} \tbl{disk\_metrics} \arrowto{} \fn{read\_series} (\texttt{date\_bin})\\
RAG & \cls{DataHubPage} \arrowto{} \route{POST /storage/documents} \arrowto{} \fn{knowledge\_service.create\_document} \arrowto{} extraction, \fn{chunk\_text}, \fn{embed\_passages} \arrowto{} Chroma \texttt{storage\_kb}\\
Demandes & \cls{RequestsPage} \arrowto{} \route{/requests} \arrowto{} \fn{request\_service} (direction déduite du rôle, transitions contrôlées) \arrowto{} \tbl{service\_requests} + \tbl{request\_events}\\
Recherche & \cls{GlobalSearch} \arrowto{} \route{GET /storage/search?q=} \arrowto{} 4 requêtes ILIKE (incidents, demandes, conversations, documents)\\
Audit HTTP & Toute requête \arrowto{} \cls{AuditMiddleware} \arrowto{} \tbl{audit\_logs} (méthode, chemin, statut, durée, IP)\\
\bottomrule
\end{longtable}

% =====================================================================
%  Chapitres suivants
% =====================================================================
% =====================================================================
\chapter{Arborescence du projet}
\label{ch:arbo}

Ce chapitre parcourt l'arborescence réelle. Le format demandé est respecté :
\textbf{dossier/fichier \arrowto{} rôle \arrowto{} dépendances \arrowto{}
utilisation}.

\section{Vue générale}

\begin{lstlisting}[caption={Arborescence réelle --- les fichiers vides sont signalés}]
DXC-OASIS-Intelligent-IT-Operations-Assistant/
|-- docker-compose.yml        postgres:16 + chromadb  (SEULEMENT ces deux-la)
|-- README.md                 vue "Architecture Freeze V1" (anterieure au code actuel)
|-- .env                      POSTGRES_*, DATABASE_URL, CHROMA_*, JWT_*, GROQ_*, ANTHROPIC_*
|-- .claude/launch.json       lancement backend (8000) / frontend (5173)
|-- nginx/nginx.conf          [VIDE - 0 octet]
|-- .github/workflows/ci.yml  [VIDE - 0 octet]
|
|-- backend/
|   |-- Dockerfile            [VIDE - 0 octet]
|   |-- requirements.txt      dependances (dont mcp==2.0.0, paramiko, groq, chromadb)
|   |-- alembic.ini           configuration des migrations
|   |-- alembic/
|   |   |-- env.py
|   |   `-- versions/         14 migrations, chaine lineaire
|   `-- app/
|       |-- main.py           FastAPI, 11 routeurs, lifespan (3 taches de fond)
|       |-- api/
|       |   |-- deps.py       get_db, get_current_user, require_role
|       |   `-- v1/           11 modules de routes
|       |-- core/
|       |   |-- config.py       Settings pydantic-settings (lit .env)
|       |   |-- security.py     bcrypt + JWT HS256 + token cross-domain
|       |   |-- logging.py
|       |   |-- permissions.py  [VIDE]
|       |   `-- constants.py    [VIDE]
|       |-- database/
|       |   |-- base.py       declarative_base()
|       |   |-- session.py    engine + SessionLocal + get_db
|       |   |-- seed.py       3 roles + 3 comptes de test
|       |   `-- init_db.py    [VIDE]
|       |-- models/           16 classes SQLAlchemy sur 11 fichiers (+ 2 vides)
|       |-- schemas/          schemas Pydantic (plusieurs vides)
|       |-- services/         16 services metier
|       |-- personas/
|       |   |-- base_persona.py       [MORT : importe par personne]
|       |   `-- storage/
|       |       |-- agent.py          StoragePersona (445 lignes) : coeur de l'IA
|       |       |-- prompts.py        STORAGE_SYSTEM_PROMPT
|       |       `-- incident_rules.py seuils de detection
|       |-- mcp/
|       |   `-- storage/
|       |       |-- server.py       TOOL_REGISTRY + DESTRUCTIVE/REMEDIATION_TOOLS
|       |       |-- tool_server.py  MCPServer("oasis-storage-mcp"), transport stdio
|       |       |-- client.py       StorageMCPClient + audit mcp_calls
|       |       |-- ssh.py          paramiko : run(), run_many(), quote()
|       |       `-- tools/          capacity, backup, snapshot, restore,
|       |                           repair, health, performance
|       |-- rag/
|       |   |-- embedding_service.py  E5 multilingue, prefixes query:/passage:
|       |   |-- vectorstore.py        Chroma + garde-fou EmbeddingModelMismatch
|       |   |-- chunker.py            500 mots, overlap 50
|       |   |-- retriever.py          point d'entree unique des personas
|       |   |-- query_pipeline.py     assemblage du contexte RAG
|       |   |-- indexing_pipeline.py  indexation d'un dossier
|       |   |-- reindex.py            reconstruction complete des collections
|       |   |-- storage_kb/           4 documents markdown (NetApp, Veeam, TrueNAS, AWS)
|       |   `-- storage_runbooks/     4 runbooks correspondants
|       |-- dashboard/kpi_service.py  agregation du dashboard (461 lignes)
|       |-- orchestrator/
|       |   |-- llm_router.py           [MORT : Claude non branche]
|       |   |-- orchestrator.py         [VIDE]
|       |   `-- cross_domain_manager.py [VIDE]
|       |-- middleware/
|       |   |-- audit_middleware.py   ACTIF : journalise toute requete
|       |   |-- rbac_middleware.py    [MORT : doublon de deps.require_role]
|       |   `-- auth_middleware.py    [VIDE]
|       |-- llm/groq_client.py        client Groq + tool calling + fallback modele
|       `-- reports/generator.py      [VIDE - la vraie generation est dans services/]
|
|-- frontend/
|   |-- Dockerfile            [VIDE - 0 octet]
|   |-- vite.config.ts        alias @ -> src, proxy /api -> 127.0.0.1:8000
|   |-- tailwind.config.ts
|   `-- src/
|       |-- App.tsx           routes /storage/*, RequireAuth
|       |-- auth/AuthContext.tsx   session, signIn, signOut
|       |-- lib/
|       |   |-- api.ts         fetch + Bearer + download blob + upload FormData
|       |   |-- types.ts       259 lignes de types miroirs des schemas Pydantic
|       |   |-- useResource.ts hook de chargement + reload
|       |   |-- format.ts      dates, humanize
|       |   `-- cn.ts          clsx + tailwind-merge
|       |-- components/
|       |   |-- ActionProposal.tsx  boutons Confirm / Reject
|       |   |-- layout/    AppLayout, Sidebar, Topbar, GlobalSearch, PageHeader
|       |   |-- dashboard/ KpiCard, CapacityOverview, PerformanceChart,
|       |   |              RecentActivity, OasisInsights, AttentionBanner
|       |   `-- ui/        Badge, Button, Card, Field, Overlay, States, Toast
|       `-- pages/         Login, Dashboard, Chat, Incidents, Requests,
|                          DataHub, Reports, Settings
|
`-- vm-lab/
    |-- 00-network.sh         IP fixe host-only + sshd
    |-- 10-storage-lab.sh     LVM, LUN, restic, timers systemd
    |-- 20-nas-object.sh      NFS, Samba, MinIO, Object Lock
    |-- inject-incident.sh    provoque une panne pour la demonstration
    `-- README.md             correspondance "equipement d'entreprise" -> "objet Linux"
\end{lstlisting}

\clearpage

\section{Backend : le noyau}

\begin{longtable}{@{}p{3.2cm}p{4.1cm}p{3.5cm}p{3.4cm}@{}}
\toprule
\textbf{Fichier} & \textbf{Rôle} & \textbf{Dépendances} & \textbf{Quand}\\
\midrule
\endhead
\file{app/main.py} & Assemble l'application : middlewares, 11 routeurs, 3 tâches de fond & Tous les routeurs, \file{kpi\_service}, \file{incident\_watcher}, \file{metrics\_collector} & Au démarrage d'uvicorn\\
\file{core/config.py} & Objet \cls{Settings} unique lisant \file{.env} & \texttt{pydantic-settings} & Importé partout via \fn{settings}\\
\file{core/security.py} & \fn{hash\_password}, \fn{verify\_password}, \fn{create\_access\_token}, \fn{decode\_access\_token}, \fn{create\_cross\_domain\_token} & \texttt{passlib[bcrypt]}, \texttt{python-jose} & Connexion, chaque requête, audit\\
\file{api/deps.py} & \fn{get\_current\_user} et fabrique \fn{require\_role} & \file{security.py}, \file{session.py}, \cls{User} & Dépendance de presque toutes les routes\\
\file{database/session.py} & \fn{engine}, \fn{SessionLocal}, générateur \fn{get\_db} & SQLAlchemy & Chaque requête, chaque tâche de fond\\
\file{database/base.py} & \fn{Base = declarative\_base()} & SQLAlchemy & Hérité par les 16 classes de modèle\\
\file{database/seed.py} & Crée 3 rôles et 3 comptes de test & \cls{Role}, \cls{User}, passlib & Script manuel\\
\file{middleware/audit\_middleware.py} & Journalise chaque requête HTTP & \file{security.py}, \cls{AuditLog} & Toutes les requêtes sauf \file{/docs}, \file{/openapi.json}, \file{/redoc}, \file{/favicon.ico}\\
\bottomrule
\end{longtable}

\section{Backend : routes (\file{app/api/v1/})}

\begin{longtable}{@{}p{2.7cm}p{3.1cm}p{4.5cm}p{4.0cm}@{}}
\toprule
\textbf{Fichier} & \textbf{Préfixe} & \textbf{Rôle} & \textbf{Appelle}\\
\midrule
\endhead
\file{auth.py} & \route{/auth} & Connexion, profil courant & \fn{auth\_service.login}, \fn{get\_current\_user}\\
\file{storage.py} & \route{/storage} & Chat : CRUD des conversations + tour de dialogue & \fn{chat\_service}, \fn{storage\_persona}, \fn{incident\_investigation}\\
\file{actions.py} & \route{/storage/actions} & Propose / confirme / rejette une action & \fn{incident\_service}, \cls{StorageMCPClient}, \fn{incident\_verification}\\
\file{incidents.py} & \route{/incidents} & Liste, fiche détaillée, investigation, statut & \fn{incident\_service}, \fn{incident\_investigation}\\
\file{dashboard.py} & \route{/storage/dashboard} & Overview, performance, activité, insights & \fn{kpi\_service}\\
\file{knowledge.py} & \route{/storage/documents} & Data Hub : dépôt, lecture, suppression & \fn{knowledge\_service}\\
\file{reports.py} & \route{/storage/reports} & PDF global et PDF d'incident & \fn{report\_service}\\
\file{requests.py} & \route{/requests} & Demandes ingénieur / admin + pièces jointes & \fn{request\_service}\\
\file{search.py} & \route{/storage/search} & Recherche globale multi-entités & \fn{search\_service}\\
\file{admin\_routes.py} & \route{/admin} & Comptes, rôles, journaux d'audit (lecture) & \fn{admin\_service}, \fn{audit\_service}\\
\file{cross\_domain.py} & \route{/cross-domain} & Demandes d'accès temporaire & \fn{access\_request\_service}\\
\bottomrule
\end{longtable}

\section{Backend : services (\file{app/services/})}

\begin{longtable}{@{}p{3.5cm}p{1.0cm}p{10.5cm}@{}}
\toprule
\textbf{Fichier} & \textbf{L.} & \textbf{Rôle et moment d'utilisation}\\
\midrule
\endhead
\file{incident\_service.py} & 367 & Cœur transactionnel : \fn{create\_action\_request}, \fn{confirm\_and\_execute\_action}, \fn{reject\_action\_request}, \fn{create\_ticket\_from\_diagnosis}, \fn{resolve\_ticket}, \fn{add\_note}, \fn{has\_completed\_action}. Appelé par l'agent, les routes actions et incidents, et le watcher.\\
\file{incident\_investigation.py} & 254 & Ouvre le fil d'investigation, fait tourner l'agent avec le contexte du ticket, synthétise \texttt{root\_cause} et \texttt{impact\_summary} par un second appel Groq, propose une action de repli.\\
\file{incident\_verification.py} & 151 & Relit sur la VM l'indicateur qui avait ouvert le ticket. Renvoie un \cls{Verdict} (\texttt{resolved}, \texttt{detail}, \texttt{conclusive}).\\
\file{incident\_watcher.py} & 264 & Tâche de fond (300\,s) : détecte les jobs en échec, ouvre les tickets, referme ceux dont la remédiation confirmée a produit son effet.\\
\file{chat\_service.py} & 184 & Persistance des fils : \fn{save\_turn}, \fn{history\_for\_llm} (10 derniers messages), \fn{attach\_actions}, \fn{title\_from\_message}.\\
\file{request\_service.py} & 494 & Demandes : direction déduite du rôle, table \fn{ALLOWED\_TRANSITIONS}, pièces jointes stockées sous UUID.\\
\file{report\_service.py} & 489 & Génération PDF avec ReportLab : \fn{build\_global\_report}, \fn{build\_incident\_report}.\\
\file{knowledge\_service.py} & 235 & Data Hub : validation, extraction de texte, découpage, embeddings, \texttt{upsert} Chroma, statut d'indexation.\\
\file{search\_service.py} & 233 & Recherche ILIKE échappée sur 4 entités, avec extrait contextuel.\\
\file{metrics\_service.py} & 144 & \fn{collect\_once}, \fn{purge\_old}, \fn{read\_series} (agrégation \texttt{date\_bin} en deux passes).\\
\file{metrics\_collector.py} & 64 & Boucle de fond (300\,s) et purge quotidienne (rétention 31 jours).\\
\file{mcp\_audit.py} & 70 & Point d'écriture unique de \tbl{mcp\_calls}.\\
\file{access\_request\_service.py} & 80 & Workflow cross-domain : création, approbation avec jeton temporaire, refus.\\
\file{auth\_service.py} & 21 & \fn{authenticate\_user} et \fn{login}.\\
\file{admin\_service.py} & 15 & Lecture des comptes et des rôles.\\
\file{audit\_service.py} & 16 & Lecture des journaux d'audit.\\
\bottomrule
\end{longtable}

\section{Frontend}

\begin{longtable}{@{}p{3.4cm}p{4.2cm}p{3.4cm}p{3.2cm}@{}}
\toprule
\textbf{Fichier} & \textbf{Rôle} & \textbf{Dépendances} & \textbf{Quand}\\
\midrule
\endhead
\file{main.tsx} & Monte \cls{App} dans \texttt{\#root} & React DOM & Chargement\\
\file{App.tsx} & Routeur + \cls{AuthProvider} + \cls{ToastProvider} + \cls{RequireAuth} & react-router-dom & Toujours\\
\file{auth/AuthContext.tsx} & Session : \fn{signIn}, \fn{signOut}, reprise via \route{/auth/me} & \file{lib/api.ts} & Au montage et à la connexion\\
\file{lib/api.ts} & Client HTTP unique & \texttt{fetch} & Chaque appel réseau\\
\file{lib/useResource.ts} & Hook \texttt{data / error / loading / reload} & \cls{ApiError} & Dashboard, incidents, demandes, Data Hub\\
\file{components/ActionProposal.tsx} & Confirmation humaine & \file{api.ts}, \cls{Badge}, \cls{Button}, \cls{Toast} & Chat et fiche d'incident\\
\file{components/layout/GlobalSearch.tsx} & Barre de recherche et navigation & \route{/storage/search} & Topbar\\
\file{pages/ChatPage.tsx} & 738 lignes : liste de fils, envoi, renommage, suppression & \file{api.ts}, \cls{ActionProposal} & Route \route{/storage/chat}\\
\file{pages/IncidentsPage.tsx} & 521 lignes : filtres, fiche, cycle de vie, investigation & \file{api.ts}, \cls{ActionProposal} & Route \route{/storage/incidents}\\
\file{pages/RequestsPage.tsx} & 1318 lignes : boîtes « sent » et « received », pièces jointes & \file{api.ts} & Route \route{/storage/requests}\\
\bottomrule
\end{longtable}

\section{Laboratoire \file{vm-lab/}}

Ce dossier n'est pas de la documentation : ce sont les scripts qui
\emph{fabriquent} l'infrastructure que les outils interrogent. Le principe est
énoncé dans son README : « l'application ne simule rien \ldots{} c'est donc la
machine qu'il faut équiper --- pas le code ».

\begin{longtable}{@{}p{4.6cm}p{10.2cm}@{}}
\toprule
\textbf{Équipement d'entreprise} & \textbf{Équivalent construit dans la VM}\\
\midrule
\endhead
Disques d'une baie SAN & Fichiers loopback creux (\texttt{truncate} + \texttt{losetup}) déclarés en PV\\
Pool de la baie & Volume group \texttt{vg\_prod}\\
LUN présenté à un hyperviseur & Logical volume formaté, monté sous \file{/srv/luns/}\\
Volume NFS d'un NAS & \texttt{nfs-kernel-server}, export \file{/srv/nas/nfs\_prod}\\
Partage SMB & Samba, partage \texttt{finance} authentifié\\
Stockage objet S3 & MinIO\\
Logiciel de sauvegarde (Veeam) & \texttt{restic} piloté par des timers systemd\\
Cible cloud hors-site & Dépôt restic sur MinIO\\
Sauvegarde immuable & Object Lock MinIO en mode COMPLIANCE\\
Snapshot de baie et restauration & \texttt{lvcreate -{}-snapshot}, \texttt{lvconvert -{}-merge}\\
\bottomrule
\end{longtable}

\begin{infobox}
Le README du laboratoire liste aussi honnêtement \textbf{ce qui n'est pas
simulable} : Fibre Channel et zonage (aucune émulation FC n'existe), robots de
bandes LTO (\texttt{mhvtl} exigerait la compilation d'un module noyau) et
microcodes de baie. C'est exactement le genre de limite qu'un jury apprécie de
voir assumée.
\end{infobox}

% =====================================================================
\chapter{Le backend en profondeur}
\label{ch:backend}

Ce chapitre traite les fichiers backend importants selon la grille demandée :
objectif, imports, classes, fonctions, paramètres, valeurs de retour, appels
internes, accès base de données, endpoints, gestion des erreurs,
authentification et dépendances. Les parties complexes sont décomposées étape
par étape.

% ---------------------------------------------------------------------
\section{Configuration --- \file{app/core/config.py}}

\subsection*{Objectif}
Fournir un objet \fn{settings} unique, validé, construit à partir des variables
d'environnement et du fichier \file{.env}.

\subsection*{Imports importants}
\texttt{pydantic\_settings.BaseSettings} et \texttt{SettingsConfigDict}.

\subsection*{Classe \cls{Settings}}

\begin{longtable}{@{}p{4.4cm}p{2.4cm}p{7.9cm}@{}}
\toprule
\textbf{Champ} & \textbf{Défaut} & \textbf{Rôle}\\
\midrule
\endhead
\texttt{database\_url} & \emph{requis} & Chaîne DSN PostgreSQL\\
\texttt{jwt\_secret\_key} & \emph{requis} & Clé de signature HS256\\
\texttt{jwt\_algorithm} & \texttt{HS256} & Algorithme JWT\\
\texttt{jwt\_expire\_minutes} & \texttt{60} & Durée de vie du jeton\\
\texttt{groq\_api\_key} & \emph{requis} & Clé de l'API Groq\\
\texttt{anthropic\_api\_key} & \texttt{""} & Déclarée, non utilisée dans le chemin actif\\
\texttt{chroma\_host} / \texttt{chroma\_port} & \texttt{localhost} / \texttt{8001} & Base vectorielle\\
\texttt{storage\_vm\_host} & \texttt{None} & Hôte SSH de la VM Storage\\
\texttt{storage\_vm\_port} & \texttt{22} & Port SSH\\
\texttt{storage\_vm\_user} & \texttt{None} & Compte SSH\\
\texttt{storage\_vm\_ssh\_key\_path} & \texttt{None} & Chemin de la clé privée\\
\texttt{upload\_dir} & \texttt{uploads} & Racine des fichiers déposés\\
\texttt{embedding\_model} & \texttt{intfloat/multilingual-e5-small} & Modèle d'embedding\\
\texttt{embedding\_query\_prefix} & \texttt{"query: "} & Préfixe E5 pour les questions\\
\texttt{embedding\_passage\_prefix} & \texttt{"passage: "} & Préfixe E5 pour les documents\\
\bottomrule
\end{longtable}

\begin{keybox}
\textbf{Décision de conception explicite dans le code.} Les quatre variables
SSH sont \emph{optionnelles}. Le commentaire l'assume : « absents,
\file{capacity.py} renvoie une erreur explicite au moment de l'appel plutôt que
de faire planter tout le backend au démarrage pour qui n'a pas encore la VM
configurée ».
\end{keybox}

% ---------------------------------------------------------------------
\section{Sécurité --- \file{app/core/security.py}}

\subsection*{Fonctions}

\begin{longtable}{@{}p{4.6cm}p{4.2cm}p{5.9cm}@{}}
\toprule
\textbf{Fonction} & \textbf{Paramètres} & \textbf{Retour et comportement}\\
\midrule
\endhead
\fn{hash\_password} & \texttt{plain\_password: str} & \texttt{str} --- empreinte bcrypt\\
\fn{verify\_password} & \texttt{plain}, \texttt{hashed} & \texttt{bool}\\
\fn{create\_access\_token} & \texttt{subject: str}, \texttt{role: str} & \texttt{str} --- JWT contenant \texttt{sub}, \texttt{role}, \texttt{exp}\\
\fn{decode\_access\_token} & \texttt{token: str} & \texttt{dict} ou \texttt{None} si \cls{JWTError} (jeton invalide \emph{ou} expiré)\\
\fn{create\_cross\_domain\_token} & \texttt{subject}, \texttt{workspace}, \texttt{minutes=30} & \texttt{str} --- JWT avec \texttt{scope="cross\_domain\_read\_only"} et \texttt{workspace}\\
\bottomrule
\end{longtable}

\subsection*{Gestion des erreurs}
\fn{decode\_access\_token} ne lève jamais : elle attrape \cls{JWTError} et
renvoie \texttt{None}. C'est l'appelant (\fn{get\_current\_user}) qui transforme
ce \texttt{None} en HTTP 401. Ce choix évite de disperser la gestion
d'exception JWT dans toutes les routes.

% ---------------------------------------------------------------------
\section{Dépendances FastAPI --- \file{app/api/deps.py}}

\subsection*{\fn{get\_current\_user} : décomposition pas à pas}

\begin{enumerate}
  \item \texttt{HTTPBearer()} extrait l'en-tête \texttt{Authorization: Bearer
  <token>}. En son absence, FastAPI répond 403 avant même d'entrer dans la
  fonction.
  \item \fn{decode\_access\_token(token)} : si le résultat est \texttt{None},
  levée d'une \cls{HTTPException} 401 avec
  \texttt{WWW-Authenticate: Bearer} et le message
  « Token invalide ou expire ».
  \item \texttt{payload.get("sub")} donne l'e-mail. \textbf{Accès base :}
  \texttt{db.query(User).filter(User.email == email).first()}.
  \item Si l'utilisateur est absent \emph{ou} \texttt{is\_active} est faux :
  401 « Utilisateur introuvable ou desactive ».
  \item Retour : l'entité \cls{User} rattachée à la session de la requête, donc
  \texttt{current\_user.role.name} est chargeable par relation.
\end{enumerate}

\begin{infobox}
\textbf{Détail important pour la soutenance.} Le rôle est présent dans le jeton,
mais \fn{require\_role} ne le lit \emph{pas} dans le jeton : il lit
\texttt{current\_user.role.name}, donc la base. Conséquence : changer le rôle
d'un utilisateur en base prend effet immédiatement, sans attendre l'expiration
de son jeton.
\end{infobox}

\subsection*{\fn{require\_role} : une fabrique de dépendances}

\begin{lstlisting}[language=Python,caption={api/deps.py --- contrôle de rôle}]
def require_role(*allowed_roles: str):
    def role_checker(current_user: User = Depends(get_current_user)) -> User:
        if current_user.role.name not in allowed_roles:
            raise HTTPException(status_code=403, detail="Acces refuse pour ce role")
        return current_user
    return role_checker
\end{lstlisting}

\noindent
Chaque routeur crée sa dépendance une seule fois, au niveau du module :
\texttt{\_storage\_access = require\_role("storage\_engineer", "administrator")}.

% ---------------------------------------------------------------------
\section{Authentification --- \file{api/v1/auth.py} et \file{services/auth\_service.py}}

\subsection*{Endpoints}

\begin{longtable}{@{}p{3.4cm}p{2.4cm}p{4.2cm}p{4.5cm}@{}}
\toprule
\textbf{Endpoint} & \textbf{Auth} & \textbf{Entrée} & \textbf{Sortie}\\
\midrule
\endhead
\route{POST /auth/login} & aucune & \cls{LoginRequest} (\texttt{EmailStr}, \texttt{password}) & \cls{TokenResponse} (\texttt{access\_token}, \texttt{token\_type="bearer"})\\
\route{GET /auth/me} & \fn{get\_current\_user} & --- & \cls{UserResponse} (\texttt{id}, \texttt{email}, \texttt{full\_name}, \texttt{role})\\
\bottomrule
\end{longtable}

\subsection*{Chaîne d'appels}
\route{POST /auth/login} \arrowto{} \fn{auth\_service.login(db, email, password)}
\arrowto{} \fn{authenticate\_user} \arrowto{} requête SQL sur \tbl{users}
\arrowto{} \fn{verify\_password} \arrowto{} \fn{create\_access\_token}.

\subsection*{Gestion des erreurs}
\fn{login} renvoie \texttt{None} dans trois cas indistincts : e-mail inconnu,
mot de passe faux, compte désactivé. La route les traduit tous en un seul 401
« Email ou mot de passe incorrect ». C'est délibéré : distinguer les cas
révélerait quels comptes existent.

\subsection*{Note sur \cls{UserResponse}}
\file{admin\_routes.py} construit ce schéma \textbf{à la main} et le commente :
« \texttt{UserResponse.role} attend une chaîne, alors que \texttt{User.role} est
une relation SQLAlchemy (objet \cls{Role}) ». Le mapping automatique
\texttt{from\_attributes} échouerait ici.

% ---------------------------------------------------------------------
\section{Le chat --- \file{api/v1/storage.py}}

\subsection*{Endpoints}

\begin{longtable}{@{}p{6.0cm}p{8.8cm}@{}}
\toprule
\textbf{Endpoint} & \textbf{Comportement}\\
\midrule
\endhead
\route{GET /storage/chat/conversations} & Fils de l'utilisateur, le plus récemment utilisé en tête\\
\route{GET /storage/chat/conversations/\{id\}/messages} & Messages du fil, avec leurs \cls{ActionRequestOut}\\
\route{PATCH /storage/chat/conversations/\{id\}} & Renomme (sans toucher à \texttt{updated\_at})\\
\route{DELETE /storage/chat/conversations/\{id\}} & 204 ; les messages suivent en cascade\\
\route{POST /storage/chat} & Un tour de dialogue complet\\
\bottomrule
\end{longtable}

\subsection*{Isolation par utilisateur}
Toutes ces routes passent par \fn{\_require\_conversation}, qui appelle
\fn{chat\_service.get\_conversation(db, user.id, conversation\_id)} --- filtré
sur \texttt{user\_id}. Le commentaire justifie le code de retour : « Le fil
demandé, ou 404 --- y compris s'il appartient à un collègue : dire "interdit"
plutôt qu'"introuvable" confirmerait son existence. »

\subsection*{\route{POST /storage/chat} : sept étapes}

\begin{enumerate}
  \item \textbf{Validation} --- \cls{ChatRequest} : \texttt{message: str},
  \texttt{conversation\_id: int | None}.
  \item \textbf{Résolution du fil} --- si \texttt{conversation\_id} est fourni,
  \fn{\_require\_conversation} puis \fn{chat\_service.history\_for\_llm} (10
  derniers messages).
  \item \textbf{Contexte d'incident} --- si le fil porte un
  \texttt{incident\_ticket\_id}, on recharge le ticket et on construit un
  \cls{IncidentContext} via \fn{incident\_investigation.context\_for}.
  \item \textbf{Appel de l'agent} ---
  \fn{await storage\_persona.handle\_message(message, created\_by=user.id,
  history=..., incident=...)}.
  \item \textbf{Création différée du fil} --- \emph{après} la réponse
  seulement. Commentaire : « Rien n'est écrit avant la réponse : si l'agent
  échoue, l'exception remonte sans laisser dans la liste un fil vide. »
  \item \textbf{Écriture du tour} --- \fn{chat\_service.save\_turn} écrit les
  deux messages dans la même transaction, puis \fn{attach\_actions} relie les
  \tbl{action\_requests} au message de l'assistant.
  \item \textbf{Réponse} --- \cls{ChatTurnOut} : \texttt{response},
  \texttt{conversation}, \texttt{user\_message}, \texttt{assistant\_message}.
\end{enumerate}

% ---------------------------------------------------------------------
\section{Le cœur de l'IA --- \file{app/personas/storage/agent.py}}

C'est le fichier le plus dense du projet (445 lignes). Il est décomposé
intégralement.

\subsection{Structures de données}

\subsubsection*{\cls{IncidentContext} (dataclass)}
Champs : \texttt{ticket\_id}, \texttt{ticket\_number}, \texttt{title},
\texttt{description}, \texttt{severity}, \texttt{affected\_system},
\texttt{mcp\_data}.

\begin{keybox}
\textbf{Pourquoi une dataclass et non l'entité ORM ?} Le commentaire l'explique :
« l'agent ouvre ses propres sessions et vit plus longtemps que celle de la
requête HTTP --- une entité détachée y deviendrait inutilisable à la première
lecture d'attribut ». C'est un problème classique de SQLAlchemy
(\cls{DetachedInstanceError}), résolu par la copie des champs.
\end{keybox}

\noindent
\textbf{Propriété \fn{technical\_id}} --- cherche successivement
\texttt{job\_id}, \texttt{volume\_id}, \texttt{mountpoint} dans
\texttt{mcp\_data}. Elle distingue l'\emph{identifiant technique}
(\texttt{backup-app-critical.timer}) du \emph{libellé affiché}
(\texttt{Sauvegarde de l'application critique}, qui est la \texttt{Description}
systemd). Confondre les deux ferait échouer l'action au moment précis où
l'ingénieur la confirme.

\noindent
\textbf{Méthode \fn{as\_prompt\_block}} --- produit le bloc de prompt injecté en
mode investigation : titre, ressource, cible technique, constat initial, relevé
MCP en JSON, puis trois consignes numérotées (lire les preuves, énoncer la cause
racine, proposer \emph{une} remédiation en appelant réellement l'outil).

\subsubsection*{\cls{AgentTurn} (dataclass)}
\texttt{text: str} et \texttt{proposed\_action\_ids: list[uuid.UUID]}. Le
commentaire dit pourquoi ce type existe : « \fn{handle\_message} renvoyait une
simple chaîne. Les propositions d'action doivent remonter jusqu'à la route pour
être rattachées au message qui les porte --- sans quoi l'interface ne saurait
pas quels boutons afficher. »

\subsection{\fn{\_propose\_action} --- le garde-fou central}

\textbf{Paramètres} : \texttt{db}, \texttt{tool\_name}, \texttt{arguments},
\texttt{created\_by}, \texttt{incident}.\\
\textbf{Retour} : \texttt{tuple[dict, uuid.UUID | None]} --- le dictionnaire est
réinjecté dans la conversation comme « résultat » de l'outil, l'UUID sert à
rattacher la proposition au message.

\begin{enumerate}
  \item \textbf{Extraction de la cible} parmi \texttt{("target", "volume\_id",
  "job\_id")}. Absente : retour \texttt{\{"status": "error", ...\}} et pas
  d'action créée.
  \item \textbf{Correction de la cible} : si \texttt{target ==
  incident.affected\_system} et qu'un \fn{technical\_id} existe, on substitue
  l'identifiant technique, avec journalisation
  (\texttt{"Cible de \%s corrigee : \%r -> \%r"}).
  \item \textbf{Nettoyage} : la clé \texttt{confirm} est retirée des paramètres.
  Elle sera ajoutée par \fn{confirm\_and\_execute\_action}, et la laisser ici la
  passerait deux fois.
  \item \textbf{Écriture} : \fn{incident\_service.create\_action\_request(...)}
  \arrowto{} une ligne \tbl{action\_requests} au statut \texttt{PENDING}.
  \item \textbf{Retour au modèle} : un dictionnaire dont le message est
  volontairement directif ---
  \emph{« Action NOT executed: it has been proposed to the engineer, who
  confirms or rejects it from the interface. Present it as a recommendation,
  never as a fact. »}
\end{enumerate}

\subsection{\fn{handle\_message} --- la boucle de tool calling}

\textbf{Signature} : \fn{async handle\_message(user\_message, created\_by=None,
history=None, incident=None) -> AgentTurn}.

\begin{figure}[h!]
\centering
\begin{tikzpicture}[node distance=5mm and 9mm]
\node[bluebox,minimum width=34mm] (rag) {1. RAG\\\tiny retriever.retrieve(storage\_kb, question, n=3)};
\node[bluebox,below=of rag,minimum width=34mm] (prompt) {2. Prompt système\\\tiny STORAGE\_SYSTEM\_PROMPT + date + contexte RAG + bloc incident};
\node[bluebox,below=of prompt,minimum width=34mm] (hist) {3. Historique\\\tiny seuls les tours user/assistant sont rejoués};
\node[ambox,below=of hist,minimum width=34mm] (open) {4. \texttt{async with StorageMCPClient()}\\\tiny list\_tools\_for\_llm()};
\node[vibox,below=of open,minimum width=34mm] (loop) {5. Boucle (max 6 tours)\\\tiny groq\_client.chat\_with\_tools(messages, tools)};
\node[greenbox,right=16mm of loop,minimum width=30mm] (read) {5a. Outil de lecture\\\tiny mcp.call\_tool(...)};
\node[redbox,below=6mm of read,minimum width=30mm] (rem) {5b. Outil de remédiation\\\tiny \_propose\_action(...)};
\node[bluebox,below=16mm of loop,minimum width=34mm] (rules) {6. evaluate\_tool\_result\\\tiny détection d'anomalie};
\node[bluebox,below=of rules,minimum width=34mm] (recover) {7. Rattrapage\\\tiny \_recover\_written\_call()};
\node[bluebox,below=of recover,minimum width=34mm] (tickets) {8. Création des tickets\\\tiny create\_ticket\_from\_diagnosis};
\node[box,below=of tickets,minimum width=34mm,draw=oasisdark] (out) {\textbf{AgentTurn}\\\tiny text + proposed\_action\_ids};

\draw[flowb] (rag)--(prompt); \draw[flowb] (prompt)--(hist);
\draw[flowb] (hist)--(open);  \draw[flowb] (open)--(loop);
\draw[flow] (loop)--(read);   \draw[flow] (loop)-|(rem);
\draw[flow] (read.east) to[bend right=45] node[lbl,right]{résultat} (rules.east);
\draw[flow] (rem.east) to[bend right=35] (rules.east);
\draw[flowb] (loop)--(rules); \draw[flowb] (rules)--(recover);
\draw[flowb] (recover)--(tickets); \draw[flowb] (tickets)--(out);
\end{tikzpicture}
\caption{Les huit temps de \fn{handle\_message}.}
\end{figure}

\clearpage

\subsubsection*{Étape 1 --- RAG}
\fn{retriever.retrieve(collection\_name="storage\_kb", question=user\_message,
n\_results=3)} \arrowto{} \fn{get\_context\_for\_query} \arrowto{}
\fn{vectorstore.search} \arrowto{} \fn{embed\_query} (préfixe
\texttt{"query: "}) \arrowto{} \texttt{collection.query}. Le résultat est une
chaîne où chaque passage est préfixé de \texttt{[Source: <fichier>]}.

\subsubsection*{Étape 2 --- Prompt système}
Il concatène quatre blocs : le \fn{STORAGE\_SYSTEM\_PROMPT}, la date et l'heure
du serveur, une consigne explicite sur MCP, et le contexte RAG. En mode
investigation s'ajoute \fn{incident.as\_prompt\_block()}.

\begin{lstlisting}[caption={Consigne anti-hallucination injectée à chaque tour}]
Use MCP only for current storage information. The tools that change the VM
(run_backup, create_snapshot, restore_from_backup) never execute from the chat:
calling one records an action proposal that the engineer confirms with a click.
So call them normally when you want to propose one, but never copy a tool call
as JSON into your answer -- written that way it creates nothing. And never say
a remediation has been launched: it is waiting on a human decision.
\end{lstlisting}

\subsubsection*{Étape 3 --- Historique}
Seuls les messages de rôle \texttt{user} et \texttt{assistant} sont rejoués. Le
commentaire justifie : « les appels d'outils d'une conversation précédente ne
sont pas rejouables (leurs résultats ne sont plus à jour), le modèle les
redemande s'il en a besoin ».

\subsubsection*{Étape 4 --- Session MCP et session d'audit}
\texttt{async with StorageMCPClient() as mcp} démarre le sous-processus.
\texttt{audit\_db = SessionLocal() if created\_by is not None else None} : sans
utilisateur identifié, on n'audite pas plutôt que d'inventer un auteur.

\subsubsection*{Étape 5 --- La boucle (au plus \fn{\_MAX\_TOOL\_ROUNDS} = 6)}

\begin{enumerate}
  \item \fn{groq\_client.chat\_with\_tools(messages, tools)}.
  \item Si \texttt{reply.tool\_calls} est vide \arrowto{} \texttt{final\_text =
  reply.content}, sortie de boucle.
  \item Sinon, le message assistant \emph{avec ses \texttt{tool\_calls}} est
  réinjecté dans \texttt{messages} (indispensable au protocole OpenAI/Groq).
  \item Pour chaque appel : \texttt{json.loads(arguments)} dans un
  \texttt{try/except JSONDecodeError} qui retombe sur \texttt{\{\}}.
  \item Aiguillage : \texttt{REMEDIATION\_TOOLS} \arrowto{} proposition ;
  sinon \arrowto{} exécution MCP réelle, avec \texttt{RuntimeError} traduite en
  \texttt{\{"status": "error", "message": ...\}} plutôt que propagée.
  \item Le résultat est réinjecté sous forme d'un message de rôle
  \texttt{tool}, avec le \texttt{tool\_call\_id} correspondant.
\end{enumerate}

\noindent
Si la boucle atteint 6 tours sans réponse finale, un dernier appel sans outils
produit le texte, avec un message de repli explicite (« MCP collection limit
reached; try a more targeted question »).

\subsubsection*{Étape 6 --- Détection déterministe}
\fn{evaluate\_tool\_result(name, result)} applique les règles de
\file{incident\_rules.py}. \textbf{Aucune décision d'incident n'est prise par le
modèle.}

\subsubsection*{Étape 7 --- \fn{\_recover\_written\_call} : le rattrapage}

C'est une des idées les plus fines du projet. Le défaut visé est décrit dans le
code : « le diagnostic est juste, la remédiation est la bonne, mais elle est
recopiée en JSON dans la réponse. Écrite ainsi elle ne crée rien, et l'ingénieur
lit une analyse complète sans avoir aucun bouton à cliquer ».

\begin{enumerate}
  \item Une expression régulière \fn{\_JSON\_OBJECT} extrait les objets JSON du
  texte (un niveau d'imbrication toléré).
  \item \fn{\_candidate\_calls} accepte trois formes rencontrées selon les
  modèles : \texttt{\{"tool", "args"\}}, \texttt{\{"name", "arguments"\}},
  \texttt{\{"tool", "parameters"\}}, plus la forme où le nom de l'outil est
  directement une clé.
  \item \textbf{Double validation} : l'outil doit appartenir à
  \texttt{REMEDIATION\_TOOLS}, \emph{et} la cible doit être
  \texttt{incident.affected\_system} ou \texttt{incident.technical\_id}. Une
  cible inventée ou portant sur une autre machine est ignorée.
  \item Ce qui en sort reste une proposition \texttt{pending}, soumise au même
  clic humain.
\end{enumerate}

\subsubsection*{Étape 8 --- Création des tickets détectés}
Une nouvelle session (\texttt{with SessionLocal() as db}) écrit les tickets via
\fn{create\_ticket\_from\_diagnosis}. Sans \texttt{created\_by}, on journalise un
avertissement et on ne crée rien.

\begin{warnbox}
\textbf{Gestion d'erreurs à souligner.} Chaque bloc à risque est entouré d'un
\texttt{try/except} qui journalise via \fn{logger.exception} et \textbf{ne fait
jamais échouer le tour}. Une panne d'audit, de création de ticket ou de
rattrapage ne doit pas priver l'ingénieur de sa réponse.
\end{warnbox}

% ---------------------------------------------------------------------
\section{Client LLM --- \file{app/llm/groq\_client.py}}

\subsection*{Classe \cls{GroqClient}}
\texttt{self.client = Groq(api\_key=settings.groq\_api\_key)} ;
\texttt{self.model = getattr(settings, "groq\_model", "openai/gpt-oss-120b")}.

\begin{infobox}
\texttt{groq\_model} n'est \textbf{pas} un champ de \cls{Settings} : le
\fn{getattr} retombe donc toujours sur le défaut
\texttt{"openai/gpt-oss-120b"}. C'est une lecture du code, pas une supposition.
\end{infobox}

\subsection*{\fn{chat\_with\_tools(messages, tools=None)}}
Construit une liste de modèles candidats
(\texttt{openai/gpt-oss-120b}, \texttt{openai/gpt-oss-20b},
\texttt{llama-3.3-70b-versatile}), dédoublonnée en conservant l'ordre, puis les
essaie successivement. À la première réussite, \texttt{self.model} est mis à
jour --- le modèle qui fonctionne devient le défaut pour les appels suivants. Si
tous échouent, la dernière exception est relevée.

\noindent
Options envoyées : \texttt{temperature=0.3} et, si des outils sont fournis,
\texttt{tools} et \texttt{tool\_choice="auto"}.

\subsection*{\fn{chat(prompt, system\_prompt=None)}}
Enveloppe simple qui renvoie \texttt{.content}. Utilisée par
\fn{incident\_investigation.\_summarise} pour la synthèse
\texttt{root\_cause} / \texttt{impact\_summary}.

% ---------------------------------------------------------------------
\section{Actions sensibles --- \file{api/v1/actions.py}}

\subsection*{Endpoints}

\begin{longtable}{@{}p{4.6cm}p{2.6cm}p{7.5cm}@{}}
\toprule
\textbf{Endpoint} & \textbf{Auth} & \textbf{Comportement}\\
\midrule
\endhead
\route{GET /storage/actions} & \fn{get\_current\_user} & Actions de l'utilisateur courant (filtre \texttt{requested\_by}), pagination \texttt{limit} (\(\le 200\)) et \texttt{offset}\\
\route{POST /storage/actions/} & \fn{get\_current\_user} & Enregistre une proposition (\texttt{pending}). \cls{ValueError} \arrowto{} 400\\
\route{POST /storage/actions/\{id\}/confirm} & \fn{get\_current\_user} & \textbf{Seul chemin d'exécution réelle}\\
\route{POST /storage/actions/\{id\}/reject} & \fn{get\_current\_user} & Passe en \texttt{REJECTED} et écrit une note sur le ticket\\
\route{GET /storage/actions/\{id\}} & \fn{get\_current\_user} & Lecture d'une action\\
\bottomrule
\end{longtable}

\subsection*{\fn{confirm\_action} pas à pas}

\begin{enumerate}
  \item \texttt{mcp = StorageMCPClient()} puis \texttt{await mcp.start()}. Si le
  sous-processus ne démarre pas : \textbf{HTTP 503}, et l'action \emph{reste}
  \texttt{PENDING}, donc rejouable.
  \item \fn{await incident\_service.confirm\_and\_execute\_action(db,
  action\_id, user\_id, mcp\_client=mcp)}.
  \item \fn{await \_close\_the\_loop(db, action, user, mcp)}.
  \item \texttt{finally: await mcp.stop()} --- le sous-processus est toujours
  fermé.
\end{enumerate}

\begin{keybox}
\textbf{Le correctif documenté dans le fichier.} Le docstring raconte un bug
réel : « cette route utilisait une instance module-level de
\cls{StorageMCPClient} dont personne n'appelait jamais \fn{start()} \ldots{}
\fn{call\_tool()} levait donc systématiquement "Storage MCP indisponible",
l'action retombait en FAILED et l'API répondait quand même 200 : le seul chemin
autorisé pour exécuter une action destructive ne fonctionnait pas. »
C'est un excellent exemple à citer en soutenance.
\end{keybox}

\subsection*{\fn{\_close\_the\_loop} : la fermeture du cycle}

\begin{enumerate}
  \item Sans \texttt{incident\_ticket\_id} \arrowto{} rien à faire.
  \item Si l'action est \texttt{FAILED} \arrowto{} note
  « \ldots{} confirmed by X but failed: <raison> », et le ticket reste ouvert.
  \item Sinon \arrowto{} \fn{await incident\_verification.verify(ticket,
  mcp=mcp, db=db, user\_id=user.id)}.
  \item Une note consigne le verdict ; si \texttt{verdict.resolved} est vrai,
  \fn{incident\_service.resolve\_ticket}.
  \item \texttt{except Exception} global : la fermeture du cycle ne doit
  \textbf{jamais} faire échouer la requête. « L'action, elle, a bien eu lieu, et
  la surveillance périodique repassera vérifier. »
\end{enumerate}

% ---------------------------------------------------------------------
\section{Service transactionnel --- \file{services/incident\_service.py}}

\subsection*{\fn{confirm\_and\_execute\_action} : le point critique}

\begin{lstlisting}[language=Python,caption={Extrait réel --- transition, exécution, qualification du résultat}]
if action.status != ActionRequestStatus.PENDING:
    raise ValueError(f"Action already handled (current status: {action.status.value})")

action.status = ActionRequestStatus.CONFIRMED
action.confirmed_by = user_id
action.confirmed_at = datetime.utcnow()
db.commit()

try:
    result = await mcp_client.call_tool(
        action.action_type, db=db, user_id=user_id,
        action_request_id=action.id, confirm=True, **action.parameters,
    )
    if isinstance(result, dict) and result.get("error"):
        action.status = ActionRequestStatus.FAILED
    else:
        action.status = ActionRequestStatus.COMPLETED
    action.result = result
except Exception as error:
    action.status = ActionRequestStatus.FAILED
    action.result = {"error": str(error)}
\end{lstlisting}

\begin{keybox}
\textbf{Trois protections dans dix lignes.}

\emph{(1)} Le test \texttt{status != PENDING} empêche la double exécution : deux
clics sur « Confirm » ne lancent pas deux restaurations.

\emph{(2)} Le passage en \texttt{CONFIRMED} est \texttt{commit} \textbf{avant}
l'appel, donc la trace de l'intention survit à un crash pendant l'exécution.

\emph{(3)} Le test \texttt{result.get("error")} est essentiel : les outils
Storage ne lèvent pas d'exception quand la VM refuse une commande, ils renvoient
\texttt{\{"error": ...\}}. Sans ce test, \texttt{systemctl start} sur un job qui
échoue aussitôt marquerait l'action \texttt{COMPLETED}, et
\fn{has\_completed\_action} prendrait une remédiation ratée pour une remédiation
faite --- jusqu'à clore le ticket sur cette base.
\end{keybox}

\subsection*{\fn{\_generate\_ticket\_number}}
Le numéro suivant est dérivé du \textbf{maximum} existant, pas d'un
\texttt{COUNT}. Justification donnée : « compter suppose une suite sans trou. Un
seul ticket supprimé, et le compteur repointe sur un numéro déjà pris ---
l'insertion viole alors l'unicité de \texttt{ticket\_number} et fait tomber tout
le passage de surveillance, donc la détection elle-même. » Format :
\texttt{INC-<année>-0001}.

\subsection*{\fn{\_open\_ticket\_for} : déduplication}
Recherche d'abord par identifiant technique dans le JSONB
(\texttt{mcp\_data["job\_id"]}, puis \texttt{mcp\_data["volume\_id"]}), et
\emph{seulement ensuite} par \texttt{affected\_system}. Motif : renommer l'unit
sur la VM suffisait à faire ouvrir un second ticket pour le même job.

\subsection*{\fn{has\_completed\_action}}
Retourne vrai s'il existe au moins une \tbl{action\_requests} en
\texttt{COMPLETED} sur ce ticket. C'est la garde de l'auto-résolution : « un
incident ne se referme pas tout seul parce que la VM va mieux, mais parce
qu'une action confirmée par un humain a produit cet effet ».

% ---------------------------------------------------------------------
\section{Investigation --- \file{services/incident\_investigation.py}}

\subsection*{\fn{start(db, ticket, user)} --- séquence}

\begin{enumerate}
  \item \fn{find\_conversation(db, ticket\_id, user\_id)} : un fil existant est
  renvoyé tel quel, sans relancer l'agent. Le fil est \textbf{par utilisateur},
  car une conversation appartient à une personne et n'est lisible que par elle.
  \item \fn{context\_for(ticket)} \arrowto{} \cls{IncidentContext}.
  \item \fn{\_investigation\_prompt(ticket)} rédige la question à la place de
  l'ingénieur --- et elle est enregistrée comme un vrai message utilisateur,
  « en relisant le fil six mois plus tard, on voit la question qui a lancé
  l'analyse, et non un diagnostic tombé de nulle part ».
  \item L'agent tourne \textbf{avant} toute écriture.
  \item Création du fil, rattachement à \texttt{incident\_ticket\_id},
  \fn{save\_turn}.
  \item Si l'agent n'a proposé aucune action \arrowto{} \fn{\_fallback\_action}.
  \item \fn{attach\_actions}, puis \fn{mark\_in\_progress} (le ticket passe de
  \texttt{OPEN} à \texttt{IN\_PROGRESS}).
  \item \fn{\_summarise(turn.text)} : second appel Groq qui renvoie un JSON
  \texttt{\{"root\_cause", "impact\_summary"\}} recopié sur le ticket.
  \item Note d'audit nommée, puis retour \texttt{\{"conversation\_id",
  "created"\}}.
\end{enumerate}

\subsection*{\fn{\_fallback\_action} : filet de sécurité déterministe}
Parcourt \texttt{ticket.recommendations} (calculées par
\file{incident\_rules.py} à la détection), retient la première dont
\texttt{action\_type} appartient à \texttt{REMEDIATION\_TOOLS}, écarte les
doublons (\texttt{pending}, \texttt{confirmed}, \texttt{completed}) et crée
l'action.

\begin{infobox}
Le commentaire résume la philosophie : « on ne demande donc au modèle que ce que
lui seul sait faire, le diagnostic, et on tient la proposition de ce qui est
déterministe ».
\end{infobox}

\subsection*{\fn{\_parse\_summary}}
Découpe entre la première \texttt{\{} et la dernière \texttt{\}} avant
\texttt{json.loads}, pour survivre à un modèle qui aurait ajouté un bloc de code
markdown. Ne conserve que les deux clés attendues.

% ---------------------------------------------------------------------
\section{Vérification --- \file{services/incident\_verification.py}}

\subsection*{Classe \cls{Verdict}}
Trois attributs : \texttt{resolved: bool}, \texttt{detail: str},
\texttt{conclusive: bool}. Le troisième distingue « toujours en panne » de
« état pas encore lisible » (job encore en cours).

\subsection*{\fn{verify(ticket, mcp, db, user\_id)} --- aiguillage}
Le type d'incident se déduit de \texttt{ticket.mcp\_data} :
\texttt{job\_id} \arrowto{} \fn{\_verify\_backup} ;
\texttt{volume\_id} ou \texttt{mountpoint} \arrowto{} \fn{\_verify\_capacity} ;
sinon \arrowto{} verdict non concluant, clôture manuelle.

\subsection*{\fn{\_verify\_backup} --- boucle d'attente}
5 tentatives espacées de 3 secondes tant que le job est \texttt{running}, car
« un job systemd de type \texttt{oneshot} est encore \texttt{running} juste
après un \texttt{systemctl start} ». Puis la \textbf{même règle} que la
détection est réappliquée : si \fn{check\_backup\_status(job)} renvoie
\texttt{None}, le problème a disparu.

\begin{keybox}
\textbf{Point d'architecture élégant.} La vérification ne réécrit pas une
variante des seuils : elle réutilise \fn{check\_capacity} et
\fn{check\_backup\_status}. « Ce qui a ouvert le ticket est exactement ce qui
doit avoir disparu pour le fermer. »
\end{keybox}

% ---------------------------------------------------------------------
\section{Surveillance de fond --- \file{services/incident\_watcher.py}}

\subsection*{\fn{scan\_once()} --- un passage, sans exception}

\begin{enumerate}
  \item \fn{list\_backups()} --- appel \textbf{direct}, pas via MCP.
  \item Si \texttt{error} \arrowto{} rapport
  \texttt{\{"error", "created": [], "skipped": [], "closed": []\}}. Une VM
  éteinte est « une situation normale, pas une panne de l'application ».
  \item \fn{\_system\_user\_id(db)} --- compte de service.
  \item \fn{\_verify\_resolved} \textbf{d'abord} : referme les tickets dont la
  remédiation a produit son effet. L'ordre est délibéré, « un ticket clos au
  même cycle ne doit pas être rouvert par la boucle de détection qui suit ».
  \item Pour chaque job en échec : \fn{check\_backup\_status}, test
  \fn{\_already\_handled}, puis \fn{create\_ticket\_from\_diagnosis} et une
  première note d'audit.
\end{enumerate}

\subsection*{Le compte de service}
\texttt{oasis-system@oasis.local}, rôle \texttt{storage\_engineer},
\texttt{is\_active=False}, mot de passe aléatoire jamais communiqué. Il ne peut
pas servir à se connecter (\fn{get\_current\_user} rejette les comptes
désactivés) : il n'existe que pour être l'auteur des tickets automatiques, car
\texttt{incident\_tickets.created\_by} est \texttt{NOT NULL}.

\subsection*{\fn{\_already\_handled} --- anti-rouvertures}
Sans ce test, résoudre manuellement un ticket sur un job encore en échec en
rouvrirait un identique au cycle suivant, « le bouton résoudre n'aurait aucun
effet visible ». Règle : on ne rouvre qu'à partir d'une exécution
\emph{postérieure} au ticket précédent --- un ticket par échec, pas par cycle.

% ---------------------------------------------------------------------
\section{Dashboard --- \file{app/dashboard/kpi\_service.py}}

\subsection*{Objectif}
Agréger, \textbf{sans aucun appel LLM}, la santé de la VM et les compteurs de la
base. Le fichier le déclare : « Aucun appel LLM ici : le dashboard doit rester
rapide et déterministe. »

\subsection*{Le cache de santé : trois fonctions coopérantes}

\begin{longtable}{@{}p{4.6cm}p{10.2cm}@{}}
\toprule
\textbf{Fonction} & \textbf{Comportement}\\
\midrule
\endhead
\fn{refresh\_health\_snapshot} & Prend \fn{\_refresh\_lock}, revérifie la fraîcheur, appelle \fn{get\_storage\_health()}, écrit le cache sous \fn{\_health\_lock}\\
\fn{\_storage\_health(force)} & Sert le dernier relevé ; s'il dépasse 60\,s et qu'aucun rafraîchissement n'est en cours, lance un thread de fond et rend la main immédiatement\\
\fn{\_refresh\_in\_background} & Cible du thread, avale les exceptions (la VM peut être injoignable)\\
\fn{\_invalidate\_health\_snapshot} & Vide le cache, utilisé par \texttt{?refresh=true}\\
\bottomrule
\end{longtable}

\begin{infobox}
\textbf{Motif documenté.} « \fn{get\_storage\_health()} ouvre deux connexions SSH
vers la VM Storage : la réponse met plusieurs secondes, et le dashboard la
demandait à chaque affichage --- c'est ce qui rendait le changement de page si
lent. » Le verrou évite qu'« dix onglets ouverts en même temps sur un cache vide
déclenchent dix sessions SSH pour la même donnée ».
\end{infobox}

\subsection*{Constructeurs de blocs}

\begin{itemize}
  \item \fn{\_build\_capacity} --- totalise les volumes, calcule
  \texttt{reserved\_gb} (l'écart entre total et \(used + available\)
  correspondant aux blocs réservés au superutilisateur), attribue un statut par
  volume aux seuils 80\,\% / 90\,\%.
  \item \fn{\_build\_backups} --- compte
  \texttt{successful} / \texttt{failed} / \texttt{running} / \texttt{scheduled},
  et distingue \texttt{configured=False} (aucun timer) de \texttt{error} (VM
  injoignable).
  \item \fn{\_build\_incidents} --- agrégation SQL \texttt{GROUP BY severity} sur
  les statuts ouverts, plus le compte des résolus sur 7 jours.
  \item \fn{\_health\_score} --- \(100 - \sum \text{poids}\) avec
  \texttt{critical}=12, \texttt{high}=6, \texttt{medium}=2, \texttt{low}=1, plus
  une pénalité de capacité. Renvoie \texttt{None} si la VM est injoignable :
  « afficher 100/100 sans avoir rien pu mesurer donnerait une fausse
  assurance ».
  \item \fn{\_headline} --- une phrase en langage naturel, par ordre de priorité
  décroissante.
  \item \fn{get\_insights} --- observations avec priorité, détail, et lien vers
  l'écran qui permet d'agir.
\end{itemize}

% ---------------------------------------------------------------------
\section{Métriques --- \file{tools/performance.py} et \file{services/metrics\_service.py}}

\subsection*{\fn{sample\_disk\_performance()}}
\texttt{/proc/diskstats} n'expose que des compteurs cumulatifs : une seule
lecture ne dit rien. La fonction exécute donc, \textbf{en une seule session
SSH} :

\begin{lstlisting}[caption={La commande réellement envoyée à la VM}]
cat /proc/diskstats; echo '---'; sleep 2; cat /proc/diskstats
\end{lstlisting}

\noindent
Puis \fn{\_parse\_diskstats} filtre les périphériques virtuels
(\texttt{loop}, \texttt{ram}, \texttt{sr}, \texttt{fd}) et les partitions
(regex \fn{\_PARTITION}) --- « leurs E/S sont déjà comptées dans le disque
parent, les inclure reviendrait à compter deux fois ».
\fn{\_rates} calcule IOPS, débit (secteurs \(\times\) 512 octets) et latence
moyenne, et renvoie \texttt{None} si un compteur a reculé (redémarrage entre les
deux lectures).

\subsection*{\fn{read\_series(db, range\_key)} --- agrégation en deux passes}
\texttt{RANGE\_WINDOWS} définit fenêtre et granularité :
\texttt{24h}\,\arrowto{}\,1\,h, \texttt{7d}\,\arrowto{}\,6\,h,
\texttt{30d}\,\arrowto{}\,1\,jour. La largeur vient de cette table, jamais de
l'utilisateur --- ce qui ferme toute injection dans le \texttt{literal\_column}.

\begin{enumerate}
  \item Sous-requête : moyenne \emph{par device} et par intervalle.
  \item Requête principale : somme des devices pour IOPS et débit, moyenne pour
  la latence.
\end{enumerate}

\noindent
Justification : « agréger en une seule passe donnerait un total faussé dès
qu'un device a plus d'échantillons qu'un autre ».
La tendance \texttt{iops\_trend\_percent} compare la seconde moitié des points à
la première.

\begin{keybox}
\textbf{Le point d'honnêteté.} \fn{get\_performance} renvoie
\texttt{collecting = not series["points"]}. Tant que le collecteur n'a rien
enregistré, la série est vide et l'interface l'annonce, au lieu de fabriquer une
courbe.
\end{keybox}

% ---------------------------------------------------------------------
\section{Data Hub --- \file{services/knowledge\_service.py}}

\subsection*{\fn{create\_document} --- validations en cascade}

\begin{enumerate}
  \item Fichier vide \arrowto{} \cls{DocumentError}.
  \item Taille \(>\) \texttt{MAX\_UPLOAD\_BYTES} (10\,Mo) \arrowto{}
  \cls{DocumentError}.
  \item Extension hors de \texttt{ALLOWED\_EXTENSIONS}
  (\texttt{.md .txt .log .json .yaml .yml .csv .pdf}) \arrowto{}
  \cls{DocumentError}.
  \item Écriture sur disque sous \texttt{<uuid><extension>} --- \textbf{jamais}
  sous le nom fourni par le client. C'est ce qui rend toute traversée de
  répertoire impossible.
  \item Ligne \tbl{knowledge\_documents} en statut \texttt{PENDING}, commit.
  \item Extraction du texte, \fn{\_index\_document}, statut \texttt{INDEXED} et
  \texttt{indexed\_at}.
  \item En cas d'échec : statut \texttt{FAILED} et message tronqué à 500
  caractères dans \texttt{error}. \textbf{Le document reste en base}, l'interface
  affiche pourquoi il n'est pas indexé.
\end{enumerate}

\subsection*{\fn{\_index\_document}}
Refuse un document de moins de \texttt{MIN\_INDEXABLE\_WORDS} (20) mots utiles,
avec un message qui diagnostique le cas courant : « S'il s'agit d'un PDF scanné
ou d'une image exportée en PDF, il n'a pas de couche texte : il faut lui
appliquer un OCR avant de le déposer. » Puis \fn{chunk\_text} (500 mots,
recouvrement 50), \fn{embed\_passages}, et \texttt{collection.upsert} avec des
identifiants dérivés de l'UUID du document --- ce qui rend la suppression et la
réindexation possibles.

% ---------------------------------------------------------------------
\section{Demandes de service --- \file{services/request\_service.py}}

\subsection*{Les deux règles qui portent toute la logique}

\begin{enumerate}
  \item \textbf{Le sens est déduit du rôle de l'auteur, jamais fourni par le
  client.} Un \texttt{storage\_engineer} ne peut créer qu'une
  \texttt{ENGINEER\_TO\_ADMIN} ; un \texttt{administrator}, qu'une
  \texttt{ADMIN\_TO\_ENGINEER}. « Le client ne peut donc pas fabriquer une
  demande qui semblerait venir de l'autre camp. »
  \item \textbf{Seul le destinataire fait avancer le statut.}
  \texttt{request.assigned\_to != user.id} \arrowto{} \cls{RequestForbidden}.
\end{enumerate}

\subsection*{Table \fn{ALLOWED\_TRANSITIONS}}

\begin{longtable}{@{}p{3.0cm}p{5.4cm}p{6.2cm}@{}}
\toprule
\textbf{Statut} & \textbf{Ingénieur \arrowto{} Admin} & \textbf{Admin \arrowto{} Ingénieur}\\
\midrule
\endhead
\texttt{pending} & \texttt{in\_review}, \texttt{approved}, \texttt{rejected} & \texttt{acknowledged}, \texttt{in\_progress}\\
\texttt{in\_review} & \texttt{approved}, \texttt{rejected} & ---\\
\texttt{approved} & \texttt{completed} & ---\\
\texttt{acknowledged} & --- & \texttt{in\_progress}, \texttt{completed}\\
\texttt{in\_progress} & --- & \texttt{completed}\\
\texttt{rejected} & \(\varnothing\) (terminal) & ---\\
\texttt{completed} & \(\varnothing\) (terminal) & \(\varnothing\) (terminal)\\
\bottomrule
\end{longtable}

\noindent
Une transition refusée produit un message qui \emph{énumère} les transitions
possibles, ce qui rend l'erreur exploitable côté interface.

\subsection*{Horodatages distincts}
\texttt{decided\_at} est figé à l'arrivée dans \texttt{APPROVED} ou
\texttt{REJECTED}, distinct de \texttt{updated\_at} : « le détail annonce
"refusée le \ldots", pas "modifiée le \ldots" ».

\subsection*{Pièces jointes}
Plafond 10\,Mo, éventail d'extensions plus large que le Data Hub (« ce qui
circule ici n'a pas vocation à être indexé, seulement lu par l'autre partie »),
stockage sous UUID dans \file{upload\_dir/requests}, suppression réservée au
déposant.

% ---------------------------------------------------------------------
\section{Recherche globale --- \file{services/search\_service.py}}

\subsection*{\fn{search(db, user, query, limit\_per\_type)}}
Refuse les requêtes de moins de 2 caractères (« elle renverrait la moitié de la
base pour rien »), puis concatène quatre recherches :
incidents, demandes, conversations, documents.

\subsection*{\fn{\_pattern} --- échappement}

\begin{lstlisting}[language=Python,caption={Protection contre les jokers SQL}]
escaped = query.replace("\\", "\\\\").replace("%", "\\%").replace("_", "\\_")
return f"%{escaped}%"
\end{lstlisting}

\noindent
Puis \texttt{.ilike(pattern, escape="\textbackslash\textbackslash")}. Sans cela,
« un \texttt{\%} tapé dans la barre de recherche deviendrait un joker et
remonterait toute la base ».

\subsection*{Cloisonnement}
Les demandes et les conversations sont filtrées sur l'utilisateur courant
(\texttt{created\_by == user\_id or assigned\_to == user\_id} pour les demandes,
\texttt{user\_id} pour les conversations). Les incidents et les documents sont
communs.

\subsection*{Identifiants retournés}
Pour un incident, c'est l'\textbf{UUID} qui est renvoyé et non le numéro
affiché : « la page Incidents recharge le ticket par la route
\route{/incidents/\{id\}}, ce qui l'ouvre même quand il ne figure pas dans la
liste actuellement filtrée ».

% ---------------------------------------------------------------------
\section{Rapports PDF --- \file{services/report\_service.py}}

\subsection*{Deux entrées}
\fn{build\_global\_report(db, author)} et
\fn{build\_incident\_report(db, ticket\_id, author)} renvoient des
\texttt{bytes}. La seconde renvoie \texttt{None} si le ticket n'existe pas,
traduit en 404 par la route.

\subsection*{Mise en forme}
Le fichier définit une charte : couleurs (\texttt{\_BRAND}, \texttt{\_INK},
\texttt{\_MUTED}, \texttt{\_LINE}, \texttt{\_ZEBRA}), styles de paragraphe
(\fn{\_styles}), un constructeur de tableau (\fn{\_table}), un pied de page
(\fn{\_footer}) et un moteur de rendu (\fn{\_render}).

\subsection*{Détail HTTP à connaître}
La route ajoute \texttt{Access-Control-Expose-Headers: Content-Disposition}.
Sans cet en-tête, « le nom de fichier n'est pas lisible par le frontend, qui lit
la réponse en blob depuis un fetch authentifié ».

% ---------------------------------------------------------------------
\section{Audit HTTP --- \file{middleware/audit\_middleware.py}}

\subsection*{\fn{dispatch}}

\begin{enumerate}
  \item Court-circuit sur \texttt{EXCLUDED\_PATHS} (\file{/docs},
  \file{/openapi.json}, \file{/redoc}, \file{/favicon.ico}).
  \item Chronomètre \texttt{time.monotonic()} autour de \fn{call\_next}.
  \item \textbf{Session dédiée} \fn{SessionLocal()} --- pas celle de la requête,
  qui est déjà refermée à ce stade.
  \item \fn{\_extract\_user\_id} redécode le \texttt{Bearer} pour attribuer la
  ligne.
  \item Écriture d'une \cls{AuditLog} : \texttt{action} = méthode,
  \texttt{resource} = chemin, \texttt{details} =
  \texttt{\{status\_code, duration\_ms, client\_ip\}}.
  \item \texttt{except Exception: db.rollback()} --- un échec d'audit ne casse
  jamais la réponse.
\end{enumerate}

\begin{keybox}
\textbf{Le corps n'est jamais journalisé.} Commentaire du fichier : « on ne
journalise jamais le corps de la requête (donc jamais un mot de passe en clair
envoyé à \route{/auth/login}) --- seulement des métadonnées ». En contrepartie,
les tentatives refusées en 401 et 403 \emph{sont} enregistrées, ce qui est
« précieux pour la supervision de sécurité ».
\end{keybox}

% =====================================================================
\chapter{La base de données}
\label{ch:db}

\section{Organisation générale}

Le schéma comporte \textbf{16 tables} réparties en cinq familles, définies par
16 classes réparties sur 11 fichiers de modèle. Toutes héritent de
\fn{Base = declarative\_base()} (\file{app/database/base.py}) et sont importées
dans \file{app/models/\_\_init\_\_.py}, ce qui garantit qu'Alembic les voit
toutes.

\begin{longtable}{@{}p{3.4cm}p{11.4cm}@{}}
\toprule
\textbf{Famille} & \textbf{Tables}\\
\midrule
\endhead
Identité \& sécurité & \tbl{roles}, \tbl{users}, \tbl{audit\_logs}, \tbl{access\_requests}\\
Incidents \& actions & \tbl{incident\_tickets}, \tbl{incident\_ticket\_notes}, \tbl{action\_requests}, \tbl{mcp\_calls}\\
Conversation & \tbl{chat\_conversations}, \tbl{chat\_messages}\\
Demandes de service & \tbl{service\_requests}, \tbl{request\_messages}, \tbl{request\_events}, \tbl{request\_attachments}\\
Connaissance \& mesures & \tbl{knowledge\_documents}, \tbl{disk\_metrics}\\
\bottomrule
\end{longtable}

\begin{infobox}
\textbf{Deux conventions de clé primaire coexistent.} Les tables héritées de la
V1 (\tbl{users}, \tbl{roles}, \tbl{audit\_logs}, \tbl{access\_requests})
utilisent des entiers auto-incrémentés ; les tables métier ajoutées ensuite
utilisent des \texttt{UUID}. Les commentaires du code portent la trace de cette
transition : plusieurs modèles annoncent « \textsc{corrigé} : Integer (users.id
est un Integer, pas un UUID) », après un premier jet où les clés étrangères
vers \tbl{users} avaient été déclarées en UUID --- ce que PostgreSQL refuse.

Enfin, \tbl{chat\_conversations} et \tbl{chat\_messages} utilisent des
\texttt{BigInteger} auto-incrémentés, et ce choix est justifié : c'est
l'identifiant qui donne \emph{l'ordre du fil}, parce que les deux messages d'un
même tour partagent le \texttt{created\_at} renvoyé par \texttt{now()}.
\end{infobox}

\section{Schéma relationnel}

\begin{figure}[h!]
\centering
\begin{tikzpicture}[node distance=9mm and 13mm, every node/.style={font=\scriptsize\sffamily}]

\node[bluebox,minimum width=24mm] (roles) {\textbf{roles}\\\tiny id · name};
\node[bluebox,below=of roles,minimum width=24mm] (users) {\textbf{users}\\\tiny id · email · hashed\_password\\\tiny full\_name · is\_active · role\_id};

\node[redbox,right=20mm of users,minimum width=30mm] (tickets) {\textbf{incident\_tickets}\\\tiny id(UUID) · ticket\_number\\\tiny severity · status · mcp\_data(JSONB)\\\tiny recommendations(JSONB) · created\_by};
\node[redbox,below=of tickets,minimum width=30mm] (notes) {\textbf{incident\_ticket\_notes}\\\tiny id · ticket\_id · author\_id · content};
\node[ambox,above=of tickets,minimum width=30mm] (actions) {\textbf{action\_requests}\\\tiny id(UUID) · action\_type · target\\\tiny parameters(JSONB) · status · result(JSONB)\\\tiny requested\_by · confirmed\_by\\\tiny incident\_ticket\_id · chat\_message\_id};
\node[ambox,above=of actions,minimum width=30mm] (mcpcalls) {\textbf{mcp\_calls}\\\tiny id(UUID) · persona · tool\_name\\\tiny arguments(JSONB) · result(JSONB)\\\tiny status · correlation\_id\\\tiny user\_id · action\_request\_id};

\node[greenbox,right=20mm of tickets,minimum width=28mm] (conv) {\textbf{chat\_conversations}\\\tiny id(BigInt) · user\_id · title\\\tiny incident\_ticket\_id};
\node[greenbox,above=of conv,minimum width=28mm] (msg) {\textbf{chat\_messages}\\\tiny id(BigInt) · conversation\_id\\\tiny role · content};

\node[vibox,below=of notes,minimum width=28mm] (sreq) {\textbf{service\_requests}\\\tiny id(UUID) · direction · status\\\tiny created\_by · assigned\_to};
\node[vibox,right=20mm of sreq,minimum width=26mm] (satt) {\textbf{request\_attachments}\\\tiny id · request\_id · stored\_path};

\draw[flow] (roles)--(users);
\draw[flow] (users)--node[lbl,above]{created\_by}(tickets);
\draw[flow] (tickets)--(notes);
\draw[flow] (tickets)--(actions);
\draw[flow] (actions)--(mcpcalls);
\draw[flow] (tickets)--node[lbl,above]{investigation}(conv);
\draw[flow] (conv)--(msg);
\draw[flow] (msg.west) to[bend right=20] node[lbl,left]{chat\_message\_id} (actions.east);
\draw[flow] (users.south) to[bend right=25] (sreq.west);
\draw[flow] (sreq)--(satt);
\end{tikzpicture}
\caption{Relations principales. \tbl{audit\_logs}, \tbl{access\_requests},
\tbl{knowledge\_documents} et \tbl{disk\_metrics} ne sont pas représentées :
elles n'ont au plus qu'une clé étrangère vers \tbl{users}.}
\end{figure}

\clearpage

\section{Table par table}

% -------------------------------------------------- roles
\subsection{\tbl{roles} --- \file{models/role.py}}

\textbf{Rôle} : liste fermée des rôles applicatifs.\\
\textbf{Clé primaire} : \texttt{id} (Integer).\\
\textbf{Colonnes} : \texttt{id}, \texttt{name} (String(50), \texttt{unique},
\texttt{NOT NULL}).\\
\textbf{Relations} : \texttt{users} (un-à-plusieurs, \texttt{back\_populates}).\\
\textbf{Données} : trois lignes créées par \file{database/seed.py} ---
\texttt{linux\_engineer}, \texttt{storage\_engineer}, \texttt{administrator}.\\
\textbf{Utilisée dans} : \fn{require\_role}, \fn{request\_service.\_users\_with\_role},
\fn{admin\_service.list\_roles}, \fn{incident\_watcher.\_system\_user\_id}.

% -------------------------------------------------- users
\subsection{\tbl{users} --- \file{models/user.py}}

\textbf{Rôle} : comptes applicatifs.\\
\textbf{Clé primaire} : \texttt{id} (Integer, indexé).\\
\textbf{Clé étrangère} : \texttt{role\_id} \arrowto{} \texttt{roles.id}
(\texttt{NOT NULL}).\\
\textbf{Colonnes} : \texttt{email} (unique, indexé), \texttt{hashed\_password}
(bcrypt), \texttt{full\_name}, \texttt{is\_active} (défaut \texttt{True}),
\texttt{created\_at}.\\
\textbf{Utilisée dans} : \fn{authenticate\_user}, \fn{get\_current\_user},
\fn{\_extract\_user\_id} du middleware, et comme cible de \emph{toutes} les clés
étrangères d'auteur.

\begin{infobox}
Le mot de passe n'est jamais stocké en clair : \fn{hash\_password} utilise
\texttt{CryptContext(schemes=["bcrypt"])}. Le compte de service
\texttt{oasis-system@oasis.local} reçoit
\texttt{hash\_password(secrets.token\_urlsafe(32))} et
\texttt{is\_active=False}.
\end{infobox}

% -------------------------------------------------- audit_logs
\subsection{\tbl{audit\_logs} --- \file{models/audit\_log.py}}

\textbf{Rôle} : journal immuable des requêtes HTTP.\\
\textbf{Clé primaire} : \texttt{id} (Integer).\\
\textbf{Clé étrangère} : \texttt{user\_id} \arrowto{} \texttt{users.id}
(\emph{nullable} --- une requête anonyme est aussi journalisée).\\
\textbf{Colonnes} : \texttt{action} (méthode HTTP), \texttt{resource} (chemin),
\texttt{details} (JSON : \texttt{status\_code}, \texttt{duration\_ms},
\texttt{client\_ip}), \texttt{timestamp}.\\
\textbf{Écrite par} : \cls{AuditMiddleware} uniquement.\\
\textbf{Lue par} : \fn{audit\_service.list\_recent} \arrowto{}
\route{GET /admin/audit-logs}.

\begin{keybox}
Le docstring du modèle pose une règle : « cette table est pensée pour être
\emph{insert-only} --- on n'update et on ne supprime JAMAIS une ligne d'audit
log, on ne fait qu'en ajouter ». Aucun code du dépôt ne la modifie ni ne la
supprime : la règle est tenue.
\end{keybox}

% -------------------------------------------------- access_requests
\subsection{\tbl{access\_requests} --- \file{models/access\_request.py}}

\textbf{Rôle} : demande d'accès temporaire en lecture seule à un autre
workspace.\\
\textbf{Clé primaire} : \texttt{id} (Integer).\\
\textbf{Clés étrangères} : \texttt{requested\_by\_id} et \texttt{reviewed\_by\_id}
\arrowto{} \texttt{users.id} (deux relations distinctes, désambiguïsées par
\texttt{foreign\_keys=[...]}).\\
\textbf{Colonnes} : \texttt{target\_workspace} (texte libre), \texttt{reason},
\texttt{status} (\texttt{pending} / \texttt{approved} / \texttt{rejected} /
\texttt{expired}), \texttt{granted\_token}, \texttt{expires\_at},
\texttt{created\_at}, \texttt{reviewed\_at}.\\
\textbf{Utilisée dans} : \file{services/access\_request\_service.py} et
\route{/cross-domain}.

\begin{warnbox}
Le jeton généré à l'approbation est stocké dans \texttt{granted\_token}, mais
\textbf{aucun code du dépôt ne le vérifie ensuite}. Le statut \texttt{EXPIRED} de
l'énumération n'est jamais affecté non plus. Ce workflow est complet côté
enregistrement, incomplet côté application des droits.
\end{warnbox}

% -------------------------------------------------- incident_tickets
\subsection{\tbl{incident\_tickets} --- \file{models/incident\_ticket.py}}

\textbf{Rôle} : le ticket d'incident, pivot du cycle de vie.\\
\textbf{Clé primaire} : \texttt{id} (UUID).\\
\textbf{Clé étrangère} : \texttt{created\_by} \arrowto{} \texttt{users.id}
(\texttt{NOT NULL}).

\begin{longtable}{@{}p{3.6cm}p{2.6cm}p{8.4cm}@{}}
\toprule
\textbf{Colonne} & \textbf{Type} & \textbf{Contenu}\\
\midrule
\endhead
\texttt{ticket\_number} & String unique & \texttt{INC-2026-0001}, calculé par \fn{\_generate\_ticket\_number}\\
\texttt{title} & String & Ex. « Backup job backup-app-critical.timer failed »\\
\texttt{description} & Text & Phrase générée par la règle de détection\\
\texttt{severity} & Enum & \texttt{critical} / \texttt{high} / \texttt{medium} / \texttt{low}\\
\texttt{status} & Enum & \texttt{open} / \texttt{in\_progress} / \texttt{resolved} / \texttt{closed}\\
\texttt{affected\_system} & String & \textbf{Libellé lisible} (Description systemd)\\
\texttt{root\_cause} & Text & Rempli par la synthèse post-investigation\\
\texttt{impact\_summary} & Text & Idem\\
\texttt{mcp\_data} & \textbf{JSONB} & Relevé brut de l'outil au moment de la détection --- porte \texttt{job\_id} ou \texttt{volume\_id}, c'est-à-dire l'\textbf{identifiant technique}\\
\texttt{recommendations} & \textbf{JSONB} & Liste d'objets \texttt{\{order, description, action\_type, target, requires\_confirmation\}}\\
\texttt{correlation\_id} & UUID & Relie le ticket aux appels MCP de la même analyse\\
\texttt{created\_at} / \texttt{updated\_at} / \texttt{resolved\_at} & Timestamptz & Horodatages\\
\bottomrule
\end{longtable}

\textbf{Relations} : \texttt{creator} (\cls{User}), \texttt{notes} (cascade
\texttt{all, delete-orphan}), \texttt{action\_requests}.

\begin{keybox}
\textbf{La distinction à retenir absolument.} \texttt{affected\_system} est un
\emph{libellé} (« Sauvegarde de l'application critique »),
\texttt{mcp\_data["job\_id"]} est l'\emph{identifiant technique}
(\texttt{backup-app-critical.timer}). C'est le second que les outils attendent.
Toute la logique de \fn{IncidentContext.technical\_id} et de correction de cible
dans \fn{\_propose\_action} existe pour ne pas les confondre.
\end{keybox}

% -------------------------------------------------- notes
\subsection{\tbl{incident\_ticket\_notes}}

\textbf{Rôle} : piste d'audit lisible du ticket.\\
\textbf{Clé primaire} : \texttt{id} (UUID).\\
\textbf{Clés étrangères} : \texttt{ticket\_id} \arrowto{}
\texttt{incident\_tickets.id} ; \texttt{author\_id} \arrowto{} \texttt{users.id}.\\
\textbf{Écrite à} : détection automatique, ouverture d'investigation,
confirmation, rejet, échec d'action, clôture.\\
\textbf{Pourquoi} : « c'est ce qui permet de relire après coup pourquoi un
incident a été clos, et par qui --- l'horodatage seul des colonnes du ticket ne
le dit pas ».

% -------------------------------------------------- action_requests
\subsection{\tbl{action\_requests} --- \file{models/action\_request.py}}

\textbf{Rôle} : \textbf{la table qui matérialise la confirmation humaine.}\\
\textbf{Clé primaire} : \texttt{id} (UUID).

\begin{longtable}{@{}p{3.8cm}p{2.9cm}p{7.9cm}@{}}
\toprule
\textbf{Colonne} & \textbf{Type} & \textbf{Contenu}\\
\midrule
\endhead
\texttt{action\_type} & String indexé & Nom exact de l'outil MCP\\
\texttt{target} & String & Cible technique\\
\texttt{parameters} & \textbf{JSONB} & Arguments, \emph{sans} \texttt{confirm}\\
\texttt{status} & Enum & \texttt{pending} / \texttt{confirmed} / \texttt{completed} / \texttt{failed} / \texttt{rejected}\\
\texttt{requested\_by} & FK users & Auteur de la proposition\\
\texttt{confirmed\_by} & FK users, nullable & Qui a cliqué « Confirm »\\
\texttt{incident\_ticket\_id} & FK, nullable & Ticket rattaché\\
\texttt{chat\_message\_id} & FK BigInteger, \texttt{ON DELETE SET NULL} & Message qui porte la proposition\\
\texttt{result} & \textbf{JSONB} & Réponse brute de l'outil, ou \texttt{\{"error": ...\}}\\
\texttt{created\_at} / \texttt{confirmed\_at} / \texttt{completed\_at} & Timestamptz & Les trois temps de l'action\\
\bottomrule
\end{longtable}

\textbf{Pourquoi \texttt{chat\_message\_id}} : le commentaire du modèle est
explicite --- « Sert à réafficher les boutons Confirmer/Rejeter quand on rouvre
le fil : sans lui, une proposition disparaîtrait de l'écran au rechargement
alors qu'elle reste en attente en base. »

% -------------------------------------------------- mcp_calls
\subsection{\tbl{mcp\_calls} --- \file{models/mcp\_call.py}}

\textbf{Rôle} : preuve qu'une réponse de l'agent s'appuie sur un relevé réel.\\
\textbf{Clé primaire} : \texttt{id} (UUID).\\
\textbf{Clés étrangères} : \texttt{user\_id} \arrowto{} \texttt{users.id}
(\texttt{NOT NULL}, indexé) ; \texttt{action\_request\_id} \arrowto{}
\texttt{action\_requests.id} (nullable).\\
\textbf{Colonnes} : \texttt{persona} (« storage »), \texttt{tool\_name},
\texttt{arguments} (JSONB), \texttt{result} (JSONB), \texttt{status}
(\texttt{success} / \texttt{failed} / \texttt{denied}), \texttt{error\_message},
\texttt{correlation\_id} (indexé), \texttt{created\_at}.\\
\textbf{Écrite par} : \fn{mcp\_audit.log\_mcp\_call}, appelée \emph{uniquement}
depuis \fn{StorageMCPClient.call\_tool}.

\begin{infobox}
Le commentaire de \texttt{arguments} précise : ce sont « les arguments réellement
envoyés au tool (après normalisation et garde-fou, donc reflète le
\texttt{confirm=True/False} réel, pas ce que le LLM avait proposé) ». C'est
exactement ce qu'il faut pour un audit de sécurité.
\end{infobox}

% -------------------------------------------------- chat
\subsection{\tbl{chat\_conversations} et \tbl{chat\_messages}}

\textbf{\tbl{chat\_conversations}} --- clé primaire \texttt{id} (BigInteger) ;
\texttt{user\_id} \arrowto{} \texttt{users.id} ; \texttt{title}
(\(\le\) 120 caractères, jamais vide) ; \texttt{incident\_ticket\_id}
(\texttt{ON DELETE SET NULL}) ; \texttt{created\_at} ; \texttt{updated\_at}.

\textbf{\tbl{chat\_messages}} --- clé primaire \texttt{id} (BigInteger) ;
\texttt{conversation\_id} \arrowto{} \texttt{chat\_conversations.id}
(\texttt{ON DELETE CASCADE}) ; \texttt{role} (\texttt{"user"} ou
\texttt{"assistant"}, stocké en texte) ; \texttt{content} ; \texttt{created\_at}.

\begin{keybox}
\textbf{Index unique partiel.} La migration \file{b6d0f3a91c47} crée
\texttt{uq\_chat\_conversations\_incident\_ticket} sur
\texttt{(incident\_ticket\_id, user\_id)}, unique, avec
\texttt{postgresql\_where = incident\_ticket\_id IS NOT NULL}. Justification :
« Un ingénieur n'a qu'un fil d'investigation par ticket : l'unicité est ce qui
permet de rouvrir l'existant sans arbitrer entre plusieurs candidats \ldots{}
L'index est partiel, la colonne étant nulle pour les fils ordinaires, qui sont
la majorité. »
\end{keybox}

\noindent
\textbf{Pourquoi le propriétaire n'est pas dupliqué sur le message} : « Le
propriétaire se lit sur la conversation : le dupliquer ici ouvrirait la porte à
un message rattaché à un fil qui n'est pas le sien. »

% -------------------------------------------------- service_requests
\subsection{\tbl{service\_requests} et ses satellites}

\textbf{\tbl{service\_requests}} --- clé primaire UUID.
Clés étrangères : \texttt{created\_by} et \texttt{assigned\_to} \arrowto{}
\texttt{users.id} (les deux \texttt{NOT NULL} et indexées).
Colonnes notables : \texttt{direction} (Enum), \texttt{title},
\texttt{description} (nullable), \texttt{request\_type}, \texttt{priority},
\texttt{status}, \texttt{deadline}, \texttt{administrator\_response},
\texttt{decided\_at}, \texttt{acknowledged\_at}, \texttt{completed\_at}.

\begin{longtable}{@{}p{4.0cm}p{10.8cm}@{}}
\toprule
\textbf{Table satellite} & \textbf{Statut réel}\\
\midrule
\endhead
\tbl{request\_events} & \textbf{Active}. Une ligne par transition (\texttt{from\_status}, \texttt{to\_status}, \texttt{note}, \texttt{actor\_id}). Jamais renvoyée au client : « l'interface se contente du statut courant et de son horodatage ».\\
\tbl{request\_attachments} & \textbf{Active}. Fichier écrit sous l'UUID de la ligne ; le nom d'origine ne survit que comme métadonnée, « ce qui rend toute traversée de répertoire impossible ».\\
\tbl{request\_messages} & \textbf{Gelée}. Le docstring l'assume : « le fil a été retiré de l'interface et de l'API. La table est conservée pour ne pas détruire les échanges déjà enregistrés, mais plus rien n'y écrit. »\\
\bottomrule
\end{longtable}

\begin{infobox}
\textbf{Trois modèles au nom proche, à ne pas confondre --- le code le dit
lui-même.} \cls{AccessRequest} = accès temporaire cross-domain ;
\cls{ActionRequest} = exécution d'un outil MCP sensible ;
\cls{ServiceRequest} = canal de communication humain, qui « ne déclenche aucune
action technique ».
\end{infobox}

% -------------------------------------------------- knowledge / metrics
\subsection{\tbl{knowledge\_documents} --- \file{models/knowledge\_document.py}}

\textbf{Rôle} : métadonnées des documents du Data Hub. Le fichier lui-même vit
sur disque ; « cette table ne garde que les métadonnées et l'état d'indexation
Chroma ».\\
\textbf{Clé primaire} : \texttt{id} (UUID) --- c'est aussi le nom du fichier sur
disque.\\
\textbf{Clé étrangère} : \texttt{uploaded\_by} \arrowto{} \texttt{users.id}.\\
\textbf{Colonnes} : \texttt{filename}, \texttt{stored\_path},
\texttt{content\_type}, \texttt{size\_bytes}, \texttt{collection} (défaut
\texttt{storage\_kb}, indexé), \texttt{chunk\_count}, \texttt{status}
(\texttt{pending} / \texttt{indexed} / \texttt{failed}), \texttt{error},
\texttt{created\_at}, \texttt{indexed\_at}.

\subsection{\tbl{disk\_metrics} --- \file{models/metric.py}}

\textbf{Rôle} : un échantillon de performance, pour un device, à un instant.\\
\textbf{Clé primaire} : \texttt{id} (UUID).\\
\textbf{Aucune clé étrangère.}\\
\textbf{Colonnes} : \texttt{device} (indexé), \texttt{iops},
\texttt{throughput\_mbps}, \texttt{latency\_ms}, \texttt{recorded\_at} (indexé).\\
\textbf{Écrite par} : \fn{metrics\_service.collect\_once}.
\textbf{Lue par} : \fn{read\_series}. \textbf{Purgée par} : \fn{purge\_old}
(rétention 31 jours).\\
\textbf{Principe} : « l'historique 24h/7j/30j se construit au fil des
prélèvements, il n'est jamais fabriqué ».

\section{Les migrations Alembic}

Quatorze migrations forment une \textbf{chaîne linéaire}, sans branche :

\begin{lstlisting}[caption={Ordre réel des révisions}]
a5bd2724f8a3  create users and roles tables            (base, down_revision = None)
a7678a14436e  add audit_logs table
8aab52e3789f  add access_requests table
cfa440d22089  add incidents table            <-- MIGRATION VIDE (upgrade: pass)
c93419b63533  add incident_tickets and action_requests tables
d41f7c2b9e10  add knowledge_documents table (Data Hub)
e58a3d9c4712  add mcp_calls table
f7c1b0d5a983  add disk_metrics table (real performance history)
b2f4d6c81a05  add chat_messages table (historique de conversation)
c7e1a94b3d20  add chat_conversations (plusieurs fils par utilisateur)
d9a3f17c6b42  add service_requests, request_messages, request_events
f4b28c50e173  add decided_at to service_requests
a1c7e94d2f80  add request_attachments
b6d0f3a91c47  link investigation chat to incidents            <-- HEAD
\end{lstlisting}

\begin{warnbox}
\textbf{Détail à ne pas cacher.} La révision \file{cfa440d22089}, malgré son
titre « add incidents table », a un \fn{upgrade()} qui ne contient qu'un
\texttt{pass} : c'est un squelette généré puis abandonné. Les vraies tables
d'incident arrivent avec \file{c93419b63533}.
\end{warnbox}

\section{Workflow base de données : le cycle complet d'une action}

Reprenons la demande de l'énoncé --- \emph{création \arrowto{} stockage
\arrowto{} récupération \arrowto{} modification du statut \arrowto{} affichage}
--- appliquée à une \cls{ActionRequest} réelle.

\begin{longtable}{@{}p{0.7cm}p{4.4cm}p{9.6cm}@{}}
\toprule
\textbf{\#} & \textbf{Étape} & \textbf{Effet exact en base}\\
\midrule
\endhead
1 & Le LLM appelle \fn{run\_backup} & Rien encore : l'appel est intercepté par \fn{\_propose\_action}\\
2 & \fn{create\_action\_request} & \texttt{INSERT} dans \tbl{action\_requests} : \texttt{status='pending'}, \texttt{requested\_by}, \texttt{target}, \texttt{parameters}, \texttt{incident\_ticket\_id}. \texttt{commit} + \texttt{refresh}\\
3 & \fn{chat\_service.save\_turn} & \texttt{INSERT} de deux \tbl{chat\_messages} dans la même transaction ; \texttt{chat\_conversations.updated\_at = now()}\\
4 & \fn{attach\_actions} & \texttt{UPDATE action\_requests SET chat\_message\_id = <id du message assistant> WHERE id IN (...)}\\
5 & Rechargement du fil & \texttt{SELECT} des messages avec leur relation \texttt{action\_requests} \arrowto{} \cls{ChatMessageOut} \arrowto{} boutons rendus par \cls{ActionProposal}\\
6 & Clic « Confirm » & \texttt{UPDATE} : \texttt{status='confirmed'}, \texttt{confirmed\_by}, \texttt{confirmed\_at}. \textbf{commit immédiat}\\
7 & Appel MCP & \texttt{INSERT} dans \tbl{mcp\_calls} : \texttt{tool\_name}, \texttt{arguments} (avec \texttt{confirm=True}), \texttt{result}, \texttt{status}, \texttt{action\_request\_id}\\
8 & Qualification & \texttt{UPDATE} : \texttt{status='completed'} ou \texttt{'failed'}, \texttt{result} (JSONB), \texttt{completed\_at}\\
9 & \fn{\_close\_the\_loop} & \texttt{INSERT} dans \tbl{incident\_ticket\_notes} avec le verdict\\
10 & Si résolu & \texttt{UPDATE incident\_tickets} : \texttt{status='resolved'}, \texttt{resolved\_at} ; plus une note de clôture\\
11 & Sinon, plus tard & \fn{incident\_watcher.\_verify\_resolved} refait le test au passage suivant (300\,s) et clôt si l'indicateur est redevenu sain\\
\bottomrule
\end{longtable}

\section{Workflow base de données : naissance d'un incident}

\begin{enumerate}
  \item \fn{list\_backups()} lit les timers systemd par SSH.
  \item \fn{check\_backup\_status(job)} renvoie un dictionnaire de ticket, ou
  \texttt{None}.
  \item \fn{\_latest\_ticket} + \fn{\_already\_handled} : \texttt{SELECT} du
  dernier ticket sur ce système, pour ne pas rouvrir un incident déjà traité.
  \item \fn{create\_ticket\_from\_diagnosis} \arrowto{} \fn{\_open\_ticket\_for}
  fait un \texttt{SELECT} sur \texttt{mcp\_data['job\_id']} (opérateur JSONB
  \texttt{->>}) parmi les statuts \texttt{open} et \texttt{in\_progress}.
  \begin{itemize}
    \item \textbf{Trouvé} \arrowto{} \texttt{UPDATE} du titre, de la description,
    de la sévérité, de \texttt{mcp\_data} et des recommandations. Un seul ticket
    pour une anomalie qui persiste.
    \item \textbf{Absent} \arrowto{} \fn{\_generate\_ticket\_number}
    (\texttt{SELECT max(substr(ticket\_number, 10))}) puis \texttt{INSERT}.
  \end{itemize}
  \item Si le ticket est neuf : \texttt{INSERT} d'une première note d'audit.
\end{enumerate}

% =====================================================================
\chapter{Les workflows métier}
\label{ch:workflows}

Chaque workflow est présenté selon la chaîne demandée, en n'employant que les
composants réellement présents dans le code. Quand une étape n'existe pas (pas
de LLM, pas de MCP), elle est simplement absente plutôt que mentionnée pour la
forme.

% =====================================================================
\section{Workflow 1 --- Authentification}

\begin{longtable}{@{}p{2.9cm}p{11.9cm}@{}}
\toprule
\textbf{Étape} & \textbf{Ce qui se passe}\\
\midrule
\endhead
Utilisateur & Saisit e-mail, mot de passe, coche ou non « se souvenir »\\
Frontend & \cls{LoginPage} \arrowto{} \fn{signIn(email, password, remember)} de \cls{AuthContext}\\
API & \route{POST /api/v1/auth/login} --- \textbf{aucune dépendance d'authentification}\\
Backend & \fn{auth\_service.login} \arrowto{} \fn{authenticate\_user} : \texttt{SELECT} sur \tbl{users}, \fn{verify\_password} (bcrypt), contrôle \texttt{is\_active}\\
Résultat & \fn{create\_access\_token(subject=email, role=role.name)} \arrowto{} JWT HS256 avec \texttt{sub}, \texttt{role}, \texttt{exp}\\
Frontend & \fn{tokenStore.set(token, remember)} : \texttt{localStorage} si « se souvenir », sinon \texttt{sessionStorage}. Puis \route{GET /auth/me} pour charger le profil\\
Ensuite & \cls{RequireAuth} laisse passer ; \cls{AppLayout} s'affiche\\
\bottomrule
\end{longtable}

\subsection*{Reprise de session au rechargement}
\cls{AuthProvider} n'accorde pas sa confiance au seul jeton stocké : il appelle
\route{GET /auth/me}. Commentaire du code : « un jeton en stockage ne suffit
pas, il faut que le backend le reconnaisse encore ». En cas d'échec,
\fn{tokenStore.clear()}.

\subsection*{Expiration en cours d'usage}
Toute réponse 401 déclenche \fn{onUnauthorized()} dans \file{lib/api.ts}, câblé
sur \fn{signOut}. Le hook \fn{useResource} \emph{n'affiche pas} l'erreur dans ce
cas : « Un 401 est déjà traité par le handler global : ne pas l'afficher. »

% =====================================================================
\section{Workflow 2 --- Dashboard}

\begin{longtable}{@{}p{2.9cm}p{11.9cm}@{}}
\toprule
\textbf{Étape} & \textbf{Ce qui se passe}\\
\midrule
\endhead
Utilisateur & Arrive sur \route{/storage/dashboard}\\
Frontend & Quatre appels : \route{/overview}, \route{/performance?range=24h}, \route{/activity?limit=6}, \route{/insights}\\
API & \file{api/v1/dashboard.py}, protégé par \fn{require\_role("storage\_engineer","administrator")}. Routes en \texttt{def} synchrone, donc exécutées dans le threadpool (SSH bloquant)\\
Backend & \fn{kpi\_service.get\_overview} : \fn{\_storage\_health()} sert le cache (\(\le 60\)\,s), sinon rafraîchit\\
Tool & \fn{get\_storage\_health()} --- \textbf{appel direct, pas via MCP} --- lance \fn{get\_capacity} et \fn{list\_backups} en parallèle dans un \cls{ThreadPoolExecutor}\\
Infrastructure & Deux sessions SSH : \texttt{df -T -B1 -P} et \texttt{systemctl list-timers} + \texttt{systemctl show}\\
Résultat & Capacité, jobs, alertes déterministes (seuils 90/95\,\% et 24\,h), \texttt{overall\_status}\\
Backend (suite) & Agrégats SQL sur \tbl{incident\_tickets}, calcul du \texttt{health\_score} et du \texttt{headline}\\
Frontend & \cls{KpiCard}, \cls{CapacityOverview}, \cls{PerformanceChart}, \cls{RecentActivity}, \cls{OasisInsights}, \cls{AttentionBanner}\\
\bottomrule
\end{longtable}

\begin{keybox}
\textbf{Pourquoi ce workflow ne passe pas par MCP ni par le LLM.} Le dashboard
doit s'afficher vite et rendre toujours la même valeur pour le même état.
Passer par le sous-processus MCP ajouterait une latence de démarrage, et passer
par le LLM rendrait l'affichage non déterministe. Le fichier le déclare :
« Aucun appel LLM ici. »
\end{keybox}

% =====================================================================
\section{Workflow 3 --- Chat avec l'assistant}

\begin{figure}[h!]
\centering
\begin{tikzpicture}[node distance=6mm]
\node[box,minimum width=100mm] (u) {\textbf{Utilisateur} --- « Quels volumes approchent de leur limite ? »};
\node[bluebox,below=of u,minimum width=100mm] (f) {\textbf{ChatPage} \quad api.post('/storage/chat', \{message, conversation\_id\})};
\node[bluebox,below=of f,minimum width=100mm] (a) {\textbf{API} \quad require\_role \arrowto{} history\_for\_llm(10 msg) \arrowto{} contexte incident éventuel};
\node[vibox,below=of a,minimum width=100mm] (ag) {\textbf{StoragePersona.handle\_message} \quad RAG Chroma + prompt système + historique};
\node[vibox,below=of ag,minimum width=100mm] (g1) {\textbf{Groq} \quad chat\_with\_tools(messages, tools) \arrowto{} tool\_calls = [get\_capacity]};
\node[ambox,below=of g1,minimum width=100mm] (m) {\textbf{MCP} \quad client stdio \arrowto{} MCPServer \arrowto{} get\_capacity()};
\node[redbox,below=of m,minimum width=100mm] (v) {\textbf{VM} \quad ssh: df -T -B1 -P \arrowto{} \{"volumes": [\ldots]\}};
\node[greenbox,below=of v,minimum width=100mm] (au) {\textbf{Audit} \quad INSERT mcp\_calls (tool\_name, arguments, result, correlation\_id)};
\node[vibox,below=of au,minimum width=100mm] (g2) {\textbf{Groq} \quad second tour : rédige la réponse à partir du résultat};
\node[bluebox,below=of g2,minimum width=100mm] (s) {\textbf{save\_turn} \quad INSERT 2 chat\_messages + updated\_at};
\node[box,below=of s,minimum width=100mm,draw=oasisdark] (r) {\textbf{ChatTurnOut} \arrowto{} affichage React};
\foreach \x/\y in {u/f,f/a,a/ag,ag/g1,g1/m,m/v,v/au,au/g2,g2/s,s/r} \draw[flowb] (\x)--(\y);
\end{tikzpicture}
\caption{Un tour de chat avec un outil de lecture.}
\end{figure}

\clearpage

\subsection*{Variante : le modèle veut relancer une sauvegarde}
À l'étape MCP, le nom appartient à \texttt{REMEDIATION\_TOOLS}. L'agent
n'exécute rien : il crée une \tbl{action\_requests} \texttt{pending} et renvoie
au modèle un résultat qui \emph{lui dit} de présenter la chose comme une
recommandation. Le tour se termine par un message accompagné d'un composant
\cls{ActionProposal}.

\subsection*{Effet de bord : création automatique de ticket}
Si \fn{evaluate\_tool\_result} détecte une anomalie (volume \(\ge\) 85\,\%, job
en échec), un \tbl{incident\_tickets} est créé à la fin du tour, attribué à
l'utilisateur courant, avec le \texttt{correlation\_id} du tour --- ce qui relie
le ticket aux \tbl{mcp\_calls} qui l'ont motivé.

% =====================================================================
\section{Workflow 4 --- Détection automatique d'un incident}

\textbf{Aucun utilisateur n'est impliqué.} C'est une tâche de fond démarrée par
le \fn{lifespan} de \file{main.py}.

\begin{longtable}{@{}p{2.9cm}p{11.9cm}@{}}
\toprule
\textbf{Étape} & \textbf{Ce qui se passe}\\
\midrule
\endhead
Déclencheur & \fn{incident\_watcher.run\_forever}, boucle toutes les 300\,s\\
Isolation & \fn{asyncio.to\_thread(scan\_once)} --- la lecture SSH est bloquante et ne doit pas figer la boucle d'événements\\
Tool & \fn{list\_backups()} en appel direct\\
Infrastructure & \texttt{systemctl list-timers -{}-all}, puis \texttt{systemctl show} sur toutes les units \textbf{en une seule commande} (séparateur \texttt{\#\#\#OASIS-UNIT=})\\
Règle & \fn{check\_backup\_status(job)} : seulement si \texttt{status == "failed"}\\
Vérification d'abord & \fn{\_verify\_resolved} referme les tickets dont une action \texttt{completed} existe et dont le job est redevenu sain\\
Anti-doublon & \fn{\_already\_handled} : on ne rouvre qu'à partir d'une exécution postérieure au ticket précédent\\
Backend & \fn{create\_ticket\_from\_diagnosis} avec \texttt{created\_by = compte de service}\\
Audit & Première note : « Detected automatically by monitoring: job \ldots{} came back failed (systemd result: \ldots) »\\
Frontend & Le ticket apparaît au prochain chargement de \route{/storage/incidents}\\
\bottomrule
\end{longtable}

\begin{infobox}
\textbf{Pourquoi cette tâche existe.} Le docstring le justifie : « les tickets ne
sont créés que par le persona Storage, pendant une conversation. Un job qui
échoue la nuit ne produit donc un ticket que si quelqu'un pense à interroger
l'agent le lendemain. Les alertes du dashboard, elles, sont recalculées à chaque
appel et ne laissent aucune trace. »
\end{infobox}

% =====================================================================
\section{Workflow 5 --- Investigation d'un incident}

\begin{longtable}{@{}p{2.9cm}p{11.9cm}@{}}
\toprule
\textbf{Étape} & \textbf{Ce qui se passe}\\
\midrule
\endhead
Utilisateur & Ouvre un ticket, clique « Investigate with OASIS »\\
Frontend & \route{POST /incidents/\{id\}/investigate}\\
API & \fn{require\_role("storage\_engineer","administrator")}. Route \texttt{async}, volontairement bloquante : « l'ingénieur vient de cliquer et attend le diagnostic, un accusé de réception vide ne lui apprendrait rien »\\
Backend & \fn{incident\_investigation.start} : fil existant \arrowto{} retour immédiat, \texttt{created=False}\\
Agent & \fn{handle\_message} avec \cls{IncidentContext} : le prompt contient le numéro du ticket, la cible technique, le relevé MCP et trois consignes\\
MCP & Le modèle appelle typiquement \fn{get\_backup\_logs} puis \fn{get\_backup\_job\_status}\\
Infrastructure & \texttt{journalctl -u <service> -n 20 -o short-iso}, \texttt{systemctl show}\\
Proposition & Le modèle appelle un outil de remédiation \arrowto{} \tbl{action\_requests} \texttt{pending}\\
Filet & Si aucune proposition n'a été enregistrée : \fn{\_recover\_written\_call} puis \fn{\_fallback\_action}\\
Écriture & Conversation créée, \texttt{incident\_ticket\_id} posé, \fn{save\_turn}, \fn{attach\_actions}, \fn{mark\_in\_progress}\\
Synthèse & Second appel Groq \arrowto{} JSON \texttt{\{root\_cause, impact\_summary\}} recopié sur le ticket\\
Audit & Note « Investigation opened by \ldots{} proposed N action(s) for confirmation »\\
Frontend & \cls{InvestigationOut} \arrowto{} ouverture du fil, ticket rechargé avec sa fiche complète\\
\bottomrule
\end{longtable}

\begin{keybox}
\textbf{Trois filets successifs pour ne pas perdre la proposition.}
\emph{(1)} l'appel d'outil normal ;
\emph{(2)} \fn{\_recover\_written\_call}, qui récupère un appel écrit en JSON
dans la prose du modèle ;
\emph{(3)} \fn{\_fallback\_action}, qui reprend l'\texttt{action\_type} calculé
par la règle de détection au moment de l'ouverture du ticket.
Le cycle ne doit pas « s'arrêter à l'avant-dernière marche ».
\end{keybox}

% =====================================================================
\section{Workflow 6 --- Confirmation humaine et clôture}

C'est le workflow central du projet. Il répond exactement à la chaîne
\emph{demande \arrowto{} analyse \arrowto{} proposition \arrowto{} confirmation
\arrowto{} exécution \arrowto{} résultat}.

\begin{figure}[h!]
\centering
\begin{tikzpicture}[node distance=5.5mm]
\node[box,minimum width=108mm] (n1) {\textbf{1. Analyse} --- l'agent lit la VM et identifie la cause racine};
\node[ambox,below=of n1,minimum width=108mm] (n2) {\textbf{2. Proposition} --- \_propose\_action() \arrowto{} action\_requests(status=pending), \textbf{rien n'est exécuté}};
\node[bluebox,below=of n2,minimum width=108mm] (n3) {\textbf{3. Affichage} --- ActionProposal : libellé lisible + action\_type(target) + avertissement si destructif};
\node[greenbox,below=of n3,minimum width=108mm] (n4) {\textbf{4. Décision humaine} --- clic « Confirm » ou « Reject »};
\node[ambox,below=of n4,minimum width=108mm] (n5) {\textbf{5. Démarrage MCP} --- StorageMCPClient().start() ; échec \arrowto{} HTTP 503, action toujours PENDING};
\node[redbox,below=of n5,minimum width=108mm] (n6) {\textbf{6. Exécution} --- call\_tool(action\_type, confirm=True, **parameters) \arrowto{} SSH \arrowto{} VM};
\node[greenbox,below=of n6,minimum width=108mm] (n7) {\textbf{7. Qualification} --- result contient "error" ? FAILED : COMPLETED ; audit dans mcp\_calls};
\node[bluebox,below=of n7,minimum width=108mm] (n8) {\textbf{8. Vérification} --- incident\_verification.verify() relit l'indicateur sur la VM};
\node[box,below=of n8,minimum width=108mm,draw=oasisdark] (n9) {\textbf{9. Clôture} --- si verdict.resolved : resolve\_ticket + note ; sinon le ticket reste ouvert avec la raison};
\foreach \x/\y in {n1/n2,n2/n3,n3/n4,n4/n5,n5/n6,n6/n7,n7/n8,n8/n9} \draw[flowb] (\x)--(\y);
\end{tikzpicture}
\caption{Les neuf temps de la confirmation humaine, tels qu'ils sont codés.}
\end{figure}

\clearpage

\subsection*{Ce que voit l'ingénieur avant de cliquer}
\cls{ActionProposal} traduit l'appel d'outil en phrase :

\begin{lstlisting}[caption={frontend/src/components/ActionProposal.tsx}]
const ACTION_LABELS = {
  run_backup:               (t) => `Re-run backup job ${t} now`,
  create_snapshot:          (t) => `Take a snapshot of volume ${t}`,
  restore_from_backup:      (t) => `Restore ${t} from a backup`,
  repair_backup_repository: (t) => `Repair the backup repository configured for ${t}`,
};
const DESTRUCTIVE_ACTIONS = new Set(['restore_from_backup']);
\end{lstlisting}

\noindent
Sous la phrase, l'appel technique brut est affiché en police à chasse fixe :
\texttt{action\_type(target)}. Pour une action destructive, un avertissement
rouge s'ajoute : « This action overwrites existing data. Check the target before
confirming. »

\subsection*{La branche « Reject »}
\route{POST /storage/actions/\{id\}/reject} passe l'action en \texttt{REJECTED}
et, si elle porte un ticket, écrit une note nommée avec le motif éventuel :
« Un rejet est une décision, pas un non-événement : il doit rester lisible dans
l'historique du ticket. »

% =====================================================================
\section{Workflow 7 --- Data Hub et RAG}

\begin{longtable}{@{}p{2.9cm}p{11.9cm}@{}}
\toprule
\textbf{Étape} & \textbf{Ce qui se passe}\\
\midrule
\endhead
Utilisateur & Dépose un fichier sur \route{/storage/data-hub}\\
Frontend & \fn{api.upload('/storage/documents', file)} --- \texttt{FormData}\\
API & \route{POST /storage/documents}, \fn{require\_role}, \texttt{await file.read()}\\
Backend & \fn{knowledge\_service.create\_document} : taille \(\le\) 10\,Mo, extension autorisée, écriture sous \texttt{<uuid><ext>}\\
Extraction & Texte brut, ou \texttt{pypdf} pour un PDF ; refus si \(<\) 20 mots utiles\\
Découpage & \fn{chunk\_text} : 500 mots, recouvrement 50\\
Vectorisation & \fn{embed\_passages} avec le préfixe \texttt{"passage: "}\\
Chroma & \texttt{collection.upsert} dans \texttt{storage\_kb}, identifiants dérivés de l'UUID\\
Résultat & \texttt{status=indexed} et \texttt{chunk\_count}, ou \texttt{status=failed} avec le message d'erreur conservé\\
Frontend & La ligne apparaît avec son statut ; un document en échec explique pourquoi\\
Plus tard & Chaque question du chat déclenche \fn{embed\_query} (préfixe \texttt{"query: "}) et les 3 passages les plus proches entrent dans le prompt\\
\bottomrule
\end{longtable}

\begin{keybox}
\textbf{Pourquoi deux fonctions d'embedding et non une seule.} Les modèles E5
sont \emph{asymétriques} : entraînés avec \texttt{"query: "} devant la question
et \texttt{"passage: "} devant le document. Le fichier le dit : « une fonction
unique \fn{embed()} serait forcément fausse d'un côté ou de l'autre ».
\end{keybox}

% =====================================================================
\section{Workflow 8 --- Demandes de service}

\begin{longtable}{@{}p{2.9cm}p{11.9cm}@{}}
\toprule
\textbf{Étape} & \textbf{Ce qui se passe}\\
\midrule
\endhead
Utilisateur & Ingénieur Storage : « J'ai besoin de 500\,Go supplémentaires »\\
Frontend & \route{POST /requests} avec titre, description, type, priorité\\
API & \fn{require\_role("storage\_engineer","administrator")} --- le Linux Engineer n'a pas accès à ce canal\\
Backend & \fn{create\_request} : la \textbf{direction} est déduite du rôle, jamais du corps de la requête\\
Destinataire & S'il n'y a qu'un administrateur actif, il est choisi automatiquement ; s'il y en a plusieurs et qu'aucun n'est précisé, la demande est refusée\\
Base & \texttt{INSERT service\_requests} (\texttt{status=pending}) + \texttt{INSERT request\_events} (\texttt{to\_status=pending}, note « Demande créée »)\\
Administrateur & Voit la demande dans sa boîte « received » (\route{GET /requests?box=received})\\
Transition & \route{PATCH /requests/\{id\}/status} : contrôle \texttt{assigned\_to == user.id}, puis table \fn{ALLOWED\_TRANSITIONS}\\
Décision & \texttt{approved} ou \texttt{rejected} \arrowto{} \texttt{decided\_at} figé, \texttt{administrator\_response} enregistrée\\
Audit & Une ligne \tbl{request\_events} par transition, jamais renvoyée au client\\
Frontend & Le statut et la date de décision s'affichent dans la liste et le détail\\
\bottomrule
\end{longtable}

\noindent
\textbf{Aucun LLM, aucun MCP, aucune infrastructure} n'intervient : c'est un
canal humain.

% =====================================================================
\section{Workflow 9 --- Recherche globale}

Utilisateur \arrowto{} \cls{GlobalSearch} (Topbar) \arrowto{}
\route{GET /storage/search?q=} \arrowto{} \fn{search\_service.search}
\arrowto{} quatre requêtes \texttt{ILIKE} échappées (incidents, demandes,
conversations, documents) \arrowto{} liste de \cls{SearchHitOut}
(\texttt{type}, \texttt{id}, \texttt{title}, \texttt{subtitle},
\texttt{timestamp}) \arrowto{} le frontend traduit le \texttt{type} en route :

\begin{lstlisting}[caption={frontend/src/components/layout/GlobalSearch.tsx}]
const DESTINATION = {
  incident:     (id) => `/storage/incidents?incident=${id}`,
  request:      (id) => `/storage/requests?request=${id}`,
  conversation: (id) => `/storage/chat?conversation=${id}`,
  document:     (id) => `/storage/data-hub?document=${id}`,
};
\end{lstlisting}

% =====================================================================
\section{Workflow 10 --- Rapports PDF}

Utilisateur \arrowto{} \cls{ReportsPage} \arrowto{} \fn{api.download} \arrowto{}
\route{GET /storage/reports/global} ou
\route{GET /storage/reports/incidents/\{id\}} \arrowto{}
\fn{report\_service.build\_*\_report} (lecture SQL + lecture de capacité par SSH)
\arrowto{} \texttt{bytes} PDF \arrowto{} \texttt{Response} avec
\texttt{Content-Disposition} et \texttt{Access-Control-Expose-Headers}
\arrowto{} le frontend lit le blob, extrait le nom de fichier et déclenche le
téléchargement.

% =====================================================================
\section{Workflow 11 --- Accès cross-domain}

\begin{enumerate}
  \item \route{POST /cross-domain/requests} --- n'importe quel utilisateur
  connecté crée une demande (\texttt{target\_workspace}, \texttt{reason}).
  \item \route{GET /cross-domain/requests/pending} --- réservé à
  \texttt{administrator}.
  \item \route{POST /cross-domain/requests/\{id\}/decision} --- l'admin approuve
  ou refuse. En cas d'approbation, \fn{create\_cross\_domain\_token} génère un
  JWT de 30 minutes avec \texttt{scope="cross\_domain\_read\_only"} et le
  workspace, stocké dans \texttt{granted\_token}.
\end{enumerate}

\begin{warnbox}
\textbf{Le workflow s'arrête là.} Aucun code ne vérifie ensuite ce scope, et
aucune interface frontend ne l'expose. C'est une brique préparée et testable,
pas une fonctionnalité utilisable de bout en bout aujourd'hui.
\end{warnbox}

% =====================================================================
\section{Workflow 12 --- Collecte de métriques}

\begin{longtable}{@{}p{2.9cm}p{11.9cm}@{}}
\toprule
\textbf{Étape} & \textbf{Ce qui se passe}\\
\midrule
\endhead
Déclencheur & \fn{metrics\_collector.run\_forever}, toutes les 300\,s\\
Isolation & \fn{asyncio.to\_thread(\_collect)}\\
Tool & \fn{sample\_disk\_performance()} --- \textbf{appel direct} ; cet outil n'est \emph{pas} dans le \fn{TOOL\_REGISTRY} et n'est donc pas exposé au LLM\\
Infrastructure & Une session SSH : deux lectures de \texttt{/proc/diskstats} espacées de 2\,s\\
Calcul & \fn{\_rates} : IOPS, débit (secteurs \(\times\) 512), latence moyenne ; \texttt{None} si un compteur a reculé\\
Base & Une ligne \tbl{disk\_metrics} par device\\
Purge & Tous les 288 cycles (\(\approx\) une fois par jour), suppression au-delà de 31 jours\\
Lecture & \route{GET /storage/dashboard/performance?range=} \arrowto{} \fn{read\_series} (\texttt{date\_bin}, agrégation en deux passes)\\
Frontend & \cls{PerformanceChart} (recharts) ; si \texttt{collecting} est vrai, un message remplace la courbe\\
\bottomrule
\end{longtable}

% =====================================================================
\section{Synthèse : quel workflow mobilise quoi ?}

\begin{longtable}{@{}p{3.6cm}ccccc@{}}
\toprule
\textbf{Workflow} & \textbf{LLM} & \textbf{RAG} & \textbf{MCP} & \textbf{SSH} & \textbf{Confirmation}\\
\midrule
\endhead
Authentification & --- & --- & --- & --- & ---\\
Dashboard & --- & --- & --- & \(\bullet\) & ---\\
Chat & \(\bullet\) & \(\bullet\) & \(\bullet\) & \(\bullet\) & si remédiation\\
Détection auto & --- & --- & --- & \(\bullet\) & ---\\
Investigation & \(\bullet\) & \(\bullet\) & \(\bullet\) & \(\bullet\) & \(\bullet\)\\
Confirmation & --- & --- & \(\bullet\) & \(\bullet\) & \(\bullet\)\\
Data Hub & --- & \(\bullet\) & --- & --- & ---\\
Demandes & --- & --- & --- & --- & ---\\
Recherche & --- & --- & --- & --- & ---\\
Rapports & --- & --- & --- & \(\bullet\) & ---\\
Cross-domain & --- & --- & --- & --- & ---\\
Métriques & --- & --- & --- & \(\bullet\) & ---\\
\bottomrule
\end{longtable}

% =====================================================================
\chapter{Le Model Context Protocol dans ce projet}
\label{ch:mcp}

Ce chapitre répond point par point aux questions posées sur le MCP, avec le
code réel sous les yeux.

% ---------------------------------------------------------------------
\section{Qu'est-ce que le MCP Server dans \emph{ce} projet ?}

C'est un \textbf{processus Python séparé} qui expose les dix fonctions
d'exploitation Storage sous forme d'outils décrits, interrogeables par un modèle
de langage. Il n'est pas un serveur réseau : il ne écoute aucun port, il
dialogue par son entrée et sa sortie standard.

\begin{keybox}
\textbf{En une phrase.} Le MCP Server d'OASIS est le \emph{catalogue exécutable}
des capacités Storage : il dit au modèle quels outils existent, avec quels
paramètres, et il les exécute quand on le lui demande --- sans jamais décider
lui-même de ce qui est permis.
\end{keybox}

\noindent
La dernière partie de cette phrase est importante et le code l'assume
explicitement : \file{server.py} porte le commentaire « Applique par le persona
et la route Actions ; le serveur MCP reste neutre. » La politique de sécurité
vit \emph{au-dessus} du serveur, pas dedans.

% ---------------------------------------------------------------------
\section{Où se trouve son code ?}

\begin{longtable}{@{}p{4.6cm}p{2.0cm}p{7.8cm}@{}}
\toprule
\textbf{Fichier} & \textbf{Lignes} & \textbf{Rôle}\\
\midrule
\endhead
\file{app/mcp/storage/tool\_server.py} & 12 & Le serveur : instancie \cls{MCPServer} et enregistre les outils\\
\file{app/mcp/storage/server.py} & 60 & Le registre \fn{TOOL\_REGISTRY} et la classification du risque\\
\file{app/mcp/storage/client.py} & 118 & Le client \cls{StorageMCPClient} et le point d'audit\\
\file{app/mcp/storage/ssh.py} & 130 & L'accès SSH mutualisé par tous les outils\\
\file{app/mcp/storage/tools/*.py} & 7 fichiers & Les implémentations\\
\bottomrule
\end{longtable}

% ---------------------------------------------------------------------
\section{Comment est-il lancé ?}

Il n'est \textbf{jamais} lancé à la main. Le client le démarre comme
sous-processus, à la demande.

\begin{lstlisting}[language=Python,caption={app/mcp/storage/client.py --- paramètres de lancement}]
_BACKEND_DIR = Path(__file__).resolve().parents[3]      # -> backend/
_SERVER_PARAMS = StdioServerParameters(
    command=sys.executable,                              # le meme interpreteur Python
    args=["-m", "app.mcp.storage.tool_server"],          # module lance en -m
    cwd=str(_BACKEND_DIR),                               # pour que "app.*" soit importable
)
\end{lstlisting}

\begin{lstlisting}[language=Python,caption={app/mcp/storage/tool\_server.py --- le serveur complet}]
from mcp.server.mcpserver import MCPServer
from app.mcp.storage.server import TOOL_REGISTRY

mcp = MCPServer("oasis-storage-mcp")
for _tool in TOOL_REGISTRY.values():
    mcp.tool()(_tool)

if __name__ == "__main__":
    mcp.run(transport="stdio")
\end{lstlisting}

\subsection*{Cycle de vie côté client}

\begin{lstlisting}[language=Python,caption={Démarrage et arrêt}]
async def start(self) -> None:
    if self._session is not None:
        return
    self._stack = AsyncExitStack()
    read, write = await self._stack.enter_async_context(stdio_client(_SERVER_PARAMS))
    self._session = await self._stack.enter_async_context(ClientSession(read, write))
    await self._session.initialize()
\end{lstlisting}

\begin{keybox}
\textbf{Pourquoi aucune instance partagée au niveau du module.} Le commentaire
final de \file{client.py} l'explique : « une instance créée à l'import n'est
jamais démarrée, et un client non démarré échoue à chaque appel d'outil. Chaque
appelant ouvre le sien pour la durée de son travail. »

La route de confirmation ajoute l'argument économique : « Une confirmation est
un clic humain, donc rare : le coût du démarrage du sous-processus est
négligeable, et on évite le sous-processus partagé qui reste mort jusqu'au
prochain redémarrage du backend s'il tombe. »
\end{keybox}

% ---------------------------------------------------------------------
\section{Quels outils sont disponibles ?}

\begin{longtable}{@{}p{4.4cm}p{3.3cm}p{2.0cm}p{4.4cm}@{}}
\toprule
\textbf{Outil} & \textbf{Paramètres} & \textbf{Classe} & \textbf{Commande sur la VM}\\
\midrule
\endhead
\fn{get\_capacity} & \texttt{volume\_id?} & lecture & \texttt{df -T -B1 -P}\\
\fn{list\_backups} & \texttt{target?} & lecture & \texttt{systemctl list-timers} + \texttt{systemctl show}\\
\fn{get\_backup\_job\_status} & \texttt{target?} & lecture & alias de \fn{list\_backups}\\
\fn{get\_backup\_logs} & \texttt{job\_id}, \texttt{lines=20} & lecture & \texttt{journalctl -u ... -o short-iso}\\
\fn{list\_snapshots} & \texttt{volume\_id?} & lecture & \texttt{lvs -{}-noheadings ...}\\
\fn{get\_storage\_health} & --- & lecture & agrège capacité + sauvegardes\\
\fn{run\_backup} & \texttt{target} & \textbf{remédiation} & \texttt{systemctl start <service>}\\
\fn{create\_snapshot} & \texttt{volume\_id}, \texttt{volume\_group?}, \texttt{size\_gb=1.0} & \textbf{remédiation} & \texttt{sudo lvcreate -{}-snapshot}\\
\fn{repair\_backup\_repository} & \texttt{target}, \texttt{confirm=False} & \textbf{remédiation} & drop-in systemd + \texttt{daemon-reload}\\
\fn{restore\_from\_backup} & \texttt{target}, \texttt{backup\_id}, \texttt{confirm=False} & \textbf{destructif} & \texttt{sudo lvconvert -{}-merge}\\
\bottomrule
\end{longtable}

\begin{warnbox}
\textbf{\fn{sample\_disk\_performance} n'est pas un outil MCP.} Il vit dans
\file{tools/performance.py} mais n'est \emph{pas} dans le \fn{TOOL\_REGISTRY}.
Le modèle ne le voit donc jamais ; seul le collecteur de métriques l'appelle,
directement.
\end{warnbox}

% ---------------------------------------------------------------------
\section{Comment les outils sont-ils enregistrés ?}

En deux temps.

\subsection*{Temps 1 --- le registre}
\file{server.py} importe les fonctions et les place dans un dictionnaire :

\begin{lstlisting}[language=Python,caption={app/mcp/storage/server.py}]
TOOL_REGISTRY: dict[str, Callable[..., dict[str, Any]]] = {
    "get_capacity": get_capacity,
    "list_backups": list_backups,
    "run_backup": run_backup,
    "get_backup_job_status": get_backup_job_status,
    "get_backup_logs": get_backup_logs,
    "list_snapshots": list_snapshots,
    "create_snapshot": create_snapshot,
    "get_storage_health": get_storage_health,
    "repair_backup_repository": repair_backup_repository,
    "restore_from_backup": restore_from_backup,
}
\end{lstlisting}

\subsection*{Temps 2 --- le décorateur, appliqué en boucle}
\texttt{mcp.tool()(\_tool)} est l'application \emph{manuelle} du décorateur
\fn{@mcp.tool()}. Le SDK dérive alors, pour chaque fonction :
\begin{itemize}
  \item le \textbf{nom} de l'outil, depuis \texttt{\_\_name\_\_} ;
  \item la \textbf{description}, depuis le docstring ;
  \item le \textbf{schéma JSON des paramètres}, depuis les annotations de type
  et les valeurs par défaut.
\end{itemize}

\begin{keybox}
\textbf{Conséquence concrète et vérifiable.} Le docstring de
\fn{repair\_backup\_repository} est rédigé \emph{en anglais et à destination du
modèle} : il commence par « Repair a backup job whose restic repository cannot
be opened. Use this when \ldots{} the journal shows \texttt{unable to open
config file}. » Ce texte n'est pas de la documentation interne : c'est la
description que le LLM lit pour décider s'il doit appeler cet outil.
\end{keybox}

% ---------------------------------------------------------------------
\section{Comment le client communique-t-il avec le serveur ?}

\begin{figure}[h!]
\centering
\begin{tikzpicture}[node distance=8mm and 20mm]
\node[ambox,minimum width=34mm,minimum height=30mm] (cli) {\textbf{Processus backend}\\[2mm]\tiny uvicorn\\\tiny StoragePersona\\\tiny StorageMCPClient\\\tiny (asyncio.Lock)};
\node[greenbox,right=of cli,minimum width=34mm,minimum height=30mm] (srv) {\textbf{Sous-processus}\\[2mm]\tiny python -m app.mcp.storage.tool\_server\\\tiny MCPServer("oasis-storage-mcp")\\\tiny 10 outils enregistrés};
\node[redbox,right=of srv,minimum width=30mm,minimum height=22mm] (vm) {\textbf{VM Ubuntu}\\[2mm]\tiny sshd\\\tiny df / lvs / systemctl};

\draw[flowb] ([yshift=6mm]cli.east) -- node[above,lbl]{stdin : JSON-RPC} ([yshift=6mm]srv.west);
\draw[flowb] ([yshift=-6mm]srv.west) -- node[below,lbl]{stdout : résultat} ([yshift=-6mm]cli.east);
\draw[flowr] (srv) -- node[above,lbl]{paramiko / SSH 22} (vm);
\end{tikzpicture}
\caption{Le transport stdio : pas de socket, pas de port, un tube entre deux processus.}
\end{figure}

\subsection*{\fn{list\_tools\_for\_llm} --- traduction du catalogue}

\begin{lstlisting}[language=Python,caption={Du format MCP au format attendu par Groq}]
async def list_tools_for_llm(self) -> list[dict[str, Any]]:
    response = await self._require_session().list_tools()
    return [
        {"type": "function", "function": {
            "name": tool.name,
            "description": tool.description or tool.name,
            "parameters": getattr(tool, "inputSchema", None)
                          or getattr(tool, "input_schema", {}),
        }}
        for tool in response.tools
    ]
\end{lstlisting}

\noindent
Le double \fn{getattr} absorbe la différence de nommage entre versions du SDK
(\texttt{inputSchema} en camelCase, \texttt{input\_schema} en snake\_case).

\subsection*{\fn{call\_tool} --- l'appel et son audit}

\begin{enumerate}
  \item \textbf{Verrou} --- \texttt{async with self.\_lock} : une seule requête
  MCP à la fois sur une session partagée, le transport stdio n'étant pas
  multiplexé.
  \item \textbf{Appel} --- \fn{session.call\_tool(tool\_name, arguments)}.
  \item \textbf{Détection d'erreur} --- si \texttt{isError} (ou
  \texttt{is\_error}) est vrai, les blocs de texte sont concaténés et levés en
  \cls{RuntimeError}.
  \item \textbf{Décodage} --- le contenu textuel est passé à
  \texttt{json.loads} ; s'il est vide, on retombe sur
  \texttt{raw.model\_dump(mode="json")}.
  \item \textbf{Audit} --- si \texttt{db} et \texttt{user\_id} sont fournis,
  \fn{log\_mcp\_call} écrit dans \tbl{mcp\_calls}. Cet appel est lui-même
  entouré d'un \texttt{try/except} : « Un échec d'audit ne doit pas masquer le
  résultat de l'outil. »
  \item \textbf{Propagation} --- si le statut est \texttt{FAILED}, une
  \cls{RuntimeError} est levée \emph{après} l'écriture de l'audit.
\end{enumerate}

% ---------------------------------------------------------------------
\section{Comment une requête utilisateur arrive-t-elle jusqu'au MCP ?}

\begin{lstlisting}[caption={Trajet complet, fonction par fonction}]
1.  ChatPage.send()                       [frontend]
2.  api.post('/storage/chat', {...})      [frontend]
3.  chat()                                [api/v1/storage.py]
4.  storage_persona.handle_message()      [personas/storage/agent.py]
5.  async with StorageMCPClient() as mcp  [ouverture du sous-processus]
6.  mcp.list_tools_for_llm()              [catalogue -> format Groq]
7.  groq_client.chat_with_tools(...)      [le modele choisit un outil]
8.  mcp.call_tool(name, **arguments)      [client -> serveur]
9.  MCPServer dispatche vers la fonction  [sous-processus]
10. tools/capacity.get_capacity()         [implementation]
11. ssh.run("df -T -B1 -P")               [paramiko]
12. VM Ubuntu execute la commande
13. retour JSON -> mcp_calls -> messages  [audit + reinjection]
14. groq_client redige la reponse finale
15. ChatTurnOut -> React
\end{lstlisting}

% ---------------------------------------------------------------------
\section{Comment un outil est-il sélectionné ?}

\textbf{Par le modèle, et par lui seul}, à partir du catalogue reçu. Le code ne
contient aucun classifieur, aucune correspondance mot-clé \arrowto{} outil,
aucune heuristique de routage.

\begin{lstlisting}[language=Python,caption={La seule instruction de sélection}]
options.update({"tools": tools, "tool_choice": "auto"})
\end{lstlisting}

\noindent
Ce que le modèle voit pour décider : le prompt système, la date du serveur, le
contexte RAG, le bloc incident éventuel, l'historique, la question, et la
description de chaque outil. Ce que le code \emph{contraint} ensuite : ce qui
sera exécuté et ce qui sera seulement proposé.

% ---------------------------------------------------------------------
\section{Comment les paramètres sont-ils transmis ?}

\begin{enumerate}
  \item Le modèle renvoie \texttt{call.function.arguments}, une \textbf{chaîne
  JSON}.
  \item L'agent la décode dans un \texttt{try/except json.JSONDecodeError} et
  retombe sur \texttt{\{\}} en cas d'échec ; un résultat qui ne serait pas un
  dictionnaire est également ramené à \texttt{\{\}}.
  \item Les arguments sont passés en \texttt{**kwargs} à \fn{call\_tool}, puis
  transmis au serveur MCP, qui les applique à la fonction Python.
  \item Pour une action confirmée, les paramètres viennent de la colonne JSONB
  \texttt{action\_requests.parameters}, à laquelle
  \fn{confirm\_and\_execute\_action} ajoute \texttt{confirm=True}.
\end{enumerate}

\subsection*{Échappement avant la ligne de commande}
Tout argument qui finit dans une commande shell passe par \fn{ssh.quote}, un
alias de \texttt{shlex.quote}. Exemples réels :

\begin{lstlisting}[language=Python]
run(f"systemctl start {quote(service)}")
run(f"journalctl -u {quote(service)} -n {int(lines)} --no-pager -o short-iso")
run(f"sudo lvcreate --snapshot --size {float(size_gb)}G "
    f"--name {quote(snapshot_name)} {quote(source)}")
run(f"sudo lvconvert --merge {quote(backup_id)}")
\end{lstlisting}

\begin{keybox}
Les valeurs numériques sont converties (\texttt{int(lines)},
\texttt{float(size\_gb)}) plutôt que citées : une conversion ratée lève une
exception avant d'atteindre la VM. Les valeurs textuelles sont citées. Aucun
argument venu du modèle n'est concaténé brut dans une commande.
\end{keybox}

% ---------------------------------------------------------------------
\section{Comment le tool exécute-t-il l'action ?}

Prenons \fn{run\_backup} en entier :

\begin{lstlisting}[language=Python,caption={app/mcp/storage/tools/backup.py}]
def run_backup(target: str) -> dict:
    """
    Déclenche immédiatement un job de sauvegarde existant.
    Action non destructive (ajout de données), pas de confirmation requise.
    """
    service = _service_for(target) if target.endswith(".timer") else target
    try:
        run(f"systemctl start {quote(service)}")
    except StorageVMError as error:
        return {"error": str(error), "action": "backup_failed", "target": target}
    return {"action": "backup_started", "target": service, "status": "running"}
\end{lstlisting}

\noindent
Trois choses à remarquer :
\begin{itemize}
  \item la conversion \texttt{.timer} \arrowto{} \texttt{.service} est faite dans
  l'outil, le modèle n'a pas à la connaître ;
  \item l'erreur est \textbf{retournée}, pas levée --- c'est ce qui permet à
  \fn{confirm\_and\_execute\_action} de tester \texttt{result.get("error")} ;
  \item le docstring dit « pas de confirmation requise », mais
  \fn{run\_backup} appartient tout de même à \texttt{REMEDIATION\_TOOLS} : la
  politique de l'application est plus stricte que celle de l'outil.
\end{itemize}

\subsection*{La couche SSH}
\fn{ssh.run(command, allow\_failure=False)} : ouvre une connexion, exécute,
lit \texttt{stdout}, récupère le code de retour, lit \texttt{stderr}, ferme
dans un \texttt{finally}. Un code non nul lève \cls{StorageVMError}, sauf si
\texttt{allow\_failure=True} --- utile « aux commandes dont un code de retour non
nul est une réponse valide (\texttt{systemctl} sur une unit en échec, par
exemple) ».

\noindent
\fn{ssh.run\_many(commands)} exécute plusieurs commandes sur \emph{une seule}
connexion : « Ouvrir une session par commande coûte cher. »

% ---------------------------------------------------------------------
\section{Comment le résultat revient-il ?}

\begin{enumerate}
  \item L'outil renvoie un \texttt{dict} Python.
  \item Le SDK MCP le sérialise en JSON dans un bloc de contenu textuel.
  \item Le client concatène les blocs et exécute \texttt{json.loads}.
  \item Il l'écrit dans \tbl{mcp\_calls} (colonne \texttt{result}, JSONB).
  \item L'agent le passe à \fn{evaluate\_tool\_result} pour la détection
  d'anomalie.
  \item Il le réinjecte dans \texttt{messages} sous forme d'un message de rôle
  \texttt{tool}, avec le \texttt{tool\_call\_id} correspondant.
  \item Le modèle rédige la réponse au tour suivant.
\end{enumerate}

% ---------------------------------------------------------------------
\section{Comment les erreurs sont-elles gérées ?}

\begin{longtable}{@{}p{3.5cm}p{4.2cm}p{6.9cm}@{}}
\toprule
\textbf{Niveau} & \textbf{Erreur} & \textbf{Traitement}\\
\midrule
\endhead
SSH & Config incomplète & \cls{StorageVMError} : « SSH settings are incomplete in .env: \ldots »\\
SSH & VM injoignable & \cls{StorageVMError} : « Could not reach <host>: \ldots »\\
SSH & Commande en échec & \cls{StorageVMError} avec \texttt{stderr} ou le code de retour\\
Outil & \cls{StorageVMError} attrapée & \textbf{Retour} d'un dict \texttt{\{"error": ...\}}, jamais une exception\\
Client MCP & \texttt{isError} ou exception & \texttt{status=FAILED}, audit écrit, puis \cls{RuntimeError}\\
Agent & \cls{RuntimeError} & Convertie en \texttt{\{"status": "error", "message": ...\}} et réinjectée dans la conversation --- le modèle peut réagir\\
Route confirm & Serveur MCP non démarrable & HTTP 503, action laissée \texttt{PENDING} donc rejouable\\
Route confirm & Action introuvable ou déjà traitée & \cls{ValueError} \arrowto{} HTTP 404\\
Service & Résultat contenant \texttt{error} & Action marquée \texttt{FAILED}, pas \texttt{COMPLETED}\\
Audit & Échec d'écriture & \fn{logger.exception}, le résultat de l'outil est préservé\\
\bottomrule
\end{longtable}

\begin{keybox}
\textbf{La règle qui traverse toutes ces lignes.} Une erreur d'infrastructure ne
doit jamais devenir une donnée inventée, et ne doit jamais faire taire le
système. Elle remonte, nommée, jusqu'à l'écran.
\end{keybox}

% ---------------------------------------------------------------------
\section{Comment les permissions sont-elles gérées ?}

Quatre niveaux se superposent, et \textbf{aucun n'est dans le serveur MCP}.

\begin{longtable}{@{}p{2.2cm}p{4.0cm}p{8.4cm}@{}}
\toprule
\textbf{Niveau} & \textbf{Où} & \textbf{Ce qui est contrôlé}\\
\midrule
\endhead
1 & \fn{get\_current\_user} & Jeton valide, compte actif\\
2 & \fn{require\_role} & Rôle autorisé sur la route\\
3 & \texttt{REMEDIATION\_TOOLS} + \fn{\_propose\_action} & Le chat ne peut pas exécuter un outil qui modifie la VM\\
4 & \texttt{confirm=True} dans les outils & \fn{restore\_from\_backup} et \fn{repair\_backup\_repository} refusent d'agir sans ce drapeau, seul \route{/confirm} l'accorde\\
\bottomrule
\end{longtable}

\begin{warnbox}
\textbf{Deux constats de lecture, à énoncer tels quels.}

\emph{(a)} \texttt{DESTRUCTIVE\_TOOLS} et la fonction \fn{risk\_for} sont
définis dans \file{server.py} mais \textbf{ne sont appelés nulle part} dans le
dépôt. La qualification « HIGH / LOW » n'est donc pas écrite en base
aujourd'hui. Le frontend possède sa \emph{propre} liste
(\texttt{DESTRUCTIVE\_ACTIONS = new Set(['restore\_from\_backup'])}) pour
afficher l'avertissement rouge --- une duplication à signaler.

\emph{(b)} La fonction \fn{invoke} de \file{server.py} n'est pas utilisée non
plus : le serveur MCP enregistre les fonctions directement, il ne passe pas par
elle.
\end{warnbox}

% ---------------------------------------------------------------------
\section{L'architecture LLM \arrowto{} Agent \arrowto{} MCP Client \arrowto{} MCP Server \arrowto{} Tool \arrowto{} Infrastructure correspond-elle au code ?}

\textbf{Oui, avec deux précisions.}

\subsection*{Correspondance exacte}

\begin{longtable}{@{}p{3.4cm}p{4.6cm}p{6.6cm}@{}}
\toprule
\textbf{Couche théorique} & \textbf{Réalisation dans le code} & \textbf{Fichier}\\
\midrule
\endhead
LLM & Groq, \texttt{openai/gpt-oss-120b} & \file{llm/groq\_client.py}\\
Agent & \cls{StoragePersona} & \file{personas/storage/agent.py}\\
MCP Client & \cls{StorageMCPClient} & \file{mcp/storage/client.py}\\
MCP Server & \cls{MCPServer("oasis-storage-mcp")} & \file{mcp/storage/tool\_server.py}\\
Tool & 10 fonctions Python & \file{mcp/storage/tools/}\\
Infrastructure & VM Ubuntu par SSH & \file{mcp/storage/ssh.py}\\
\bottomrule
\end{longtable}

\subsection*{Précision 1 --- l'agent n'est pas un simple relais}
Entre l'Agent et le MCP Client s'intercale un \textbf{garde-fou} : quatre des
dix outils ne traversent jamais cette frontière depuis le chat. Le schéma
générique laisserait croire que tout appel du modèle atteint l'infrastructure.
Ici, ce n'est vrai que pour les six outils de lecture.

\subsection*{Précision 2 --- une voie parallèle contourne MCP}
Le dashboard, la surveillance et le collecteur appellent les mêmes fonctions
\emph{sans} passer par le client ni le serveur MCP. La chaîne complète n'est
donc pas la seule voie d'accès à l'infrastructure --- c'est la voie de
\emph{l'agent}.

\begin{figure}[h!]
\centering
\begin{tikzpicture}[node distance=6mm and 14mm]
\node[vibox,minimum width=26mm] (llm) {LLM (Groq)};
\node[vibox,right=of llm,minimum width=26mm] (ag) {Agent\\\tiny StoragePersona};
\node[ambox,right=of ag,minimum width=26mm] (cl) {MCP Client};
\node[ambox,right=of cl,minimum width=26mm] (sv) {MCP Server};
\node[greenbox,below=14mm of sv,minimum width=26mm] (to) {Tool Python};
\node[redbox,below=of to,minimum width=26mm] (in) {Infrastructure\\\tiny VM Ubuntu};
\node[bluebox,left=of to,minimum width=30mm] (direct) {Dashboard · Watcher\\Collector\\\tiny appel direct};

\draw[flow,<->] (llm)--(ag);
\draw[flow] (ag)--(cl); \draw[flow] (cl)--(sv); \draw[flow] (sv)--(to);
\draw[flowr] (to)--(in);
\draw[flowb] (direct)--(to);
\node[lbl,below=1mm of ag] {garde-fou :\\4 outils arrêtés ici};
\end{tikzpicture}
\caption{Les deux voies d'accès aux outils.}
\end{figure}

% =====================================================================
\chapter{Actions dangereuses et confirmation humaine}
\label{ch:actions}

Ce chapitre recense \emph{toutes} les fonctions du dépôt capables de modifier un
système, de supprimer quelque chose, de démarrer un service, de toucher à un
volume, à un snapshot, à une sauvegarde ou à une configuration --- puis détaille
le mécanisme de confirmation.

% =====================================================================
\section{Inventaire des fonctions à effet de bord}

\subsection{Groupe A --- Modification de l'infrastructure (VM)}

\begin{longtable}{@{}p{3.6cm}p{3.5cm}p{7.3cm}@{}}
\toprule
\textbf{Fonction} & \textbf{Fichier} & \textbf{Ce qu'elle fait réellement}\\
\midrule
\endhead
\fn{run\_backup(target)} & \file{tools/backup.py} & \texttt{systemctl start <service>} --- démarre un job de sauvegarde existant\\
\fn{create\_snapshot(volume\_id, volume\_group, size\_gb)} & \file{tools/snapshot.py} & \texttt{sudo lvcreate -{}-snapshot -{}-size <n>G -{}-name <nom> /dev/<vg>/<lv>}\\
\fn{restore\_from\_backup(target, backup\_id, confirm)} & \file{tools/restore.py} & \texttt{sudo lvconvert -{}-merge <snapshot>} --- \textbf{écrase le volume}\\
\fn{repair\_backup\_repository(target, confirm)} & \file{tools/repair.py} & Écrit un drop-in systemd et recharge le démon\\
\fn{\_write\_drop\_in(service, body)} & \file{tools/repair.py} & \texttt{sudo -n install -d}, \texttt{sudo -n tee}, \texttt{systemctl daemon-reload}\\
\fn{ssh.run(command, allow\_failure)} & \file{mcp/storage/ssh.py} & \textbf{Exécute une commande arbitraire} sur la VM\\
\fn{ssh.run\_many(commands)} & \file{mcp/storage/ssh.py} & Idem, plusieurs commandes sur une session\\
\bottomrule
\end{longtable}

\begin{keybox}
\textbf{Point de sécurité important.} \fn{ssh.run} est la seule porte vers la
VM, et \textbf{aucune commande n'est composée à partir d'une chaîne fournie par
le modèle}. Le modèle fournit des \emph{arguments} (un nom d'unit, un nom de
volume), jamais un verbe. Les commandes sont écrites en dur dans les outils, et
chaque argument passe par \fn{quote} (\texttt{shlex.quote}) ou par une
conversion numérique.
\end{keybox}

\subsection{Groupe B --- Suppressions applicatives}

\begin{longtable}{@{}p{3.9cm}p{3.4cm}p{7.1cm}@{}}
\toprule
\textbf{Fonction} & \textbf{Fichier} & \textbf{Ce qu'elle supprime}\\
\midrule
\endhead
\fn{delete\_conversation} & \file{services/chat\_service.py} & La conversation et ses messages (\texttt{ON DELETE CASCADE})\\
\fn{delete\_document} & \file{services/knowledge\_service.py} & Les chunks dans Chroma, \emph{puis} le fichier sur disque, \emph{puis} la ligne en base\\
\fn{delete\_request} & \file{services/request\_service.py} & La demande, ses messages, son journal, ses pièces jointes, et les binaires sur disque\\
\fn{delete\_attachment} & \file{services/request\_service.py} & Une pièce jointe et son binaire\\
\fn{purge\_old} & \file{services/metrics\_service.py} & Les \tbl{disk\_metrics} de plus de 31 jours\\
\bottomrule
\end{longtable}

\begin{infobox}
\textbf{Aucune de ces suppressions ne touche l'infrastructure.} Elles restent
dans le périmètre applicatif. Elles ne passent donc pas par le mécanisme
d'\cls{ActionRequest} --- mais elles sont toutes protégées par
\fn{require\_role} et, pour les demandes et conversations, par un contrôle de
propriété.
\end{infobox}

\subsection{Groupe C --- Modifications d'état applicatif}

\fn{confirm\_and\_execute\_action}, \fn{reject\_action\_request},
\fn{resolve\_ticket}, \fn{update\_ticket\_status}, \fn{update\_status}
(demandes), \fn{approve\_request} / \fn{reject\_request} (cross-domain).
Toutes écrivent une trace : note d'incident ou ligne \tbl{request\_events}.

% =====================================================================
\section{Fiche détaillée de chaque action d'infrastructure}

% ---------------------------------------------------
\subsection{\fn{run\_backup(target)}}

\begin{longtable}{@{}p{3.2cm}p{11.6cm}@{}}
\toprule
\endhead
\textbf{Où} & \file{app/mcp/storage/tools/backup.py}\\
\textbf{Qui peut l'appeler} & Le modèle, mais l'appel est \textbf{intercepté} : \texttt{run\_backup} $\in$ \texttt{REMEDIATION\_TOOLS}. Exécution réelle uniquement via \route{POST /storage/actions/\{id\}/confirm}\\
\textbf{Paramètres} & \texttt{target: str} --- nom du timer (\texttt{backup-app-critical.timer}) ou du service\\
\textbf{Effet réel} & Convertit \texttt{.timer} en \texttt{.service}, puis \texttt{systemctl start <service>} par SSH\\
\textbf{Confirmation} & \textbf{Oui}, par la politique de l'application. Le docstring de l'outil dit pourtant « pas de confirmation requise » : l'outil se juge inoffensif, l'application le soumet quand même au clic humain\\
\textbf{Retour} & \texttt{\{"action": "backup\_started", "target", "status": "running"\}} ou \texttt{\{"error": ..., "action": "backup\_failed"\}}\\
\bottomrule
\end{longtable}

% ---------------------------------------------------
\subsection{\fn{create\_snapshot(volume\_id, volume\_group, size\_gb)}}

\begin{longtable}{@{}p{3.2cm}p{11.6cm}@{}}
\toprule
\endhead
\textbf{Où} & \file{app/mcp/storage/tools/snapshot.py}\\
\textbf{Qui} & Même circuit : proposition puis confirmation\\
\textbf{Paramètres} & \texttt{volume\_id} (obligatoire), \texttt{volume\_group} (déduit par \texttt{lvs} s'il est absent), \texttt{size\_gb=1.0}\\
\textbf{Effet réel} & Nom horodaté \texttt{<lv>-snap-AAAAMMJJ-HHMMSS}, puis \texttt{sudo lvcreate -{}-snapshot -{}-size <n>G -{}-name <nom> /dev/<vg>/<lv>}\\
\textbf{Confirmation} & Oui (appartenance à \texttt{REMEDIATION\_TOOLS}). L'outil lui-même n'exige pas \texttt{confirm=True}\\
\textbf{Nature} & Non destructive : « elle ajoute un volume, elle n'écrase rien »\\
\textbf{Garde-fou interne} & Si le logical volume n'existe pas : \texttt{\{"error": "Logical volume '<x>' not found on the VM"\}}, sans rien exécuter\\
\bottomrule
\end{longtable}

% ---------------------------------------------------
\subsection{\fn{restore\_from\_backup(target, backup\_id, confirm)} --- la plus dangereuse}

\begin{longtable}{@{}p{3.2cm}p{11.6cm}@{}}
\toprule
\endhead
\textbf{Où} & \file{app/mcp/storage/tools/restore.py}\\
\textbf{Classement} & Seul membre de \texttt{DESTRUCTIVE\_TOOLS}, et membre de \texttt{REMEDIATION\_TOOLS}\\
\textbf{Paramètres} & \texttt{target} (volume), \texttt{backup\_id} (nom du snapshot LVM), \texttt{confirm=False}\\
\textbf{Effet réel} & \texttt{sudo lvconvert -{}-merge <snapshot>} --- « elle écrase l'état courant du volume par celui du snapshot » et supprime le snapshot\\
\textbf{Double verrou} & \emph{(1)} interception par \texttt{REMEDIATION\_TOOLS} ; \emph{(2)} l'outil lui-même refuse d'agir sans \texttt{confirm=True}\\
\textbf{Sans confirmation} & Retourne \texttt{\{"action": "restore\_requires\_confirmation", "warning": "\ldots{} écrasera les données actuelles du volume et supprimera le snapshot"\}} --- \textbf{rien n'est exécuté}\\
\textbf{Note technique} & LVM ne fusionne qu'au prochain démontage/remontage : le retour le dit et invite à vérifier avec \texttt{lvs}\\
\bottomrule
\end{longtable}

\begin{lstlisting}[language=Python,caption={Le refus par défaut, dans l'outil lui-même}]
def restore_from_backup(target: str, backup_id: str, confirm: bool = False) -> dict:
    if not confirm:
        return {
            "action": "restore_requires_confirmation",
            "target": target, "backup_id": backup_id,
            "warning": (f"Fusionner le snapshot '{backup_id}' dans '{target}' ecrasera "
                        "les donnees actuelles du volume et supprimera le snapshot. "
                        "Confirmation explicite requise."),
        }
    output = run(f"sudo lvconvert --merge {quote(backup_id)}")
\end{lstlisting}

% ---------------------------------------------------
\subsection{\fn{repair\_backup\_repository(target, confirm)}}

C'est l'outil le plus élaboré du projet, et le seul qui \emph{modifie une
configuration système}. Il mérite un développement.

\subsubsection*{Le problème qu'il traite}
Le docstring l'expose : « Jusqu'ici OASIS savait relancer un job en échec, pas
corriger ce qui le faisait échouer. Sur un job dont le dépôt restic est
introuvable, relancer ne pouvait que reproduire la même erreur : l'ingénieur
voyait un diagnostic juste suivi d'une remédiation vouée à rater. »

\subsubsection*{Les trois principes déclarés}
\begin{enumerate}
  \item \textbf{Il ne devine rien.} Le dépôt de remplacement est \emph{lu sur la
  machine}, pas fourni par le modèle, et seuls les jobs dont la dernière
  exécution a réussi servent de référence.
  \item \textbf{Il ne touche pas au fichier d'unité.} Un drop-in s'annule d'un
  \texttt{systemctl revert}, ce que la réécriture d'une unité ne permet pas.
  \item \textbf{Il refuse plutôt que d'approximer.}
\end{enumerate}

\subsubsection*{Déroulé, étape par étape}

\begin{longtable}{@{}p{0.7cm}p{4.6cm}p{9.4cm}@{}}
\toprule
\textbf{\#} & \textbf{Étape} & \textbf{Détail}\\
\midrule
\endhead
1 & Lecture de la configuration & \texttt{systemctl show <service> -p Environment -p EnvironmentFiles -p FragmentPath}\\
2 & Unit inexistante & \texttt{\{"action": "repair\_failed", "error": "No systemd unit named \ldots"\}}\\
3 & Pas de \texttt{RESTIC\_REPOSITORY} direct & \texttt{"no\_repair\_possible"} : la configuration vient d'un fichier que le compte SSH ne peut pas lire\\
4 & Test du dépôt & \fn{\_repository\_state} lance \texttt{restic cat config} avec un \textbf{faux mot de passe} volontaire\\
5 & Interprétation & \texttt{"version"} dans la sortie \arrowto{} \texttt{valid} ; « wrong password » \arrowto{} \texttt{protected} ; « no such file » \arrowto{} \texttt{missing} ; sinon \texttt{unknown}\\
6 & Dépôt valide ou protégé & \texttt{"no\_repair\_needed"} : « The job fails for another reason --- read its journal with \fn{get\_backup\_logs} »\\
7 & Recherche du donneur & \fn{\_healthy\_donor} liste les services de sauvegarde, écarte les units de la distribution, lit leur \texttt{EnvironmentFiles}\\
8 & Critères stricts & Le fichier doit être celui du plus grand nombre de jobs, \textbf{strictement} (égalité \arrowto{} abstention), et au moins un job l'utilisant doit avoir réussi\\
9 & Pas de donneur & \texttt{"no\_repair\_possible"} avec la liste des voisins examinés\\
10 & \textbf{Sans \texttt{confirm}} & \texttt{"repair\_requires\_confirmation"} avec un champ \texttt{plan} en langage clair\\
11 & Avec \texttt{confirm=True} & Écriture du drop-in, \texttt{systemctl daemon-reload}, puis relecture pour vérifier la prise en compte\\
12 & Refus de sudo & \fn{\_missing\_privilege} traduit l'erreur en consigne exploitable (règle sudoers précise, ou marche à suivre manuelle)\\
\bottomrule
\end{longtable}

\begin{keybox}
\textbf{L'astuce du faux mot de passe.} \fn{\_repository\_state} passe
\texttt{-{}-password-command 'printf x'} exprès : « restic lit le fichier
\texttt{config} du dépôt avant de vérifier la clé, donc l'absence de dépôt et le
mauvais mot de passe produisent deux erreurs distinctes. C'est ce qui permet de
distinguer "rien ici" de "dépôt bien présent" sans manipuler le moindre
secret. » C'est un excellent exemple de conception défensive à citer.
\end{keybox}

\subsubsection*{Réversibilité}
Le drop-in écrit contient un en-tête qui identifie son auteur et donne la
commande d'annulation :

\begin{lstlisting}[caption={Contenu du fichier /etc/systemd/system/<service>.d/50-oasis-repair.conf}]
# Written by OASIS on 2026-08-23 14:12 UTC (repair_backup_repository).
# backup-app-critical.service pointed at /srv/repo, which is not a restic repository.
# It now reads the environment shared by this host's healthy backup jobs.
# Undo with: sudo systemctl revert backup-app-critical.service
[Service]
EnvironmentFile=/etc/oasis/restic.env
\end{lstlisting}

% =====================================================================
\section{Le mécanisme de confirmation humaine}

\subsection{Où il est implémenté}

\begin{longtable}{@{}p{4.4cm}p{10.4cm}@{}}
\toprule
\textbf{Couche} & \textbf{Implémentation}\\
\midrule
\endhead
Classification & \texttt{REMEDIATION\_TOOLS} dans \file{mcp/storage/server.py}\\
Interception & \fn{StoragePersona.\_propose\_action} dans \file{personas/storage/agent.py}\\
Persistance & Table \tbl{action\_requests}, statut \texttt{pending}\\
Rattachement & \tbl{action\_requests.chat\_message\_id} via \fn{chat\_service.attach\_actions}\\
Exposition & \cls{ActionRequestOut} dans \cls{ChatMessageOut} et \cls{IncidentTicketDetail}\\
Interface & \cls{ActionProposal} (\file{frontend/src/components/ActionProposal.tsx})\\
Exécution & \route{POST /storage/actions/\{id\}/confirm} et \fn{confirm\_and\_execute\_action}\\
Refus & \route{POST /storage/actions/\{id\}/reject} + note d'incident\\
Vérification & \fn{\_close\_the\_loop} + \file{services/incident\_verification.py}\\
\bottomrule
\end{longtable}

\subsection{Le raisonnement derrière la séparation \texttt{DESTRUCTIVE} / \texttt{REMEDIATION}}

Le commentaire de \file{server.py} explique pourquoi deux ensembles et non un :

\begin{quote}
\itshape
« La distinction avec \texttt{DESTRUCTIVE\_TOOLS} n'est pas la même : un restore
est destructif, relancer une sauvegarde ne l'est pas --- mais les deux changent
l'état du système, et aucun des deux ne doit partir sur la seule décision du
modèle. \texttt{DESTRUCTIVE\_TOOLS} continue de qualifier le risque dans
l'audit. »
\end{quote}

\noindent
Autrement dit : \texttt{REMEDIATION\_TOOLS} répond à « faut-il un humain ? »
(oui pour les quatre), \texttt{DESTRUCTIVE\_TOOLS} répond à « quel niveau de
risque ? » (HIGH pour un seul).

\subsection{Ce qui rend le contournement difficile}

\begin{enumerate}
  \item \textbf{Le modèle ne peut pas exécuter directement.} L'aiguillage est un
  test d'appartenance à un ensemble figé, en dur dans le code.
  \item \textbf{Écrire l'appel en texte ne sert à rien.} Un JSON recopié dans la
  réponse ne crée rien --- et \fn{\_recover\_written\_call} ne le transforme
  qu'en \emph{proposition}, avec double validation de l'outil et de la cible.
  \item \textbf{Le prompt le dit au modèle.} « never say a remediation has been
  launched: it is waiting on a human decision. »
  \item \textbf{Le résultat renvoyé au modèle le répète.} « Action NOT executed
  \ldots{} Present it as a recommendation, never as a fact. »
  \item \textbf{Sans utilisateur identifié, refus net.} Si \texttt{created\_by}
  est \texttt{None}, l'agent renvoie
  \texttt{\{"status": "denied", "reason": "confirmation\_required"\}} --- « on
  refuse plutôt que de l'attribuer à quelqu'un au hasard ».
  \item \textbf{Double exécution impossible.} Le test
  \texttt{status != PENDING} lève une \cls{ValueError}, traduite en 404.
  \item \textbf{Deux outils exigent en plus \texttt{confirm=True}} au niveau de
  l'implémentation.
\end{enumerate}

\begin{warnbox}
\textbf{La limite honnête à énoncer.} Les routes de \file{actions.py} n'exigent
que \fn{get\_current\_user}, pas \fn{require\_role}. Un utilisateur authentifié
avec le rôle \texttt{linux\_engineer} ne peut ni ouvrir le chat ni voir les
incidents (ces routes sont protégées), mais s'il obtenait un identifiant
d'action valide, la route \route{/confirm} ne le rejetterait pas. Ajouter
\fn{require\_role} sur ce routeur serait le durcissement le plus rentable du
projet.
\end{warnbox}

\subsection{Le workflow complet, en une table}

\begin{longtable}{@{}p{0.7cm}p{3.7cm}p{4.3cm}p{5.7cm}@{}}
\toprule
\textbf{\#} & \textbf{Étape} & \textbf{Acteur} & \textbf{Trace laissée}\\
\midrule
\endhead
1 & Demande utilisateur & Ingénieur & \tbl{chat\_messages} (rôle \texttt{user})\\
2 & Lecture de la VM & Agent + MCP + SSH & \tbl{mcp\_calls} (\texttt{success})\\
3 & Analyse & Groq & \tbl{chat\_messages} (rôle \texttt{assistant})\\
4 & Proposition & \fn{\_propose\_action} & \tbl{action\_requests} (\texttt{pending})\\
5 & Rattachement & \fn{attach\_actions} & \texttt{chat\_message\_id} renseigné\\
6 & Affichage & \cls{ActionProposal} & ---\\
7 & \textbf{Décision humaine} & \textbf{Ingénieur} & \texttt{confirmed\_by}, \texttt{confirmed\_at}\\
8 & Exécution & MCP + SSH & \tbl{mcp\_calls} avec \texttt{action\_request\_id}\\
9 & Résultat & \fn{confirm\_and\_execute\_action} & \texttt{status}, \texttt{result} (JSONB), \texttt{completed\_at}\\
10 & Vérification & \fn{verify} & \tbl{mcp\_calls} (relecture) + \tbl{incident\_ticket\_notes}\\
11 & Clôture & \fn{resolve\_ticket} & \texttt{status='resolved'}, \texttt{resolved\_at}, note\\
\bottomrule
\end{longtable}

\begin{keybox}
\textbf{Onze étapes, sept traces en base.} C'est ce qui permet de répondre, six
mois plus tard, à la question « qui a lancé cette restauration, sur la foi de
quel relevé, et qu'est-ce que ça a donné ? ».
\end{keybox}

% =====================================================================
\chapter{Les technologies utilisées}
\label{ch:technos}

Chaque technologie est justifiée par ce qu'elle apporte \emph{à ce projet},
avec le fichier qui la met en œuvre. Aucune définition générale.

% =====================================================================
\section{Backend}

\subsection{Python}
\textbf{Où} : tout le backend. Le README annonce Python 3.13.\\
\textbf{Pourquoi ici} : c'est la seule langue où cohabitent naturellement le SDK
MCP officiel, un client SSH (\texttt{paramiko}), un ORM mature (SQLAlchemy) et
les bibliothèques d'embeddings (\texttt{sentence-transformers}). Le projet
utilise aussi des constructions récentes du langage : \texttt{str | None},
\texttt{asyncio.to\_thread}, \texttt{@dataclass}, \texttt{match} implicite via
les dictionnaires de dispatch.

\subsection{FastAPI (+ Starlette, Uvicorn)}
\textbf{Où} : \file{app/main.py}, les 11 routeurs.\\
\textbf{Pourquoi ici} : trois raisons visibles dans le code.
\begin{itemize}
  \item \textbf{L'injection de dépendances} rend le contrôle d'accès
  déclaratif : \texttt{current\_user: User = Depends(require\_role(...))} suffit
  à protéger une route.
  \item \textbf{Le mélange \texttt{def} / \texttt{async def}} est exploité
  volontairement. Les routes qui font du SSH bloquant restent en \texttt{def}
  pour partir dans le threadpool (\file{dashboard.py}, \file{reports.py}) ; les
  routes qui attendent l'agent sont en \texttt{async def} (\file{storage.py},
  \route{/investigate}, \route{/confirm}). Les commentaires l'expliquent
  fichier par fichier.
  \item \textbf{Le \fn{lifespan}} héberge les trois tâches de fond avec
  annulation propre.
\end{itemize}

\subsection{Pydantic v2 et pydantic-settings}
\textbf{Où} : \file{app/schemas/}, \file{app/core/config.py}.\\
\textbf{Pourquoi ici} : la validation d'entrée est la première barrière.
\cls{ServiceRequestCreate} impose \texttt{title} entre 1 et 200 caractères ;
\cls{LoginRequest} valide le format d'e-mail (via \texttt{email-validator}) ;
\cls{ChatRequest} accepte un \texttt{conversation\_id} optionnel typé.
Côté configuration, \cls{Settings} \textbf{refuse de démarrer} si
\texttt{database\_url}, \texttt{jwt\_secret\_key} ou \texttt{groq\_api\_key}
manquent --- une panne au démarrage plutôt qu'à la première requête.

\begin{infobox}
Les schémas séparent volontairement la liste du détail : \cls{IncidentTicketOut}
pour la liste, \cls{IncidentTicketDetail} pour la fiche. Le docstring justifie :
« y joindre notes et actions ferait autant de requêtes supplémentaires que de
lignes affichées, pour des données que la liste ne montre pas. »
\end{infobox}

\subsection{SQLAlchemy (mode synchrone)}
\textbf{Où} : \file{app/models/}, \file{app/database/session.py}.\\
\textbf{Pourquoi \emph{synchrone}} : c'est un choix, et il est documenté. Les
docstrings de \file{incident\_service.py}, \file{mcp\_audit.py} et
\file{actions.py} portent tous la trace d'un correctif : une première version
avait été écrite pour \texttt{AsyncSession}, avec des \texttt{await
db.commit()} qui échouaient (« sur une session sync, \fn{db.commit()} renvoie
\texttt{None} et \texttt{await None} lève un \texttt{TypeError} »). Le projet a
convergé vers le synchrone, et FastAPI compense en exécutant les routes
\texttt{def} dans un threadpool.

\subsection{Alembic}
\textbf{Où} : \file{backend/alembic/}, 14 révisions.\\
\textbf{Pourquoi ici} : le schéma a beaucoup évolué (ajout du chat, des
conversations multiples, des demandes, des pièces jointes, du lien
investigation). Une migration comme \file{b6d0f3a91c47} fait ce qu'un
\texttt{create\_all} ne saurait pas faire : ajouter un index unique
\emph{partiel} PostgreSQL (\texttt{postgresql\_where}).

\subsection{PostgreSQL 16}
\textbf{Où} : \file{docker-compose.yml}, tous les modèles.\\
\textbf{Pourquoi ici} : JSONB interrogeable (\texttt{mcp\_data[key].astext}),
UUID natif, types ENUM, et surtout \texttt{date\_bin} --- fonction PostgreSQL
14+ sans laquelle le graphique de performance devrait être agrégé en Python.

\subsection{python-jose et passlib[bcrypt]}
\textbf{Où} : \file{app/core/security.py}.\\
\textbf{Pourquoi ici} : JWT HS256 pour une API sans état, et bcrypt pour ne
jamais stocker un mot de passe en clair. Le jeton porte le rôle, mais le contrôle
d'accès relit le rôle en base --- un changement de rôle prend effet
immédiatement.

\subsection{MCP (SDK officiel, version 2.0.0)}
\textbf{Où} : \file{app/mcp/storage/}.\\
\textbf{Pourquoi ici} : c'est le protocole qui permet d'exposer des capacités à
un modèle \emph{sans coder à la main} un catalogue de fonctions, un schéma JSON
par paramètre et un routeur d'appels. Le décorateur \texttt{mcp.tool()} dérive
tout des annotations de type et des docstrings. Le transport \texttt{stdio}
isole en outre les outils dans un processus séparé.

\subsection{Paramiko}
\textbf{Où} : \file{app/mcp/storage/ssh.py}.\\
\textbf{Pourquoi ici} : parce que le projet a fait le choix de \textbf{ne rien
simuler}. Il n'existe pas d'API REST sur une VM Ubuntu pour lire \texttt{df} ou
\texttt{lvs} : SSH est le protocole d'administration réel. Authentification par
clé, \texttt{AutoAddPolicy}, timeouts de 5\,s (connexion) et 15\,s (commande).

\subsection{Groq (SDK \texttt{groq})}
\textbf{Où} : \file{app/llm/groq\_client.py}.\\
\textbf{Pourquoi ici} : deux raisons opérationnelles. La \textbf{latence} ---
un ingénieur qui investigue un incident n'attend pas ; et le \textbf{tool
calling natif}, compatible avec le format OpenAI, ce qui permet de convertir le
catalogue MCP en \texttt{tools} d'une seule fonction. Le client essaie
successivement plusieurs modèles et retient le premier qui répond.

\subsection{ChromaDB}
\textbf{Où} : \file{app/rag/vectorstore.py}, conteneur \texttt{chromadb/chroma}.\\
\textbf{Pourquoi ici} : la base de connaissances est petite (quatre documents
d'origine plus les dépôts du Data Hub). Chroma s'installe en un conteneur,
accepte des embeddings fournis de l'extérieur, et supporte \texttt{upsert} ---
ce qui rend l'indexation relançable et la suppression d'un document possible.

\subsection{sentence-transformers}
\textbf{Où} : \file{app/rag/embedding\_service.py}.\\
\textbf{Pourquoi ce modèle précis} : le fichier l'explique. « la base de
connaissances Storage est en anglais (NetApp, Veeam, TrueNAS, AWS) alors que les
ingénieurs posent leurs questions en français \ldots{} Un modèle anglophone comme
\texttt{bge-small-en-v1.5} place ces deux langues dans des régions éloignées de
l'espace vectoriel : la question ne retrouve pas le passage qui y répond. »
D'où \texttt{intfloat/multilingual-e5-small}, avec ses préfixes asymétriques.

\subsection{ReportLab}
\textbf{Où} : \file{app/services/report\_service.py}.\\
\textbf{Pourquoi ReportLab et pas WeasyPrint} --- le docstring tranche :
« WeasyPrint s'appuie sur les bibliothèques GTK (libgobject, pango, cairo) qui
ne sont pas installées avec pip sous Windows, et son import échoue sur le poste
de développement. »

\begin{warnbox}
\textbf{Détail à corriger un jour.} \texttt{reportlab} figure dans
\file{requirements.lock.txt} mais \textbf{pas} dans \file{requirements.txt}.
Une installation faite depuis le seul \file{requirements.txt} ne pourrait pas
générer de rapport. C'est une observation, pas une supposition : la ligne est
absente.
\end{warnbox}

\subsection{pypdf}
\textbf{Où} : extraction du texte des PDF déposés au Data Hub.\\
\textbf{Pourquoi ici} : sans couche texte, pas d'indexation --- et le service le
dit à l'utilisateur en lui suggérant un OCR.

% =====================================================================
\section{Frontend}

\begin{longtable}{@{}p{3.2cm}p{2.0cm}p{9.5cm}@{}}
\toprule
\textbf{Technologie} & \textbf{Version} & \textbf{Pourquoi dans ce projet}\\
\midrule
\endhead
React & 18.3 & Interface fortement dynamique : un fil de chat qui grandit, des propositions d'action qui changent d'état après un clic, un drawer d'incident. L'état local suffit, aucun Redux n'est nécessaire\\
TypeScript & 5.7 & \file{lib/types.ts} (259 lignes) reproduit les schémas Pydantic. Renommer un champ côté backend casse la compilation côté frontend --- c'est le seul contrat automatique entre les deux moitiés du projet\\
Vite & 5.4 & Rechargement instantané, et surtout le \textbf{proxy} \route{/api} \arrowto{} \texttt{127.0.0.1:8000} qui supprime tout problème CORS en développement\\
TailwindCSS & 3.4 & Cohérence visuelle sans feuille de style séparée. Une palette \texttt{brand} / \texttt{ink} / \texttt{accent} / \texttt{danger} est réutilisée jusque dans les rapports PDF\\
react-router-dom & 6.30 & Routes imbriquées : \cls{RequireAuth} est une route parente qui rend \cls{AppLayout}, les pages sont ses enfants\\
recharts & 2.15 & Un seul graphique (\cls{PerformanceChart}), mais avec trois séries et trois plages temporelles\\
lucide-react & 0.475 & Icônes cohérentes dans la sidebar, les badges et les boutons\\
clsx + tailwind-merge & --- & \fn{cn()} : composer des classes conditionnelles sans conflit Tailwind\\
\bottomrule
\end{longtable}

% =====================================================================
\section{Infrastructure et outillage}

\subsection{Docker et Docker Compose}
\textbf{Où} : \file{docker-compose.yml}.\\
\textbf{Ce qui est réellement conteneurisé} : \texttt{postgres:16} (avec
\texttt{healthcheck} \texttt{pg\_isready} et volume nommé) et
\texttt{chromadb/chroma:latest} (port 8001 exposé, volume nommé).\\
\textbf{Pourquoi ici} : ces deux services sont pénibles à installer nativement,
surtout sous Windows. Le backend et le frontend, eux, tournent en local.

\subsection{SSH}
\textbf{Pourquoi ici} : c'est le protocole d'administration de la VM. Le projet
n'installe aucun agent sur la machine cible : il exécute les mêmes commandes
qu'un ingénieur taperait.

\subsection{Technologies Linux mobilisées sur la VM}

\begin{longtable}{@{}p{3.4cm}p{11.4cm}@{}}
\toprule
\textbf{Technologie} & \textbf{Rôle dans le projet}\\
\midrule
\endhead
LVM & Volumes logiques, snapshots (\texttt{lvcreate -{}-snapshot}), restauration (\texttt{lvconvert -{}-merge}). Lu par \texttt{lvs}\\
systemd & Les jobs de sauvegarde \emph{sont} des timers systemd. \texttt{list-timers}, \texttt{show}, \texttt{start}, les drop-ins et \texttt{revert}\\
journald & \fn{get\_backup\_logs} lit \texttt{journalctl -o short-iso}\\
\texttt{/proc/diskstats} & Source unique des métriques de performance\\
\texttt{df} & Source unique de la capacité\\
restic & Le logiciel de sauvegarde réel ; \fn{repair\_backup\_repository} sait lire l'état d'un dépôt\\
MinIO & Cible objet S3 du laboratoire, avec Object Lock\\
NFS / Samba & Partages du laboratoire NAS\\
\bottomrule
\end{longtable}

\begin{keybox}
\textbf{L'argument de fond, à retenir pour la soutenance.} Le choix
\texttt{systemd} + \texttt{restic} + \texttt{LVM} n'est pas un pis-aller : ce
sont les briques que l'outil interroge \emph{réellement}. Chaque relevé affiché
correspond à une commande qu'un jury pourrait rejouer lui-même en SSH sur la
VM.
\end{keybox}

% =====================================================================
\section{Ce qui est déclaré mais pas utilisé}

Un point d'honnêteté qui vaut mieux qu'une question embarrassante.

\begin{longtable}{@{}p{3.6cm}p{11.2cm}@{}}
\toprule
\textbf{Élément} & \textbf{Statut réel}\\
\midrule
\endhead
Anthropic Claude & Clé déclarée dans \cls{Settings}, \file{llm\_router.py} lève \texttt{NotImplementedError}. \textbf{Non branché}\\
\file{requirements.txt} & Contient un très grand nombre de paquets sans rapport (Jupyter, Streamlit, MLflow, XGBoost, OpenCV, Flask, Plotly\ldots), vestiges d'un environnement de travail global. La liste réellement nécessaire est bien plus courte\\
WeasyPrint & Présent dans les dépendances, remplacé par ReportLab pour la raison ci-dessus\\
\file{base\_persona.py}, \file{llm\_router.py}, \file{rbac\_middleware.py} & Code valide, importé par personne\\
\texttt{DESTRUCTIVE\_TOOLS}, \fn{risk\_for}, \fn{invoke} & Définis, jamais appelés\\
\tbl{request\_messages} & Table conservée, plus alimentée\\
Persona Linux & Le README l'annonce, et un commit récent (\texttt{df219f23}) porte « retrait de la persona Linux ». Seul le persona Storage existe dans le code\\
\bottomrule
\end{longtable}

% =====================================================================
\chapter{Analyse fichier par fichier}
\label{ch:fichiers}

Chaque fiche suit la même grille : objectif, fonctions, classes, dépendances,
entrées, sorties, fichiers appelés, fichiers appelants, explication du code, et
ce qu'il faut en retenir pour la présentation.

% =====================================================================
\section{\file{backend/app/main.py}}

\textbf{Objectif} --- Assembler l'application : middlewares, routeurs, tâches de
fond.

\textbf{Fonctions principales} --- \fn{lifespan} (gestionnaire de contexte
asynchrone), \fn{\_prime\_storage\_health}, \fn{root}.

\textbf{Classes principales} --- Aucune ; l'objet \texttt{app = FastAPI(...)} est
créé au niveau du module.

\textbf{Dépendances} --- FastAPI, \cls{CORSMiddleware}, \cls{AuditMiddleware},
les 11 routeurs, \file{dashboard/kpi\_service.py},
\file{services/incident\_watcher.py}, \file{services/metrics\_collector.py},
\file{app/models} (import global qui enregistre tous les modèles).

\textbf{Entrées} --- Requêtes HTTP. \textbf{Sorties} --- Réponses HTTP ; la
racine \route{GET /} renvoie \texttt{\{"status": "OASIS AI Copilot backend
running"\}}.

\textbf{Fichiers appelés} --- Tous les routeurs, les deux collecteurs.
\textbf{Fichiers qui l'appellent} --- Uvicorn
(\texttt{uvicorn app.main:app -{}-reload -{}-port 8000}, cf.
\file{.claude/launch.json}).

\textbf{Explication} --- Le \fn{lifespan} crée trois tâches
(\fn{metrics\_collector.run\_forever}, \fn{incident\_watcher.run\_forever},
\fn{\_prime\_storage\_health}), \texttt{yield}, puis les annule et attend leur
\cls{CancelledError} à l'arrêt. \fn{\_prime\_storage\_health} sort le premier
relevé SSH du chemin d'une requête : « sans lui, le premier affichage du
dashboard après un redémarrage attendrait la poignée de main SSH ».

\textbf{À retenir} --- Un backend qui \emph{travaille sans qu'on le lui
demande}. C'est ce qui distingue OASIS d'un simple chatbot : il surveille,
détecte et mesure en continu.

% =====================================================================
\section{\file{backend/app/core/config.py}}

\textbf{Objectif} --- Centraliser la configuration, validée au démarrage.

\textbf{Classes} --- \cls{Settings(BaseSettings)}, instanciée une fois en
\fn{settings}.

\textbf{Dépendances} --- \texttt{pydantic-settings}.
\textbf{Entrées} --- Variables d'environnement et \file{.env}.
\textbf{Sorties} --- L'objet \fn{settings}.

\textbf{Fichiers qui l'appellent} --- \file{security.py}, \file{session.py},
\file{ssh.py}, \file{groq\_client.py}, \file{vectorstore.py},
\file{embedding\_service.py}, \file{knowledge\_service.py},
\file{request\_service.py}.

\textbf{À retenir} --- Les variables SSH sont \emph{optionnelles} par choix, avec
une erreur explicite au moment de l'appel. Le nom du modèle d'embedding est
configurable, et le changer invalide les collections Chroma --- ce que le code
détecte au lieu de laisser passer.

% =====================================================================
\section{\file{backend/app/api/deps.py}}

\textbf{Objectif} --- Fournir les dépendances d'authentification et
d'autorisation.

\textbf{Fonctions} --- \fn{get\_current\_user(credentials, db)},
\fn{require\_role(*allowed\_roles)}.

\textbf{Entrées} --- En-tête \texttt{Authorization: Bearer}, session DB.
\textbf{Sorties} --- Une entité \cls{User}, ou \cls{HTTPException} 401/403.

\textbf{Fichiers appelés} --- \file{database/session.py},
\file{core/security.py}, \file{models/user.py}.
\textbf{Fichiers qui l'appellent} --- Les 11 routeurs.

\textbf{À retenir} --- Deux contrôles, pas un : le jeton \emph{et} l'état du
compte en base. Un compte désactivé est rejeté immédiatement, sans attendre
l'expiration de son jeton.

% =====================================================================
\section{\file{backend/app/api/v1/storage.py}}

\textbf{Objectif} --- Exposer le chat : gestion des fils et tour de dialogue.

\textbf{Fonctions} --- \fn{\_require\_conversation}, \fn{list\_conversations},
\fn{conversation\_messages}, \fn{rename\_conversation},
\fn{delete\_conversation}, \fn{chat} (\texttt{async}).

\textbf{Dépendances} --- \file{services/chat\_service.py},
\file{personas/storage/agent.py}, \file{services/incident\_investigation.py},
\file{services/incident\_service.py}, \file{schemas/chat.py}.

\textbf{Entrées} --- \cls{ChatRequest}, \cls{ConversationUpdate}, identifiants de
fil. \textbf{Sorties} --- \cls{ConversationOut}, \cls{ChatMessageOut},
\cls{ChatTurnOut}, 204.

\textbf{Explication} --- L'ordre des opérations dans \fn{chat} est le point clé :
l'agent tourne \emph{avant} toute écriture. Un échec de Groq ou de la VM
n'ajoute donc ni fil vide ni question sans réponse.

\textbf{À retenir} --- Le 404 sur un fil appartenant à un collègue est un choix de
sécurité : « dire "interdit" plutôt qu'"introuvable" confirmerait son
existence ».

% =====================================================================
\section{\file{backend/app/personas/storage/agent.py}}

\textbf{Objectif} --- Le cerveau : RAG, dialogue avec le LLM, appels d'outils,
garde-fou de sécurité, détection d'anomalie.

\textbf{Classes} --- \cls{IncidentContext}, \cls{AgentTurn},
\cls{StoragePersona} (instanciée en \fn{storage\_persona}).

\textbf{Fonctions} --- \fn{\_candidate\_calls}, \fn{\_recover\_written\_call},
\fn{StoragePersona.\_propose\_action}, \fn{StoragePersona.handle\_message}.

\textbf{Dépendances} --- \file{llm/groq\_client.py}, \file{mcp/storage/client.py},
\file{mcp/storage/server.py} (\texttt{REMEDIATION\_TOOLS}),
\file{personas/storage/incident\_rules.py},
\file{personas/storage/prompts.py}, \file{rag/retriever.py},
\file{services/incident\_service.py}, \file{database/session.py}.

\textbf{Entrées} --- \texttt{user\_message}, \texttt{created\_by},
\texttt{history}, \texttt{incident}.
\textbf{Sorties} --- \cls{AgentTurn}. Effets de bord : lignes
\tbl{mcp\_calls}, \tbl{action\_requests}, \tbl{incident\_tickets}.

\textbf{Fichiers qui l'appellent} --- \file{api/v1/storage.py},
\file{services/incident\_investigation.py}.

\textbf{Constantes} --- \texttt{\_MAX\_TOOL\_ROUNDS = 6},
\texttt{\_TARGET\_ARGUMENTS = ("target", "volume\_id", "job\_id")},
\texttt{\_JSON\_OBJECT} (expression régulière).

\textbf{À retenir} --- Trois idées à exposer : \emph{(1)} le garde-fou qui
transforme un appel d'outil en proposition ; \emph{(2)} la correction
automatique de la cible (libellé \arrowto{} identifiant technique) ; \emph{(3)}
le rattrapage d'un appel écrit en prose, avec double validation.

% =====================================================================
\section{\file{backend/app/personas/storage/prompts.py}}

\textbf{Objectif} --- Le prompt système du persona.

\textbf{Contenu} --- Une constante, \fn{STORAGE\_SYSTEM\_PROMPT}. Elle définit
quatre domaines d'expertise, quatre règles de sûreté (ne jamais prétendre avoir
agi sans résultat d'outil, demander confirmation pour une action irréversible,
séparer faits / hypothèses / recommandations, préférer la documentation interne
et ne jamais citer une documentation absente du contexte), et impose une
réponse en anglais.

\textbf{À retenir} --- Une nuance de rédaction remarquable : la spécialité
« façonne le ton et les exemples ; elle ne limite pas les questions acceptées ».
Le prompt interdit explicitement de refuser une question hors domaine --- un
défaut classique des assistants spécialisés.

% =====================================================================
\section{\file{backend/app/personas/storage/incident\_rules.py}}

\textbf{Objectif} --- Décider, de façon déterministe, ce qui constitue un
incident.

\textbf{Fonctions} --- \fn{\_extract\_items}, \fn{check\_capacity},
\fn{check\_backup\_status}, \fn{evaluate\_tool\_result}.

\textbf{Constantes} --- \texttt{CAPACITY\_WARNING\_THRESHOLD = 85},
\texttt{CAPACITY\_CRITICAL\_THRESHOLD = 95},
\texttt{\_KNOWN\_LIST\_WRAPPER\_KEYS}, \texttt{\_RULES\_BY\_TOOL}.

\textbf{Entrées} --- Nom d'outil et résultat brut.
\textbf{Sorties} --- Liste de dictionnaires prêts pour
\fn{create\_ticket\_from\_diagnosis}.

\textbf{Fichiers qui l'appellent} --- \file{agent.py},
\file{incident\_watcher.py}, \file{incident\_verification.py}.

\textbf{Explication} --- \fn{\_extract\_items} normalise trois formes
d'enveloppe (liste brute, dictionnaire enveloppant une liste connue,
dictionnaire plat). Le docstring documente un bug corrigé :
« \fn{get\_capacity()} retourne \texttt{\{"volumes": [...]\}} \ldots{}
L'ancienne version ne gérait que les deux derniers cas, donc ne détectait jamais
rien sur \fn{get\_capacity}. »

\textbf{À retenir} --- Le même fichier sert à \emph{ouvrir} et à \emph{fermer}
un ticket. C'est ce qui garantit la cohérence : ce qui a déclenché l'alerte est
exactement ce qui doit avoir disparu.

% =====================================================================
\section{\file{backend/app/mcp/storage/server.py}}

\textbf{Objectif} --- Registre des outils et classification du risque.

\textbf{Contenu} --- \fn{TOOL\_REGISTRY} (10 entrées),
\texttt{DESTRUCTIVE\_TOOLS} (1 entrée), \texttt{REMEDIATION\_TOOLS} (4 entrées),
\fn{risk\_for}, \fn{invoke}.

\textbf{Fichiers qui l'appellent} --- \file{tool\_server.py}
(\fn{TOOL\_REGISTRY}), \file{agent.py} et \file{incident\_investigation.py}
(\texttt{REMEDIATION\_TOOLS}).

\textbf{À retenir} --- 60 lignes qui portent toute la politique de sécurité
d'exécution. Signaler honnêtement que \fn{risk\_for} et \fn{invoke} ne sont pas
encore utilisés.

% =====================================================================
\section{\file{backend/app/mcp/storage/tool\_server.py}}

\textbf{Objectif} --- Le serveur MCP lui-même.

\textbf{Contenu} --- 12 lignes : import de \cls{MCPServer}, instanciation,
boucle d'enregistrement, \texttt{mcp.run(transport="stdio")}.

\textbf{Entrées} --- JSON-RPC sur \texttt{stdin}.
\textbf{Sorties} --- JSON-RPC sur \texttt{stdout}.

\textbf{Fichiers appelés} --- \file{server.py} et, transitivement, les sept
modules d'outils. \textbf{Lancé par} --- \file{client.py}, comme sous-processus.

\textbf{À retenir} --- La brièveté est l'argument : le SDK MCP dérive les
schémas des annotations de type et des docstrings, il n'y a donc rien à écrire
à la main.

% =====================================================================
\section{\file{backend/app/mcp/storage/client.py}}

\textbf{Objectif} --- Piloter le serveur MCP et auditer chaque appel.

\textbf{Classes} --- \cls{StorageMCPClient} (protocole de gestionnaire de
contexte asynchrone).

\textbf{Fonctions} --- \fn{start}, \fn{stop}, \fn{\_require\_session},
\fn{list\_tools\_for\_llm}, \fn{call\_tool}.

\textbf{Dépendances} --- \texttt{mcp.ClientSession},
\texttt{mcp.client.stdio}, \file{models/mcp\_call.py},
\file{services/mcp\_audit.py}.

\textbf{Entrées de \fn{call\_tool}} --- \texttt{tool\_name}, puis les mots-clés
\texttt{db}, \texttt{user\_id}, \texttt{correlation\_id},
\texttt{action\_request\_id}, et \texttt{**arguments}.
\textbf{Sorties} --- \texttt{dict} décodé, ou \cls{RuntimeError}.

\textbf{Fichiers qui l'appellent} --- \file{agent.py}, \file{actions.py} ;
\file{incident\_verification.py} reçoit le client déjà démarré.

\textbf{À retenir} --- Pas d'instance partagée au niveau du module : c'est la
correction d'un bug réel, documenté dans \file{actions.py}. Le verrou
\texttt{asyncio.Lock} sérialise les appels, le transport stdio n'étant pas
multiplexé.

% =====================================================================
\section{\file{backend/app/mcp/storage/ssh.py}}

\textbf{Objectif} --- Point d'accès unique à la VM.

\textbf{Classes} --- \cls{StorageVMError(RuntimeError)}.
\textbf{Fonctions} --- \fn{\_client}, \fn{run}, \fn{run\_many}, \fn{quote}.
\textbf{Constantes} --- \texttt{CONNECT\_TIMEOUT\_SECONDS = 5},
\texttt{COMMAND\_TIMEOUT\_SECONDS = 15}.

\textbf{Entrées} --- Une commande, ou un dictionnaire de commandes.
\textbf{Sorties} --- \texttt{stdout} en \texttt{str}, ou \cls{StorageVMError}.

\textbf{Fichiers qui l'appellent} --- Les sept modules d'outils.

\textbf{Explication} --- \fn{\_client} vérifie d'abord que les trois variables
SSH sont renseignées et nomme celles qui manquent. Les messages d'erreur sont
volontairement \emph{en anglais} : « Ces messages sont affichés tels quels dans
l'interface, qui est en anglais. » La connexion est fermée dans un
\texttt{finally}.

\textbf{À retenir} --- 130 lignes qui isolent tout le risque d'exécution. Aucun
outil ne construit sa propre connexion, aucun n'échappe à \fn{quote}.

% =====================================================================
\section{\file{backend/app/mcp/storage/tools/backup.py}}

\textbf{Objectif} --- Lire et piloter les jobs de sauvegarde réels.

\textbf{Fonctions} --- \fn{\_looks\_like\_backup}, \fn{\_parse\_timer\_units},
\fn{\_service\_for}, \fn{\_is\_storage\_backup}, \fn{\_show\_units},
\fn{\_parse\_show}, \fn{\_moment}, \fn{\_hours\_since},
\fn{\_job\_from\_properties}, \fn{list\_backups}, \fn{get\_backup\_job\_status},
\fn{run\_backup}, \fn{get\_backup\_logs}.

\textbf{Sorties} --- \texttt{\{"jobs": [...], "configured": bool\}} ou
\texttt{\{"error": ..., "jobs": []\}}.

\textbf{Trois finesses à connaître}
\begin{enumerate}
  \item \textbf{Filtrage des units système.} \fn{\_is\_storage\_backup} écarte
  \texttt{dpkg-db-backup.timer}, « qui sauvegarde la base de paquets d'Ubuntu »
  : une unit livrée par un paquet n'est retenue que si son nom cite un vrai
  outil de sauvegarde (restic, borg, veeam\ldots).
  \item \textbf{Lecture groupée.} \fn{\_show\_units} concatène toutes les
  commandes en une seule, séparées par le marqueur
  \texttt{\#\#\#OASIS-UNIT=}. Motif : « sur six timers, \fn{list\_backups}
  mettait près de quatre secondes, et le dashboard se figeait d'autant à chaque
  navigation. »
  \item \textbf{« Jamais lancé » n'est pas « a réussi ».} systemd renvoie
  \texttt{Result=success} sur une unit qui n'a jamais démarré. Le code teste
  \texttt{ExecMainStartTimestamp} : sans cela, « un job planifié mais jamais
  exécuté s'afficherait comme une sauvegarde réussie ».
\end{enumerate}

\textbf{À retenir} --- \texttt{configured: False} est distinct de
\texttt{error} : rien n'est configuré, ce n'est pas une panne. Cette distinction
remonte jusqu'au dashboard et jusqu'aux insights.

% =====================================================================
\section{\file{backend/app/mcp/storage/tools/repair.py}}

\textbf{Objectif} --- Réparer un job dont le dépôt restic est introuvable.

\textbf{Fonctions} --- \fn{\_service\_for}, \fn{\_parse\_environment},
\fn{\_parse\_environment\_files}, \fn{\_repository\_state},
\fn{\_healthy\_donor}, \fn{\_drop\_in\_body}, \fn{\_drop\_in\_path},
\fn{\_write\_commands}, \fn{\_write\_drop\_in}, \fn{\_missing\_privilege},
\fn{repair\_backup\_repository}.

\textbf{Entrées} --- \texttt{target}, \texttt{confirm=False}.
\textbf{Sorties} --- Un dictionnaire dont le champ \texttt{action} prend l'une de
sept valeurs : \texttt{repair\_failed}, \texttt{no\_repair\_possible},
\texttt{no\_repair\_needed}, \texttt{repair\_requires\_confirmation},
\texttt{repository\_repaired}.

\textbf{À retenir} --- Le seul outil qui \emph{modifie une configuration}, et
celui qui refuse le plus souvent d'agir. Trois idées à présenter : le
diagnostic du dépôt par un faux mot de passe, le refus en cas d'égalité entre
voisins, et la réversibilité par \texttt{systemctl revert}.

% =====================================================================
\section{\file{backend/app/api/v1/actions.py}}

\textbf{Objectif} --- Le seul endroit du système qui exécute une action sur la
VM.

\textbf{Fonctions} --- \fn{list\_my\_actions}, \fn{propose\_action},
\fn{confirm\_action} (\texttt{async}), \fn{\_close\_the\_loop}
(\texttt{async}), \fn{reject\_action}, \fn{get\_action}.

\textbf{Dépendances} --- \file{mcp/storage/client.py},
\file{services/incident\_service.py}, \file{services/incident\_verification.py},
\file{models/action\_request.py}, \file{schemas/incident.py}.

\textbf{Fichiers qui l'appellent} --- \file{frontend/src/components/ActionProposal.tsx}.

\textbf{À retenir} --- Le docstring de \fn{confirm\_action} raconte un bug de
sécurité corrigé, exemple parfait de traçabilité de conception. Signaler aussi
l'absence de \fn{require\_role} sur ce routeur.

% =====================================================================
\section{\file{backend/app/services/incident\_service.py}}

\textbf{Objectif} --- Toutes les écritures liées aux incidents et aux actions.

\textbf{Fonctions} --- \fn{create\_action\_request},
\fn{confirm\_and\_execute\_action}, \fn{reject\_action\_request},
\fn{\_generate\_ticket\_number}, \fn{\_open\_ticket\_for},
\fn{create\_ticket\_from\_diagnosis}, \fn{list\_tickets}, \fn{get\_ticket},
\fn{update\_ticket\_status}, \fn{mark\_in\_progress}, \fn{add\_note},
\fn{resolve\_ticket}, \fn{pending\_actions}, \fn{has\_completed\_action}.

\textbf{Fichiers qui l'appellent} --- \file{agent.py}, \file{actions.py},
\file{incidents.py}, \file{incident\_investigation.py},
\file{incident\_watcher.py}, \file{report\_service.py}.

\textbf{À retenir} --- Trois protections en dix lignes dans
\fn{confirm\_and\_execute\_action} : anti-double-exécution, commit de
l'intention avant l'action, et détection d'une erreur retournée plutôt que
levée.

% =====================================================================
\section{\file{backend/app/services/incident\_watcher.py}}

\textbf{Objectif} --- Détecter les échecs de sauvegarde sans intervention
humaine, et refermer les tickets vérifiés.

\textbf{Fonctions} --- \fn{\_system\_user\_id}, \fn{\_parse\_timestamp},
\fn{\_latest\_ticket}, \fn{\_already\_handled}, \fn{\_verify\_resolved},
\fn{scan\_once}, \fn{run\_forever}.

\textbf{Constantes} --- \texttt{INTERVAL\_SECONDS = 300},
\texttt{SYSTEM\_ACCOUNT\_EMAIL = "oasis-system@oasis.local"}.

\textbf{Sorties} --- Un rapport
\texttt{\{"created", "skipped", "closed", "checked"\}} ; jamais d'exception.

\textbf{À retenir} --- Le compte de service désactivé est une élégance :
l'audit distingue « ce qu'OASIS a détecté seul » de « ce qu'un humain a
demandé », sans créer de faille de connexion.

% =====================================================================
\section{\file{backend/app/dashboard/kpi\_service.py}}

\textbf{Objectif} --- Produire, sans LLM, tout ce qu'affiche le dashboard.

\textbf{Fonctions} --- \fn{\_volume\_status}, \fn{\_build\_capacity},
\fn{\_build\_backups}, \fn{\_build\_incidents}, \fn{\_health\_score},
\fn{\_headline}, \fn{refresh\_health\_snapshot}, \fn{\_storage\_health},
\fn{\_refresh\_in\_background}, \fn{\_invalidate\_health\_snapshot},
\fn{get\_overview}, \fn{get\_performance}, \fn{get\_recent\_activity},
\fn{get\_insights}.

\textbf{État partagé} --- \texttt{\_health\_cache} protégé par deux verrous
(\texttt{\_health\_lock} pour la lecture/écriture, \texttt{\_refresh\_lock} pour
sérialiser les rafraîchissements).

\textbf{À retenir} --- Le cache de 60 secondes avec rafraîchissement en arrière-
plan est la réponse à un problème mesuré : deux connexions SSH à chaque
affichage de page.

% =====================================================================
\section{\file{backend/app/rag/embedding\_service.py} et \file{vectorstore.py}}

\textbf{Objectif} --- Vectoriser, indexer et rechercher.

\textbf{Fonctions} --- \fn{get\_embedding\_model} (chargement paresseux, une
fois par processus), \fn{embed\_query}, \fn{embed\_passages} ;
\fn{get\_chroma\_client}, \fn{get\_or\_create\_collection}, \fn{add\_chunks},
\fn{search}.

\textbf{Classes} --- \cls{EmbeddingModelMismatch(RuntimeError)}.

\textbf{À retenir} --- Deux garde-fous rarement vus dans un projet d'école : la
séparation stricte query/passage imposée par l'asymétrie des modèles E5, et le
refus d'une collection indexée avec un autre modèle --- parce que deux modèles
de même dimension produiraient « des réponses fausses, mais parfaitement
plausibles ».

% =====================================================================
\section{\file{frontend/src/lib/api.ts}}

\textbf{Objectif} --- Client HTTP unique du frontend.

\textbf{Objets exportés} --- \cls{ApiError}, \fn{tokenStore},
\fn{setUnauthorizedHandler}, \fn{api} (\fn{get}, \fn{post}, \fn{patch},
\fn{delete}, \fn{upload}, \fn{download}).

\textbf{Fichiers qui l'appellent} --- \cls{AuthContext} et les huit pages.

\textbf{À retenir} --- \fn{api.download} explique bien un problème réel : « un
simple \texttt{<a href>} ne convient pas : le navigateur n'y joint pas
l'en-tête \texttt{Authorization} ». D'où le passage par un blob et une URL
objet.

% =====================================================================
\section{\file{frontend/src/components/ActionProposal.tsx}}

\textbf{Objectif} --- Matérialiser la confirmation humaine.

\textbf{Fonctions} --- \fn{describe}, \fn{resultText}, \fn{settle},
composant \cls{ActionProposal}.

\textbf{Entrées} --- \texttt{action: ActionRequest},
\texttt{onSettled?: (updated) => void}.
\textbf{Sorties} --- Appels \route{POST /storage/actions/\{id\}/confirm} ou
\route{/reject}, notifications \emph{toast}.

\textbf{Fichiers qui l'appellent} --- \cls{ChatPage} et \cls{IncidentsPage}.

\textbf{À retenir} --- Le composant est \emph{partagé} entre le chat et la fiche
d'incident, « pour que la décision se prenne au même endroit que la lecture,
quel que soit l'écran ». Il affiche à la fois la phrase lisible et l'appel
technique brut \texttt{action\_type(target)}, ce qui permet de vérifier la cible
avant de cliquer.

% =====================================================================
\section{\file{frontend/src/pages/IncidentsPage.tsx}}

\textbf{Objectif} --- Liste, fiche détaillée et pilotage du cycle de vie.

\textbf{Constantes} --- \texttt{FILTERS} (all, critical, high, medium, low,
resolved), \texttt{LIFECYCLE} (5 étapes).

\textbf{Fonctions} --- \fn{queryFor}, \fn{reachedStage}, \fn{Lifecycle},
\fn{IncidentBody}, \fn{DetailRow}.

\textbf{À retenir} --- \texttt{LIFECYCLE} n'est pas décoratif. Le commentaire du
fichier le dit : « L'ordre n'est pas décoratif : c'est le contrat de
l'application. Rien ne s'exécute sur la VM sans être passé par "Action proposée"
puis par le clic humain de "Confirmation", et un ticket ne se referme qu'une
fois l'effet vérifié. »

\begin{lstlisting}[caption={Les cinq temps affichés, et la fonction qui calcule l'étape atteinte}]
const LIFECYCLE = [
  { key: 'detected',  label: 'Detected' },
  { key: 'diagnosed', label: 'Diagnosed' },
  { key: 'proposed',  label: 'Action proposed' },
  { key: 'confirmed', label: 'Confirmed' },
  { key: 'resolved',  label: 'Resolved' },
];

function reachedStage(incident) {
  if (incident.status === 'resolved' || incident.status === 'closed') return 4;
  if (incident.actions.some((a) => a.status === 'completed'))          return 3;
  if (incident.actions.length > 0)                                     return 2;
  if (incident.root_cause)                                             return 1;
  return 0;
}
\end{lstlisting}

% =====================================================================
\chapter{Annexes}

\section{Fichiers présents mais vides}

\begin{longtable}{@{}p{7.4cm}p{7.4cm}@{}}
\toprule
\textbf{Fichier} & \textbf{Ce que son nom laisse attendre}\\
\midrule
\endhead
\file{backend/Dockerfile} & Image du backend\\
\file{frontend/Dockerfile} & Image du frontend\\
\file{nginx/nginx.conf} & Reverse proxy\\
\file{.github/workflows/ci.yml} & Intégration continue\\
\file{backend/app/core/permissions.py} & Matrice de permissions\\
\file{backend/app/core/constants.py} & Constantes partagées\\
\file{backend/app/database/init\_db.py} & Initialisation du schéma\\
\file{backend/app/orchestrator/orchestrator.py} & Orchestrateur multi-agents\\
\file{backend/app/orchestrator/cross\_domain\_manager.py} & Gestion cross-domain\\
\file{backend/app/llm/models.py}, \file{prompt\_manager.py} & Abstractions LLM\\
\file{backend/app/middleware/auth\_middleware.py} & Middleware d'authentification\\
\file{backend/app/reports/generator.py} & Générateur de rapports\\
\file{backend/app/models/report.py}, \file{llm\_usage.py} & Modèles non créés\\
\file{backend/app/schemas/} (plusieurs) & \file{audit\_log.py}, \file{llm\_usage.py}, \file{mcp\_call.py}, \file{metric.py}, \file{report.py}, \file{role.py}\\
\bottomrule
\end{longtable}

\begin{infobox}
Ces fichiers ne gênent rien : ils ne sont importés par aucun module actif. Ils
témoignent d'une architecture cible plus vaste que ce qui a été réalisé --- ce
qui est normal pour un projet de stage, et vaut mieux dit qu'esquivé.
\end{infobox}

\section{Code présent mais non appelé}

\begin{longtable}{@{}p{5.4cm}p{9.4cm}@{}}
\toprule
\textbf{Élément} & \textbf{Observation}\\
\midrule
\endhead
\file{personas/base\_persona.py} & \cls{BasePersona} n'est héritée par personne ; \cls{StoragePersona} ne l'étend pas\\
\file{orchestrator/llm\_router.py} & \fn{route()} lève \texttt{NotImplementedError} pour \texttt{COMPLEX} ; appelée uniquement par \cls{BasePersona}\\
\file{middleware/rbac\_middleware.py} & \fn{require\_role} en double de celle de \file{api/deps.py} ; c'est la seconde qui est importée partout\\
\fn{risk\_for}, \fn{invoke} & Définies dans \file{mcp/storage/server.py}, jamais appelées\\
\fn{get\_calls\_for\_correlation} & Définie dans \file{mcp\_audit.py}, jamais appelée\\
\cls{IncidentAnalysisRequest} / \cls{IncidentAnalysisResponse} & Schémas définis, aucune route ne les utilise\\
\tbl{request\_messages} & Table conservée pour l'historique, plus alimentée\\
\bottomrule
\end{longtable}

\section{Comptes de test}

Créés par \texttt{python -m app.database.seed} :

\begin{longtable}{@{}p{5.0cm}p{4.0cm}p{5.2cm}@{}}
\toprule
\textbf{E-mail} & \textbf{Mot de passe} & \textbf{Rôle}\\
\midrule
\endhead
\texttt{linux@oasis.com} & \texttt{linux123} & \texttt{linux\_engineer}\\
\texttt{storage@oasis.com} & \texttt{storage123} & \texttt{storage\_engineer}\\
\texttt{admin@oasis.com} & \texttt{admin123} & \texttt{administrator}\\
\bottomrule
\end{longtable}

\begin{warnbox}
Ces identifiants sont en clair dans \file{app/database/seed.py}. Acceptable pour
une maquette, à retirer avant tout déploiement.
\end{warnbox}

\section{Démarrage de l'environnement}

\begin{lstlisting}[caption={Séquence complète, d'après le code et .claude/launch.json}]
# 1. Dependances (racine du depot)
docker compose up -d                 # postgres:16 + chromadb

# 2. Backend
cd backend
alembic upgrade head                 # 14 migrations
python -m app.database.seed          # roles + comptes de test
python -m app.rag.reindex            # (re)construction des collections Chroma
uvicorn app.main:app --reload --port 8000

# 3. Frontend
cd frontend
npm install
npm run dev                          # http://localhost:5173, proxy /api -> :8000
\end{lstlisting}

\section{Dix points à retenir pour la soutenance}

\begin{enumerate}
  \item \textbf{Rien n'est simulé.} Chaque chiffre affiché vient d'une commande
  Linux exécutée par SSH sur une vraie VM.
  \item \textbf{Le modèle propose, l'humain décide.} Quatre outils sur dix ne
  s'exécutent jamais depuis le chat.
  \item \textbf{Les seuils sont dans le code, pas dans le modèle.}
  \file{incident\_rules.py} et \file{health.py} décident ce qui est un incident.
  \item \textbf{Un ticket ne se ferme que sur une vérification.}
  \file{incident\_verification.py} relit l'indicateur qui l'avait ouvert.
  \item \textbf{Tout est audité.} \tbl{mcp\_calls}, \tbl{audit\_logs},
  \tbl{incident\_ticket\_notes}, \tbl{request\_events}.
  \item \textbf{Le RAG est multilingue par nécessité.} KB en anglais, questions
  en français ; d'où E5 multilingue et ses préfixes asymétriques.
  \item \textbf{MCP isole les outils dans un processus séparé} et fournit le
  catalogue que le modèle lit pour choisir.
  \item \textbf{Trois filets empêchent de perdre une remédiation} : appel
  d'outil, rattrapage d'un appel écrit en prose, action de repli calculée par la
  règle.
  \item \textbf{Le système travaille seul.} Deux boucles de fond détectent,
  mesurent et referment sans qu'on le leur demande.
  \item \textbf{Les limites sont connues et nommées} : CORS ouvert, absence de
  \fn{require\_role} sur les actions, jeton cross-domain non vérifié,
  conteneurisation incomplète, Claude non branché.
\end{enumerate}

\begin{keybox}
Le dernier point est le plus fort. Un projet dont l'auteur sait exactement où
sont les limites inspire davantage confiance qu'un projet présenté comme
parfait.
\end{keybox}


\end{document}
